\documentclass[11pt,a4paper]{article}

\usepackage[utf8]{inputenc}
\usepackage[T1]{fontenc}
\usepackage[spanish,es-noshorthands,es-tabla]{babel}
\usepackage{lmodern}
\usepackage[margin=2.5cm]{geometry}
\usepackage{amsmath,amssymb,amsthm,mathtools}
\usepackage{booktabs,array,tabularx}
\usepackage{enumitem}
\usepackage{xcolor}
\usepackage{listings}
\usepackage[hidelinks]{hyperref}

\decimalpoint

% ---------- Entornos ----------
\theoremstyle{plain}
\newtheorem{teorema}{Teorema}[section]
\newtheorem{proposicion}[teorema]{Proposición}
\newtheorem{lema}[teorema]{Lema}
\newtheorem{corolario}[teorema]{Corolario}
\theoremstyle{definition}
\newtheorem{definicion}[teorema]{Definición}
\newtheorem{ejemplo}[teorema]{Ejemplo}
\newtheorem{contraejemplo}[teorema]{Contraejemplo}
\theoremstyle{remark}
\newtheorem{observacion}[teorema]{Observación}

% ---------- Notación ----------
\newcommand{\R}{\mathbb{R}}
\newcommand{\N}{\mathbb{N}}
\newcommand{\Z}{\mathbb{Z}}
\newcommand{\E}{\mathbb{E}}
\newcommand{\Prob}{\mathbb{P}}
\newcommand{\Var}{\operatorname{Var}}
\newcommand{\Cov}{\operatorname{Cov}}
\newcommand{\Corr}{\operatorname{Corr}}
\newcommand{\med}{\operatorname{med}}
\newcommand{\IQR}{\operatorname{IQR}}
\newcommand{\MAD}{\operatorname{MAD}}
\newcommand{\SC}{\operatorname{SC}}
\newcommand{\logit}{\operatorname{logit}}
\newcommand{\tr}{\operatorname{tr}}
\newcommand{\WoE}{\operatorname{WoE}}
\newcommand{\IV}{\operatorname{IV}}
\newcommand{\PMI}{\operatorname{PMI}}
\newcommand{\idf}{\operatorname{idf}}
\newcommand{\tf}{\operatorname{tf}}
\newcommand{\df}{\operatorname{df}}
\newcommand{\col}{\operatorname{col}}
\newcommand{\rango}{\operatorname{rango}}
\newcommand{\1}{\mathbf{1}}
\newcommand{\ind}[1]{\mathbf{1}\{#1\}}
\newcommand{\Dtr}{\mathcal D_{\mathrm{tr}}}
\newcommand{\Dte}{\mathcal D_{\mathrm{te}}}

% ---------- Código ----------
\definecolor{fondo}{gray}{0.96}
\lstset{
  basicstyle=\ttfamily\small,
  backgroundcolor=\color{fondo},
  frame=single, rulecolor=\color{gray!50},
  breaklines=true, columns=fullflexible, keepspaces=true,
  showstringspaces=false,
  literate={á}{{\'a}}1 {é}{{\'e}}1 {í}{{\'i}}1 {ó}{{\'o}}1 {ú}{{\'u}}1 {ñ}{{\~n}}1
           {Á}{{\'A}}1 {É}{{\'E}}1 {Í}{{\'I}}1 {Ó}{{\'O}}1 {Ú}{{\'U}}1
}
\newenvironment{aws}{\par\medskip\sloppy\noindent\textbf{Implementación en AWS.}\ }{\par\medskip}
\setlength{\emergencystretch}{3em}
\newcolumntype{Y}{>{\raggedright\arraybackslash}X}

\title{Ingeniería de características:\\ desarrollo matemático e implementación en AWS}
\author{Notas complementarias al capítulo 3 de la guía MLA-C01}
\date{Documentación de AWS verificada el 24 de septiembre de 2026}

\begin{document}
\maketitle
\tableofcontents
\newpage

% =====================================================================
\section{Marco general}
% =====================================================================

Estas notas desarrollan matemáticamente las técnicas de preparación de
datos e ingeniería de características del capítulo 3 de la guía, añaden
dos codificaciones supervisadas que el capítulo no cubre (\emph{target
encoding} y \emph{weight of evidence}) y describen, para cada técnica, el
servicio de AWS con que se implementa. El nivel de formalidad es el de los
textos clásicos de estadística y aprendizaje estadístico.

\subsection{Modelo probabilístico y notación}\label{sec:notacion}

\paragraph{Datos.} Se modela cada observación como un par aleatorio
$(X,Y)$ con distribución conjunta $P$ desconocida, donde
$X=(X_1,\dots,X_p)$ es el vector de \emph{características} (covariables)
e $Y$ la \emph{respuesta} (variable objetivo). Cada $X_j$ toma valores en
$\R$ (característica numérica) o en un conjunto finito (característica
categórica, más abajo). Una \emph{muestra} es una sucesión
$\mathcal D=((x_1,y_1),\dots,(x_n,y_n))$ de realizaciones de $n$ copias
independientes de $(X,Y)$; se escribe $x_k=(x_{k1},\dots,x_{kp})$.

\paragraph{Estadísticos de una columna numérica.} Una \emph{columna} es un
vector $z=(z_1,\dots,z_n)\in\R^n$, típicamente $z_k=x_{kj}$ para un $j$
fijo. Se usan:
\begin{itemize}[leftmargin=*]
\item la media $\bar z=n^{-1}\sum_{k=1}^nz_k$;
\item las desviaciones estándar
$s_n(z)=\bigl(n^{-1}\sum_k(z_k-\bar z)^2\bigr)^{1/2}$ y
$s_{n-1}(z)=\bigl((n-1)^{-1}\sum_k(z_k-\bar z)^2\bigr)^{1/2}$; cuando el
divisor es irrelevante o se indica en el contexto se escribe $s$;
\item los estadísticos de orden $z_{(1)}\le z_{(2)}\le\dots\le z_{(n)}$,
es decir, los valores de la columna ordenados de menor a mayor;
\item el \emph{cuantil muestral} de orden $q\in[0,1]$, definido por
interpolación lineal entre estadísticos de orden: con $h=(n-1)q+1$ y
$r=\lfloor h\rfloor$,
\[
Q(q)=z_{(r)}+(h-r)\bigl(z_{(r+1)}-z_{(r)}\bigr)\qquad(\text{con }z_{(n+1)}:=z_{(n)}).
\]
Es la definición por defecto de NumPy y pandas, las bibliotecas que usa el
capítulo. La mediana es $\med(z)=Q(1/2)$, los cuartiles son
$Q_1=Q(1/4)$ y $Q_3=Q(3/4)$, y el rango intercuartílico es
$\IQR=Q_3-Q_1$.
\end{itemize}
Por ejemplo, para la columna de consumo de energía del capítulo,
\[
z=(11,\,12,\,12.5,\,8,\,8.5,\,7.5,\,5,\,4.7,\,6,\,23),
\]
se tiene $n=10$; para $q=3/4$,
$h=7.75$, $r=7$, y como $z_{(7)}=11$ y $z_{(8)}=12$, resulta
$Q_3=11+0.75\cdot1=11.75$; análogamente $Q_1=6.375$.

\paragraph{Características categóricas.} Una característica categórica
toma valores en un conjunto finito $\mathcal C=\{c_1,\dots,c_K\}$ cuyos
elementos se llaman \emph{niveles}; $K$ es la \emph{cardinalidad}. En una
muestra, $n_i=\#\{k:x_{kj}=c_i\}$ es el número de observaciones del nivel
$c_i$.

\subsection{Transformaciones ajustadas y fuga de información}

Casi todas las técnicas de este documento reemplazan una característica
por una función de ella cuyos parámetros se calculan con datos: la media y
la desviación estándar de un escalado, el exponente $\lambda$ de Box--Cox,
la tabla de medias de un \emph{target encoding}, el vocabulario de un
vectorizador de texto, las direcciones de un PCA. El capítulo separa
«ajustar» (\texttt{fit}) de «aplicar» (\texttt{transform}); esta sección
precisa esa separación y demuestra por qué es necesaria.

\begin{definicion}[Transformación ajustada]\label{def:ajustada}
Sean $\mathcal X$ el conjunto de valores posibles de $X$, $\mathcal X'$
otro conjunto y $\Theta$ un conjunto de parámetros. Una \emph{transformación
ajustada} consta de
\begin{enumerate}[label=(\roman*)]
\item una familia de funciones $T_\theta:\mathcal X\to\mathcal X'$,
$\theta\in\Theta$, y
\item un \emph{estimador} $\hat\theta$, es decir, una regla que a cada
muestra asigna un elemento de $\Theta$.
\end{enumerate}
Dada una muestra de entrenamiento $\Dtr$, \emph{ajustar} es calcular
$\hat\theta(\Dtr)$; \emph{aplicar} es evaluar $T_{\hat\theta(\Dtr)}$ en
cualquier $x\in\mathcal X$, pertenezca o no a $\Dtr$.
\end{definicion}

\begin{ejemplo}
Estandarización de la característica $j$: $\Theta=\R\times(0,\infty)$,
$T_{(a,b)}(x)$ reemplaza $x_j$ por $(x_j-a)/b$ y deja las demás
coordenadas, y $\hat\theta(\Dtr)=(\bar z,s_n(z))$ con $z$ la columna $j$
de $\Dtr$. Un ejemplo sin parámetros es $T(x)=\log x_j$: $\Theta$ tiene un
solo elemento y ajustar no hace nada.
\end{ejemplo}

Para evaluar un modelo se fija una \emph{función de pérdida}
$L(y,\hat y)\ge0$ (por ejemplo, $(y-\hat y)^2$ en regresión o
$\ind{y\ne\hat y}$ en clasificación) y un \emph{modelo ajustado}
$\hat f:\mathcal X'\to\mathcal Y$, que también se calcula con $\Dtr$. La
cantidad de interés es el \emph{riesgo condicional}
\[
R=\E\Bigl[L\bigl(Y,\hat f(T_{\hat\theta}(X))\bigr)\,\Big|\,\Dtr\Bigr],
\]
donde $(X,Y)$ es una observación nueva independiente de $\Dtr$: el error
esperado del procedimiento ya entrenado sobre datos futuros. Se estima con
una muestra de prueba.

\begin{proposicion}[Insesgadez del error de prueba]\label{prop:test}
Sean $\Dtr$ y $\Dte=((X'_1,Y'_1),\dots,(X'_m,Y'_m))$ independientes, con
los pares $(X'_l,Y'_l)$ independientes entre sí y con la distribución $P$
de $(X,Y)$. Supóngase que $\hat\theta$ y $\hat f$ se calculan solo con
$\Dtr$ y que $R<\infty$. Entonces
$\hat R=m^{-1}\sum_{l=1}^mL\bigl(Y'_l,\hat f(T_{\hat\theta}(X'_l))\bigr)$
cumple $\E[\hat R\mid\Dtr]=R$.
\end{proposicion}

\begin{proof}
Sea $G(\mathcal D,x,y)=L\bigl(y,\hat f_{\mathcal D}(T_{\hat\theta(\mathcal
D)}(x))\bigr)$, donde el subíndice indica la muestra con que se calculan
$\hat f$ y $\hat\theta$. Se usa la siguiente propiedad de la esperanza
condicional: si $U$ y $V$ son independientes y $G\ge0$, entonces
$\E[G(U,V)\mid U]=g(U)$ con $g(u)=\E[G(u,V)]$. Con $U=\Dtr$ y
$V=(X'_l,Y'_l)$, que son independientes por hipótesis y $V$ tiene
distribución $P$, se obtiene
$\E[G(\Dtr,X'_l,Y'_l)\mid\Dtr]=g(\Dtr)$, donde
$g(\mathcal D)=\E[G(\mathcal D,X,Y)]$ con $(X,Y)\sim P$. Aplicando la misma
propiedad con $V=(X,Y)$ se tiene también $R=g(\Dtr)$. Promediando sobre
$l$ y usando la linealidad de la esperanza condicional,
$\E[\hat R\mid\Dtr]=g(\Dtr)=R$.
\end{proof}

La hipótesis de que $\hat\theta$ no use la muestra de prueba es necesaria.

\begin{contraejemplo}\label{ce:fuga}
Sea $X$ real con distribución continua e $Y\in\{0,1\}$ independiente de
$X$ con $\Prob(Y=1)=1/2$, de modo que todo clasificador tiene riesgo $1/2$
con pérdida $\ind{y\neq\hat y}$. Por la continuidad, con probabilidad $1$
las $X'_l$ son distintas entre sí. Defínase un parámetro
estimado con $\Dtr\cup\Dte$: $\hat\theta$ es la tabla
$\{(X'_l,Y'_l)\}_{l=1}^m$, y $T_{\hat\theta}(x)$ añade a $x$ la
coordenada $Y'_l$ si $x=X'_l$ y $0$ en otro caso. El clasificador que
devuelve esa coordenada añadida tiene $\hat R=0$. Una observación
verdaderamente nueva coincide con alguna $X'_l$ con probabilidad $0$
(nuevamente por la continuidad), así que el clasificador predice $0$ y su
riesgo es $\Prob(Y=1)=1/2$.
\end{contraejemplo}

El contraejemplo es extremo, pero su mecanismo es el de cualquier
parámetro que use las respuestas de prueba; las
secciones~\ref{sec:target} y~\ref{sec:woe} muestran que con
\emph{target encoding} aparece incluso dentro del entrenamiento. Cuando el
parámetro no usa respuestas (una media, un mínimo), la proposición deja
igualmente de aplicarse y la magnitud del sesgo depende del caso. La
proposición justifica también la regla operativa del capítulo: calcular
$\hat\theta$ una vez con entrenamiento y aplicar la misma
$T_{\hat\theta}$ en validación, prueba y producción. Se llama
\emph{training-serving skew} a la situación en que la transformación
aplicada en producción difiere de la $T_{\hat\theta}$ con que se entrenó
$\hat f$; entonces $\hat f$ recibe entradas con otra distribución que la de
entrenamiento, y $R$ deja de describir su desempeño.

\begin{aws}
Una transformación ajustada vive en uno de tres lugares.
(i) \textbf{Amazon SageMaker Data Wrangler} (integrado hoy en SageMaker
Canvas; la versión de Studio Classic sigue documentada): cada paso de un
\emph{data flow} que aprende de los datos (codificar categóricas,
vectorizar texto, tratar atípicos o faltantes) guarda sus parámetros,
calculados por defecto sobre la muestra importada (las primeras 50\,000
filas); al exportar puede reajustarlos (\emph{refit}) sobre el dataset
completo, y los parámetros se trasladan al \emph{inference pipeline}.
(ii) \textbf{AWS Glue DataBrew}: una \emph{receta} es una lista JSON de
acciones que un trabajo serverless aplica al dataset en Amazon S3; los
estadísticos que usan sus acciones se calculan sobre el dataset al que se
aplica la receta en cada ejecución. (iii) \textbf{SageMaker Processing}: un
contenedor (por ejemplo, el de scikit-learn) ejecuta un script propio; es
la vía para técnicas sin paso integrado y la que da control explícito de
\texttt{fit} sobre entrenamiento y \texttt{transform} sobre lo demás. Los
parámetros que deben consultarse en inferencia (por ejemplo, tablas de
codificación) pueden guardarse en S3 o en un \emph{feature group} de
\textbf{SageMaker Feature Store}, cuyo \emph{online store} los sirve con
baja latencia.
\end{aws}

\subsection{Tres propiedades de invariancia}\label{sec:invariancias}

La conveniencia de una transformación depende del modelo que la recibe.
Esta sección demuestra qué transformaciones no alteran a los árboles de
decisión y a la regresión lineal, y muestra que otros modelos sí son
sensibles; los resultados se usan en todo el documento.

\begin{definicion}[Árbol de regresión construido por CART]\label{def:cart}
Sea $\mathcal D$ una muestra con respuesta numérica. Para
$A\subset\{1,\dots,n\}$ no vacío, sean $\bar y_A$ la media de las $y_k$ con
$k\in A$ y $\SC(A)=\sum_{k\in A}(y_k-\bar y_A)^2$. Un \emph{nodo} es un
conjunto de índices $N$. Un \emph{corte} del nodo $N$ sobre la columna $j$
con umbral $t\in\R$ es la partición de $N$ en
$N_{\le}=\{k\in N:x_{kj}\le t\}$ y $N_{>}=N\setminus N_{\le}$, con ambos
no vacíos. El algoritmo parte de $N=\{1,\dots,n\}$ y, mientras una regla
de parada que depende solo de $N$ y de $(y_k)_{k\in N}$ no se cumpla,
elige entre todos los cortes de todas las columnas uno que maximice
$\Delta=\SC(N)-\SC(N_\le)-\SC(N_>)$, desempata con una regla que depende
solo de la columna y de la partición $\{N_\le,N_>\}$ (sin distinguir
cuál de los dos conjuntos es el de la izquierda), y repite en cada hijo. Las
\emph{hojas} son los nodos donde se detiene. Para clasificación, $\SC$ se
reemplaza por $|A|$ veces el índice de Gini de las respuestas en $A$, que es $1-\sum_c\hat p_c^2$ con $\hat p_c$ la proporción de la clase $c$ en $A$.
\end{definicion}

Las condiciones sobre la regla de parada y el desempate no son técnicas:
si el desempate dependiera del valor numérico del umbral $t$, el resultado
podría cambiar al transformar la columna aunque los cortes candidatos
fueran los mismos.

\begin{lema}\label{lem:cortes}
Sea $g:\R\to\R$ estrictamente creciente. Para todo nodo $N$ y toda columna
$j$, las particiones de $N$ obtenibles con cortes sobre
$(x_{kj})_{k\in N}$ coinciden con las obtenibles con cortes sobre
$(g(x_{kj}))_{k\in N}$.
\end{lema}

\begin{proof}
Sean $v_1<\dots<v_r$ los valores distintos de $(x_{kj})_{k\in N}$. Para
un umbral $t$, sea $i(t)$ el número de $v_i$ que son $\le t$; entonces
$N_\le=\{k\in N:x_{kj}\le v_{i(t)}\}$ (vacío si $i(t)=0$), de modo que las
particiones obtenibles son las
$\{k:x_{kj}\le v_i\}$, $i=1,\dots,r-1$. Como $g$ es estrictamente
creciente, $x_{kj}\le v_i$ equivale a $g(x_{kj})\le g(v_i)$, y los valores
distintos de la columna transformada son $g(v_1)<\dots<g(v_r)$. Aplicando
el mismo razonamiento a la columna transformada, sus particiones
obtenibles son las $\{k:g(x_{kj})\le g(v_i)\}=\{k:x_{kj}\le v_i\}$,
$i=1,\dots,r-1$: las mismas.
\end{proof}

\begin{proposicion}[Árboles e invariancia monótona]\label{prop:arbol}
Si se reemplaza la columna $j$ por $(g(x_{kj}))_k$ con $g$ estrictamente
creciente, el árbol de la Definición~\ref{def:cart} produce las mismas
hojas (como conjuntos de índices).
\end{proposicion}

\begin{proof}
Por inducción sobre los nodos en el orden en que el algoritmo los crea.
La raíz es la misma. Si un nodo $N$ es el mismo en ambas ejecuciones, la
regla de parada toma la misma decisión porque solo depende de $N$ y de las
respuestas. Por el Lema~\ref{lem:cortes} (aplicado a la columna $j$; las
demás columnas no cambian), el conjunto de pares (columna, partición)
candidatos es el mismo; el criterio $\Delta$ depende solo de la partición
y de las respuestas, y el desempate solo de la columna y la partición. Por
tanto se elige la misma partición y los hijos coinciden.
\end{proof}

Los umbrales numéricos cambian (en la columna transformada el umbral
aplicable es $g(v_i)$), pero ninguna observación cambia de hoja. El mismo
argumento vale para los métodos que agregan árboles de este tipo (bosques
aleatorios, \emph{gradient boosting}), siempre que sus cortes candidatos se
definan como arriba; las implementaciones que discretizan primero las
columnas en histogramas de ancho fijo no cumplen la hipótesis.

\begin{definicion}[Mínimos cuadrados con intercepto]
Sea $X\in\R^{n\times p}$ y $Z=[\1\;X]$, con $\1\in\R^n$ el vector de unos.
Un estimador de mínimos cuadrados es cualquier $\hat\beta$ que minimice
$\|y-Z\beta\|^2$; el vector de valores ajustados es $\hat y=Z\hat\beta$.
\end{definicion}

\begin{lema}\label{lem:proy}
$\hat y$ no depende del minimizador elegido: es la proyección ortogonal de
$y$ sobre $\col(Z)$, el espacio generado por las columnas de $Z$.
\end{lema}

\begin{proof}
Sea $\pi$ la proyección ortogonal de $y$ sobre $\col(Z)$, caracterizada
por $\pi\in\col(Z)$ e $y-\pi\perp\col(Z)$. Para todo $\beta$, por el
teorema de Pitágoras, $\|y-Z\beta\|^2=\|y-\pi\|^2+\|\pi-Z\beta\|^2$,
porque $\pi-Z\beta\in\col(Z)$. El mínimo es $\|y-\pi\|^2$ y se alcanza
exactamente cuando $Z\beta=\pi$.
\end{proof}

\begin{proposicion}[Mínimos cuadrados e invariancia afín]\label{prop:ols}
Sea $\tilde X=XD+\1c^\top$ con $D$ diagonal con entradas no nulas y
$c\in\R^p$, es decir, cada columna $j$ se reemplaza por
$d_jx_{\cdot j}+c_j$ con $d_j\ne0$. Los valores ajustados por mínimos
cuadrados con intercepto son los mismos con $X$ y con $\tilde X$.
\end{proposicion}

\begin{proof}
$[\1\;\tilde X]=[\1\;X]M$ con
$M=\begin{psmallmatrix}1&c^\top\\0&D\end{psmallmatrix}$, que es invertible
porque $\det M=\det D\ne0$. Si $u=[\1\;\tilde X]\alpha$, entonces
$u=[\1\;X](M\alpha)$, y si $u=[\1\;X]\beta$, entonces
$u=[\1\;\tilde X](M^{-1}\beta)$; los espacios columna coinciden. Por el
Lema~\ref{lem:proy}, los valores ajustados coinciden.
\end{proof}

La hipótesis de intercepto es necesaria para la parte de traslación: sin
él, $\col(X)$ y $\col(X+\1c^\top)$ pueden diferir (con $p=1$,
$x=(1,2)^\top$ genera la recta de dirección $(1,2)$ y $x+1=(2,3)^\top$ otra
distinta).

\begin{contraejemplo}[Regularización y vecinos más cercanos]\label{ce:ridge}
(a) Regresión \emph{ridge} con una característica y sin intercepto: el
estimador que minimiza $\sum_k(y_k-\beta x_k)^2+\lambda\beta^2$ ($\lambda>0$)
es $\hat\beta=\sum_kx_ky_k/(\sum_kx_k^2+\lambda)$, que se obtiene
igualando a cero la derivada de esa función cuadrática estrictamente
convexa; el valor ajustado en un punto $u\in\R$ es $\hat\beta u$. Si la
característica se reemplaza por $ax$, $a\ne0$ (en entrenamiento y en el
punto, que pasa a ser $au$), el valor ajustado pasa a ser
$au\cdot\frac{a\sum_kx_ky_k}{a^2\sum_kx_k^2+\lambda}
=u\frac{\sum_kx_ky_k}{\sum_kx_k^2+\lambda/a^2}$, que depende de $a$:
cambiar la unidad de la característica equivale a cambiar la penalización.
(b) El clasificador del vecino más cercano asigna a un punto $q$ la
respuesta del punto de entrenamiento más cercano en distancia euclídea.
Con $q=(0,0)$, $B=(1,0)$ y $C=(0,2)$, el más cercano es $B$ (distancias
$1$ y $2$). Si la primera coordenada se multiplica por $10$, $B=(10,0)$ y
el más cercano pasa a ser $C$.
\end{contraejemplo}

Los árboles no necesitan escalado (Proposición~\ref{prop:arbol}), la
regresión lineal sin penalizar tampoco para predecir
(Proposición~\ref{prop:ols}), y los métodos penalizados y basados en
distancias sí (Contraejemplo~\ref{ce:ridge}); la
Proposición~\ref{prop:cond} añade el descenso de gradiente a esta última
lista.

% =====================================================================
\section{Preparación previa: valores faltantes y atípicos}
% =====================================================================

\subsection{Valores faltantes}\label{sec:faltantes}

Se considera una característica numérica $X$ que puede faltar y un vector
$W$ de otras características siempre observadas. Sea $R=1$ si $X$ falta y
$R=0$ si se observa; en la muestra, $r_k$ es el valor de $R$ en la
observación $k$ y $O=\{k:r_k=0\}$ el conjunto de observaciones con $X$
presente, de tamaño $n_o$.

La práctica anterior a Rubin (1976) era el \emph{análisis de casos
completos}: descartar toda fila con algún faltante. Su costo se ve con un
cálculo: si cada una de $p$ columnas falta de forma independiente con
probabilidad $q$, una fila está completa con probabilidad $(1-q)^p$, que
para $q=0.05$ y $p=50$ vale $0.077$; se pierde más del 90\,\% de la
muestra. La alternativa, imputar, plantea la pregunta de cuándo lo que se
calcula con los valores observados describe correctamente a la población.
La respuesta depende de la relación entre $R$ y los datos.

\begin{definicion}[Mecanismos de falta]\label{def:mcar}
La falta de $X$ es
\begin{enumerate}[label=(\roman*)]
\item \emph{completamente al azar} (MCAR) si $R$ es independiente de
$(X,W)$;
\item \emph{al azar} (MAR) si $R$ y $X$ son condicionalmente independientes
dado $W$;
\item \emph{no al azar} (MNAR) si no se cumple (ii).
\end{enumerate}
\end{definicion}

MCAR implica MAR: si $R$ es independiente de $(X,W)$, entonces
$\Prob(R=1\mid X,W)=\Prob(R=1)$, que en particular no depende de $X$ una
vez fijado $W$. Ejemplos: una falla aleatoria del sensor es MCAR; que los
jóvenes declaren menos su ingreso, con la edad en $W$, es MAR; que quien
gana más declare menos su ingreso es MNAR. La jerarquía importa por lo
siguiente.

\begin{proposicion}\label{prop:mcar}
Supóngase $\E|X|<\infty$ y $\Prob(R=0)>0$.
\begin{enumerate}[label=(\alph*)]
\item Bajo MCAR, $\E[X\mid R=0]=\E[X]$.
\item Bajo MAR, $\E[X\mid W,R=0]=\E[X\mid W]$ (para los valores de $W$ con
$\Prob(R=0\mid W)>0$).
\end{enumerate}
\end{proposicion}

\begin{proof}
(a) Por la independencia de $R$ y $X$,
$\E[X\ind{R=0}]=\E[X]\Prob(R=0)$; dividiendo por $\Prob(R=0)$ se obtiene
$\E[X\mid R=0]=\E[X]$. (b) Es el mismo argumento condicionado a $W$: por
la independencia condicional,
$\E[X\ind{R=0}\mid W]=\E[X\mid W]\Prob(R=0\mid W)$, y se divide por
$\Prob(R=0\mid W)$.
\end{proof}

La parte (a) dice que bajo MCAR la media de los observados estima sin
sesgo la media poblacional; la (b), que bajo MAR la regresión de $X$ sobre
$W$ ajustada solo con las filas observadas estima la regresión poblacional,
lo que justifica imputar con una predicción a partir de $W$ en lugar de con
una constante. Los ejemplos siguientes muestran que las hipótesis no pueden
quitarse.

\begin{ejemplo}[MAR sin MCAR]
Sea $W\sim N(0,1)$, $X=W+\varepsilon$ con $\varepsilon\sim N(0,1)$
independiente de $W$, y $R=\ind{W>0}$. La falta es MAR (es función de $W$)
pero no MCAR. $\E[X\mid R=0]=\E[W\mid W\le0]=-\varphi(0)/\Phi(0)=-\sqrt{2/\pi}\approx-0.80$,
donde $\varphi$ y $\Phi$ son la densidad y la función de distribución
normales estándar (se usa $\int_{-\infty}^0w\varphi(w)\,dw=-\varphi(0)$,
consecuencia de $\varphi'(w)=-w\varphi(w)$). La media de los observados
subestima $\E[X]=0$, pero $\E[X\mid W,R=0]=W=\E[X\mid W]$, como predice
(b).
\end{ejemplo}

\begin{ejemplo}[MNAR]\label{ej:mnar}
Sea $X\sim N(0,1)$ observada solo si $X<1$ (sin $W$). Entonces
$\E[X\mid R=0]=\E[X\mid X<1]=-\varphi(1)/\Phi(1)\approx-0.288$ por el mismo
cálculo. Los valores observados tienen la distribución de $X$ condicionada
a $X<1$, y sin conocer el mecanismo de falta no hay forma de recuperar
$\E[X]$ a partir de ellos: imputar con su media hereda el sesgo.
\end{ejemplo}

\begin{definicion}[Imputación simple e indicador]
La \emph{imputación por la media} (respectivamente, por la mediana o la
moda) reemplaza cada valor faltante de una columna por la media
(respectivamente, la mediana o el valor más frecuente) de los valores
observados de esa columna en entrenamiento. La \emph{indicadora de falta}
es la columna $(r_1,\dots,r_n)$ añadida como característica.
\end{definicion}

\begin{proposicion}[Efecto de la imputación por la media]\label{prop:imput}
Sean $z\in\R^n$ una columna con $n_o\ge2$ valores observados (índices en
$O$) y $y\in\R^n$ otra columna completa. Sea $\tilde z$ la columna
imputada con $\bar z_o=n_o^{-1}\sum_{k\in O}z_k$. Con varianzas y
covarianzas muestrales de divisor $n-1$ para las columnas completas y
$n_o-1$ para las calculadas en $O$:
\begin{enumerate}[label=(\alph*)]
\item $\overline{\tilde z}=\bar z_o$;
\item $s^2(\tilde z)=\frac{n_o-1}{n-1}\,s_o^2(z)$;
\item $s(\tilde z,y)=\frac{n_o-1}{n-1}\,s_o(z,y)$, donde $s_o(z,y)$ es la
covarianza de $z$ e $y$ calculada sobre $O$;
\item $\Corr(\tilde z,y)=\sqrt{\frac{n_o-1}{n-1}}\;\frac{s_o(y)}{s(y)}\;\Corr_o(z,y)$,
donde $\Corr_o$ y $s_o(y)$ se calculan sobre $O$ y se supone $s_o(z)>0$.
\end{enumerate}
\end{proposicion}

\begin{proof}
(a) $\sum_k\tilde z_k=\sum_{k\in O}z_k+(n-n_o)\bar z_o=n_o\bar
z_o+(n-n_o)\bar z_o$. (b) Por (a), las filas imputadas valen exactamente
la media, así que $\sum_k(\tilde z_k-\bar z_o)^2=\sum_{k\in O}(z_k-\bar
z_o)^2=(n_o-1)s_o^2(z)$; se divide por $n-1$. (c) Por lo mismo,
$\sum_k(\tilde z_k-\bar z_o)(y_k-\bar y)=\sum_{k\in O}(z_k-\bar
z_o)(y_k-\bar y)$. Como $\sum_{k\in O}(z_k-\bar z_o)=0$, reemplazar
$\bar y$ por la media $\bar y_o$ de $y$ sobre $O$ no altera la suma, que
vale $(n_o-1)s_o(z,y)$; se divide por $n-1$. (d) Se divide (c) por
$s(\tilde z)s(y)$, con $s(\tilde z)$ tomada de (b), y se escribe
$s_o(z,y)=\Corr_o(z,y)s_o(z)s_o(y)$.
\end{proof}

La imputación por la media conserva la media pero contrae la varianza y la
covarianza por el factor $(n_o-1)/(n-1)$, sea cual sea el mecanismo. Bajo
MCAR, las filas de $O$ son una muestra aleatoria de la población, $s_o(y)$
y $s(y)$ estiman la misma desviación estándar, y (d) muestra que la
correlación queda multiplicada por un factor cercano a
$\sqrt{n_o/n}$: con 30\,\% de faltantes, por $\approx0.84$. En el ejemplo
del capítulo, la columna $(1,\text{falta},5)$ se imputa como $(1,3,5)$, con
varianza $4$ frente a la varianza $8$ de los observados:
$(2-1)/(3-1)=1/2$.

\begin{proposicion}[La indicadora absorbe la media del grupo con faltantes]\label{prop:indic}
Sea $\tilde z$ la columna imputada con una constante $a$ y considérese la
regresión por mínimos cuadrados de $y$ sobre $\1$, $\tilde z$ y la
indicadora $r$, con al menos una fila con $r_k=1$. El valor ajustado en
cada fila con $r_k=1$ es la media de $y$ sobre esas filas, cualquiera que
sea el mecanismo de falta.
\end{proposicion}

\begin{proof}
Por el Lema~\ref{lem:proy}, el residuo $y-\hat y$ es ortogonal a cada
columna del diseño, en particular a $r$: $\sum_kr_k(y_k-\hat y_k)=0$. En
las filas con $r_k=1$ se tiene $\tilde z_k=a$, luego
$\hat y_k=\hat\beta_0+\hat\beta a+\hat\gamma$ es la misma constante para
todas ellas; que la suma de sus residuos sea cero obliga a que esa
constante sea la media de sus $y_k$.
\end{proof}

Bajo MNAR, la respuesta media del grupo con faltantes puede diferir de la
de cualquier individuo con $z=a$; la indicadora permite al modelo estimarla
directamente, al precio de una columna por característica afectada.

Para comparar con la imputación por la mediana conviene el caso de una
constante cualquiera.

\begin{proposicion}\label{prop:imput-a}
Con la notación de la Proposición~\ref{prop:imput}, sea $\tilde z^{(a)}$ la
columna imputada con una constante $a$. Entonces
\[
(n-1)\,s^2\bigl(\tilde z^{(a)}\bigr)=(n_o-1)s_o^2(z)+\frac{n_o(n-n_o)}{n}(\bar z_o-a)^2 .
\]
\end{proposicion}

\begin{proof}
La media de $\tilde z^{(a)}$ es $\tilde m=(n_o\bar z_o+(n-n_o)a)/n$, de
donde $\bar z_o-\tilde m=\frac{n-n_o}n(\bar z_o-a)$ y
$a-\tilde m=-\frac{n_o}n(\bar z_o-a)$. Descomponiendo alrededor de $\bar
z_o$, $\sum_{k\in O}(z_k-\tilde m)^2=(n_o-1)s_o^2+n_o(\bar z_o-\tilde
m)^2$ (el término cruzado se anula porque $\sum_{k\in O}(z_k-\bar
z_o)=0$), y las $n-n_o$ filas imputadas aportan $(n-n_o)(a-\tilde m)^2$.
Sumando,
$n_o\frac{(n-n_o)^2}{n^2}(\bar z_o-a)^2+(n-n_o)\frac{n_o^2}{n^2}(\bar
z_o-a)^2=\frac{n_o(n-n_o)}{n}(\bar z_o-a)^2$.
\end{proof}

Con $a=\bar z_o$ se recupera la Proposición~\ref{prop:imput}(b). Con otra
constante, por ejemplo la mediana de una columna asimétrica, aparece un
término adicional que no proviene de variación real de los datos: los
valores imputados siguen siendo idénticos entre sí y no contienen
información sobre $y$, pero desplazan la media y añaden varianza.

\paragraph{Contraste.} La media imputada conserva la media observada
(Proposición~\ref{prop:imput}(a)) pero, como la media tiene punto de
ruptura $1/n$ (Proposición~\ref{prop:ruptura}), un solo atípico entre los
observados puede desplazar el valor imputado arbitrariamente; la mediana
tiene punto de ruptura cercano a $1/2$. La moda es la única de las tres
definida para categóricas. Toda imputación por una constante asigna el
mismo valor a filas cuyas respuestas pueden ser distintas, y por eso
atenúa la relación con $y$ (Proposición~\ref{prop:imput}(d) para la
media). Bajo MAR, la imputación condicional a $W$ tiene la justificación de la
Proposición~\ref{prop:mcar}(b); bajo MNAR, ninguna imputación calculada con los
observados tiene garantía de insesgadez (Ejemplo~\ref{ej:mnar}), y la
indicadora permite al menos que el modelo estime por separado la respuesta
del grupo con faltantes.
Algunas implementaciones de \emph{gradient boosting} (XGBoost, LightGBM,
\texttt{HistGradientBoosting} de scikit-learn) aceptan faltantes: en cada
corte eligen con los datos de entrenamiento hacia qué hijo enviarlos, lo
que equivale a aprender la indicadora dentro del árbol.

\begin{aws}
\textbf{Data Wrangler}, grupo \emph{Handle missing}: \emph{Impute} (media
o mediana para numéricas, valor más frecuente para categóricas),
\emph{Add indicator for missing} (la indicadora de la
Proposición~\ref{prop:indic}), \emph{Fill missing} (constante dada) y
\emph{Drop missing}; es un paso con parámetros ajustados, con
\emph{refit}. \textbf{Glue DataBrew}: acciones
\texttt{FILL\_WITH\_AVERAGE}, \texttt{FILL\_WITH\_MEDIAN},
\texttt{FILL\_WITH\_MOST\_FREQUENT}, entre otras. Su documentación indica
que el estadístico se calcula sobre todos los valores de la columna del
dataset al que se aplica la receta; si la receta se ejecuta por separado
sobre prueba, usa la mediana de prueba, lo que viola la separación de la
Proposición~\ref{prop:test}. Para respetarla se calcula el valor en
entrenamiento y se rellena con ese valor fijo.
\begin{lstlisting}
{"RecipeAction": {"Operation": "FILL_WITH_MEDIAN",
                  "Parameters": {"sourceColumn": "age"}}}
\end{lstlisting}
\textbf{SageMaker Processing}: \texttt{SimpleImputer(strategy=...,
add\_indicator=True)} de scikit-learn, con \texttt{fit} sobre
entrenamiento; la imputación condicional de la
Proposición~\ref{prop:mcar}(b) corresponde a \texttt{IterativeImputer}.
Para series de tiempo, Data Wrangler ofrece además relleno hacia adelante,
hacia atrás e interpolación.
\end{aws}

\subsection{Valores atípicos}\label{sec:atipicos}

Una \emph{regla de detección} asigna a cada columna $z\in\R^n$ el
conjunto de índices que declara atípicos. Las reglas del capítulo son de
umbral: declaran atípico a $z_k$ cuando una puntuación calculada con $z_k$
y con estadísticos de la columna supera un valor fijo. Se comparan por dos
propiedades: si pueden detectar algo con el tamaño de muestra disponible, y
si el atípico altera los estadísticos con que se le juzga.

\subsubsection{Regla de la puntuación $z$}

\begin{definicion}
La puntuación $z$ de $z_k$ es $\zeta_k=(z_k-\bar z)/s$, con $s>0$. La regla
declara atípico a $z_k$ si $|\zeta_k|>c$, típicamente con $c=3$.
\end{definicion}

El valor $c=3$ se elige porque si $Z\sim N(0,1)$, $\Prob(|Z|>3)\approx0.0027$.
Pero $\bar z$ y $s$ se calculan con la propia observación que se juzga, y
eso limita la regla.

\begin{proposicion}[Cota de Shiffler]\label{prop:shiffler}
Para toda columna con $n\ge2$ y $s>0$, y todo $k$,
$|\zeta_k|\le\sqrt{n-1}$ si $s=s_n$, y $|\zeta_k|\le(n-1)/\sqrt n$ si
$s=s_{n-1}$.
\end{proposicion}

\begin{proof}
Sea $d_i=z_i-\bar z$. Como $\sum_id_i=0$, $d_k=-\sum_{i\ne k}d_i$, y por
la desigualdad de Cauchy--Schwarz aplicada a los vectores $(d_i)_{i\neq
k}$ y $(1,\dots,1)\in\R^{n-1}$, $d_k^2\le(n-1)\sum_{i\ne k}d_i^2$. Luego
$\sum_id_i^2\ge d_k^2\bigl(1+\frac1{n-1}\bigr)=d_k^2\frac n{n-1}$. Si
$s=s_n$, $\sum_id_i^2=ns^2$ y resulta $\zeta_k^2\le n-1$; si $s=s_{n-1}$,
$\sum_id_i^2=(n-1)s^2$ y resulta $\zeta_k^2\le(n-1)^2/n$.
\end{proof}

La igualdad en Cauchy--Schwarz ocurre cuando todos los $d_i$, $i\ne k$,
son iguales, así que la cota se alcanza. Consecuencia: la regla con $c=3$
no puede declarar atípico a nadie si $\sqrt{n-1}\le3$, es decir, si
$n\le10$ con $s_n$ (el divisor de \texttt{scipy.stats.zscore}), ni si
$(n-1)/\sqrt n\le3$, es decir, $n\le11$, con $s_{n-1}$. Esta es la razón
aritmética de que el ejemplo del capítulo ($n=10$) no detecte nada. (La
nota del texto base atribuye la cota $(n-1)/\sqrt n$ al divisor $n$;
corresponde al divisor $n-1$. Su conclusión se mantiene.)

El segundo defecto es el \emph{enmascaramiento}: el atípico infla $s$ y
reduce su propia puntuación. En el ejemplo, el valor $23$ tiene
$\zeta=2.58$ con $\bar z$ y $s_n$ calculados con las diez observaciones,
pero $\zeta=5.31$ si se calculan con las otras nueve.

\subsubsection{Regla del rango intercuartílico}

\begin{definicion}[Vallas de Tukey]
Se declara atípico a $z_k$ si $z_k<Q_1-1.5\IQR$ o $z_k>Q_3+1.5\IQR$.
\end{definicion}

Si $Z\sim N(\mu,\sigma^2)$, los cuartiles poblacionales son
$\mu\pm0.6745\sigma$ (porque $\Phi^{-1}(0.75)=0.6745$), el $\IQR$ poblacional
es $1.349\sigma$ y la valla superior está en $\mu+2.698\sigma$; la
probabilidad de quedar fuera de las vallas es
$2(1-\Phi(2.698))\approx0.0070$. La regla equivale entonces, bajo
normalidad, a un umbral de $2.7$ desviaciones con parámetros calculados de
otro modo. La diferencia con la regla $z$ es que esos parámetros no se
mueven con el atípico:

\begin{proposicion}[Ausencia de enmascaramiento]\label{prop:tukey}
Sea $n\ge5$ y supóngase que $z_{(n)}>z_{(n-1)}$. Si se reemplaza el
máximo por cualquier valor mayor, $Q_1$, $Q_3$ y las vallas no cambian. En
particular, al aumentar el máximo acaba siendo declarado atípico.
\end{proposicion}

\begin{proof}
Para $q=3/4$, $h=\frac34(n-1)+1$ y $r=\lfloor h\rfloor$. Si $n\ge5$,
$h\le n-\frac14(n-1)\le n-1$, de modo que $r+1\le n$ y, si $r+1=n$, es
porque $h=n-1$ es entero y el peso $h-r$ de $z_{(r+1)}$ es $0$. En todos
los casos $Q_3$ depende solo de $z_{(1)},\dots,z_{(n-1)}$, que no cambian
cuando el máximo aumenta (sigue siendo el máximo). Lo mismo vale para
$Q_1$, con índices menores. Las vallas quedan fijas y el máximo termina
por superarlas.
\end{proof}

En el ejemplo del capítulo, $Q_1=6.375$ y $Q_3=11.75$, así que la valla
superior es $19.81$ y el valor $23$ queda fuera; con $100$ o $1000$ en su
lugar, las vallas son las mismas.

\subsubsection{Puntuación $z$ modificada}

\begin{definicion}[Iglewicz y Hoaglin, 1993]
Con $\MAD(z)=\med\bigl(|z_1-\med(z)|,\dots,|z_n-\med(z)|\bigr)$, la
puntuación $z$ modificada es $M_k=0.6745\,(z_k-\med(z))/\MAD(z)$. Los
autores sugieren declarar atípicos los $|M_k|>3.5$.
\end{definicion}

El factor $0.6745$ hace $M_k$ comparable con $\zeta_k$: si $Z\sim
N(\mu,\sigma^2)$, la mediana poblacional de $|Z-\mu|$ es el número $a$ con
$\Prob(|Z-\mu|\le a)=1/2$, es decir, $2\Phi(a/\sigma)-1=1/2$, $a=0.6745\sigma$;
así, $\MAD/0.6745$ estima $\sigma$. En el ejemplo, $\med=8.25$, $\MAD=3$ y
$M=3.32$ para el valor $23$: supera el umbral $3$ (el que usa DataBrew
por defecto) y no el $3.5$. Los umbrales son convenciones de los autores
y de las herramientas; ninguno es consecuencia de un resultado.

\subsubsection{Punto de ruptura}\label{sec:ruptura}

La sensibilidad de un estadístico al atípico se mide con el punto de
ruptura (Donoho y Huber, 1983).

\begin{definicion}
Sea $T:\R^n\to\R$ un estadístico y $z\in\R^n$. Su \emph{punto de ruptura}
en $z$ es $\varepsilon^*(T,z)=m^*/n$, donde $m^*$ es el menor número de
coordenadas de $z$ que, reemplazadas por valores arbitrarios, permiten
hacer $|T|$ tan grande como se quiera.
\end{definicion}

\begin{lema}[Estadísticos de orden bajo contaminación]\label{lem:orden}
Sea $1\le r\le n$ y sea $z'$ obtenido de $z$ reemplazando $m$ coordenadas.
Si $m\le\min(r-1,\,n-r)$, entonces $z'_{(r)}$ está entre el mínimo y el
máximo de las coordenadas de $z$ no reemplazadas. Si $m\ge n-r+1$, $z'_{(r)}$
puede hacerse arbitrariamente grande; si $m\ge r$, arbitrariamente
negativo.
\end{lema}

\begin{proof}
Sean $a$ y $b$ el mínimo y el máximo de las $n-m$ coordenadas no
reemplazadas. Hay al menos $n-m\ge r$ coordenadas de $z'$ menores o iguales
que $b$, de modo que el $r$-ésimo menor valor es $\le b$. Hay a lo sumo
$m\le r-1$ coordenadas menores que $a$, de modo que el $r$-ésimo menor
valor es $\ge a$. Recíprocamente, si $m\ge n-r+1$ y los valores
reemplazados se fijan en $M>b$, hay a lo sumo $n-m\le r-1$ coordenadas
menores que $M$, luego $z'_{(r)}=M$; el caso $m\ge r$ es simétrico.
\end{proof}

\begin{proposicion}\label{prop:ruptura}
Para todo $z$: la media, $s_n$, $s_{n-1}$, el máximo, el mínimo y
$\max_k|z_k|$ tienen punto de ruptura $1/n$; la mediana tiene punto de
ruptura $\lceil n/2\rceil/n$, que tiende a $1/2$; $Q_1$, $Q_3$ y el $\IQR$
tienen punto de ruptura que tiende a $1/4$ cuando $n\to\infty$.
\end{proposicion}

\begin{proof}
Reemplácese $z_1$ por $M$. La nueva media es
$\bar z'=(M+\sum_{k\ge2}z_k)/n\to\infty$ cuando $M\to\infty$. Como la
suma de cuadrados contiene el sumando $(M-\bar z')^2$ y
$M-\bar z'=\frac{n-1}nM-\frac1n\sum_{k\ge2}z_k\to\infty$, se tiene
$s_n\ge|M-\bar z'|/\sqrt n\to\infty$, y lo mismo para $s_{n-1}\ge s_n$.
El máximo y $\max_k|z_k|$ son $\ge M$; con $-M$ en lugar de $M$, el mínimo
es $\le-M$. Así, $m^*=1$.

Por la definición del cuantil, $Q(q)$ es una combinación convexa de
$z_{(r)}$ y $z_{(r+1)}$ con $r=\lfloor(n-1)q+1\rfloor$, y está acotado
si ambos lo están. Para la mediana con $n$ impar, $Q(1/2)=z_{((n+1)/2)}$ y
el Lema~\ref{lem:orden} con $r=(n+1)/2$ da $m^*=n-r+1=(n+1)/2=\lceil
n/2\rceil$. Con $n$ par, $Q(1/2)=\frac12(z_{(n/2)}+z_{(n/2+1)})$; el lema
da acotamiento para $m\le n/2-1$ y explosión de $z_{(n/2+1)}$ con
$m=n/2$, así que $m^*=n/2$. Para $Q_3$, $r=\lfloor\frac34(n-1)+1\rfloor$ y
la explosión requiere $m\ge n-r$ o $n-r+1$ según el peso de
$z_{(r+1)}$; en ambos casos $m^*/n\to1/4$. Para $Q_1$, con
$r=\lfloor\frac14(n-1)+1\rfloor$, la explosión hacia $-\infty$ requiere
$m\ge r$, y $r/n\to1/4$. Si ambos cuartiles quedan acotados, su
diferencia también; y como $\IQR=Q_3-Q_1\ge0$, el $\IQR$ crece sin límite
cuando $Q_3\to\infty$ con $Q_1$ acotado o $Q_1\to-\infty$ con $Q_3$
acotado. Su punto de ruptura es entonces el menor de los de los cuartiles,
que tiende a $1/4$.
\end{proof}

\paragraph{Tratamientos y contraste.} Ante una observación declarada
atípica caben: eliminar la fila (correcto si el valor es un error de
registro, porque entonces no describe la población); \emph{recortar}
(\emph{winsorizar}), reemplazando $z_k$ por $\min(\max(z_k,l),u)$ con
$l,u$ las vallas; imputar como si faltara; o transformar toda la columna
(\S\ref{sec:log}). La comparación entre reglas se sigue de los resultados
anteriores: la regla $z$ requiere $n>c^2+1$ para poder detectar algo
(Proposición~\ref{prop:shiffler}) y sufre enmascaramiento porque sus
parámetros tienen punto de ruptura $1/n$; las vallas de Tukey y la $z$
modificada usan estadísticos con punto de ruptura $1/4$ y $1/2$ y no
sufren enmascaramiento por un solo valor
(Proposición~\ref{prop:tukey}). Ninguna regla aplicada columna por
columna detecta atípicos conjuntos: en los puntos $(-2,-2)$, $(-1,-1)$,
$(0,0)$, $(1,1)$, $(2,2)$, $(2,-2)$, el último rompe la relación $x_2=x_1$,
pero su valor $2$ en la primera columna también lo tiene $(2,2)$ y su valor
$-2$ en la segunda también lo tiene $(-2,-2)$; una regla de umbral que
marcara esos valores marcaría también esas filas normales.

\begin{aws}
\textbf{Glue DataBrew}: \texttt{FLAG\_OUTLIERS},
\texttt{REMOVE\_OUTLIERS}, \texttt{REPLACE\_OUTLIERS},
\texttt{RESCALE\_OUTLIERS\_WITH\_Z\_SCORE} y
\texttt{RESCALE\_OUTLIERS\_WITH\_SKEW}, con \texttt{outlierStrategy} igual
a \texttt{Z\_SCORE} ($|\zeta_k|$), \texttt{MODIFIED\_Z\_SCORE} ($|M_k|$) o
\texttt{IQR} (vallas con el múltiplo dado), y \texttt{threshold} con valor
por defecto $3$:
\begin{lstlisting}
{"Action": {"Operation": "FLAG_OUTLIERS",
  "Parameters": {"sourceColumn": "energia", "targetColumn": "energia_atipico",
                 "outlierStrategy": "IQR", "threshold": "1.5",
                 "trueString": "Yes", "falseString": "No"}}}
\end{lstlisting}
\textbf{Data Wrangler}, grupo \emph{Handle outliers}: \emph{Standard
deviation numeric outliers} (regla $z$), \emph{Robust standard deviation
numeric outliers} (centro y escala calculados entre dos cuantiles),
\emph{Quantile numeric outliers} y \emph{Min-max numeric outliers}
(umbrales dados), cada uno con corrección \emph{Clip} (recortar),
\emph{Remove} o \emph{Invalidate} (convertir en faltante). Para series de
tiempo, el análisis \emph{Anomaly detection} descompone la serie en
tendencia, estacionalidad y residuo y marca los residuos que superan un
múltiplo dado de su desviación estándar.
\end{aws}

% =====================================================================
\section{Características numéricas: escalamiento lineal}
% =====================================================================

\subsection{Definición y casos}

\begin{definicion}[Escalamiento lineal]\label{def:escala}
Sea $z\in\R^n$ la columna de entrenamiento de una característica
numérica. Un \emph{escalamiento lineal} es la transformación ajustada
$T_{a,b}(u)=(u-a)/b$, $u\in\R$, con parámetros $a\in\R$ (centro) y $b>0$
(escala) calculados con $z$. Los casos del capítulo son:
\begin{center}
\begin{tabular}{@{}lll@{}}
\toprule
Nombre & $a$ & $b$\\
\midrule
Estandarización (puntuación $z$) & $\bar z$ & $s_n(z)$\\
Normalización MinMax & $z_{(1)}$ & $z_{(n)}-z_{(1)}$\\
Escalado robusto & $\med(z)$ & $\IQR(z)$\\
Escalado MaxAbs & $0$ & $\max_k|z_k|$\\
Normalización por la media & $\bar z$ & $z_{(n)}-z_{(1)}$\\
\bottomrule
\end{tabular}
\end{center}
\end{definicion}

La condición $b>0$ tiene dos funciones. Excluye $b=0$, que haría $T$
indefinida; y, frente a $b<0$, garantiza que $T$ sea creciente y conserve
el orden de los valores. No todos los casos la cumplen en las mismas
columnas: $s_n(z)>0$ y $z_{(n)}-z_{(1)}>0$ si y solo si $z$ no es
constante, y $\max|z_k|>0$ si y solo si $z$ no es nula; en cambio, el
$\IQR$ puede anularse en columnas no constantes.

\begin{contraejemplo}
Sea $z=(0,0,0,0,0,0,0,5)$, $n=8$. Para $q=1/4$, $h=2.75$ y
$Q_1=z_{(2)}+0.75(z_{(3)}-z_{(2)})=0$; para $q=3/4$, $h=6.25$ y
$Q_3=z_{(6)}+0.25(z_{(7)}-z_{(6)})=0$. Luego $\IQR=0$ aunque la columna
no es constante. En general, por la definición del cuantil, $\IQR=0$ si
y solo si todos los estadísticos de orden que intervienen en $Q_1$ y
$Q_3$ (y por tanto los intermedios) son iguales: un mismo valor ocupa las
posiciones centrales de la muestra ordenada, como en conteos con muchos
ceros. scikit-learn reemplaza entonces la escala por $1$.
\end{contraejemplo}

Sobre la muestra de entrenamiento, MinMax toma valores en $[0,1]$ y MaxAbs
en $[-1,1]$, porque $z_{(1)}\le z_k\le z_{(n)}$ y $|z_k|\le\max_i|z_i|$.
Aplicada a un valor nuevo $u$ fuera de $[z_{(1)},z_{(n)}]$, MinMax da un
resultado fuera de $[0,1]$: los parámetros son los de entrenamiento
(Proposición~\ref{prop:test}). La normalización por la media da media $0$
y valores en $[-1,1]$ (porque $|z_k-\bar z|\le z_{(n)}-z_{(1)}$), pero su
desviación estándar es $s_n/(z_{(n)}-z_{(1)})$, que no es $1$ en general.

\subsection{Qué cambia y qué no cambia}

Para $z\in\R^n$ sean los momentos centrales $m_r(z)=n^{-1}\sum_k(z_k-\bar
z)^r$, el \emph{coeficiente de asimetría} $\gamma_1(z)=m_3/m_2^{3/2}$ y la
\emph{curtosis} $\gamma_2(z)=m_4/m_2^2$ (definidos si $m_2>0$).

\begin{proposicion}[Invariancia de la forma]\label{prop:forma}
Si $b>0$, $\gamma_1(T_{a,b}(z))=\gamma_1(z)$ y
$\gamma_2(T_{a,b}(z))=\gamma_2(z)$, donde $T_{a,b}(z)$ se aplica
coordenada a coordenada. Además, la correlación de Pearson entre
$T_{a,b}(z)$ y cualquier otra columna $w$ coincide con la de $z$ y $w$.
\end{proposicion}

\begin{proof}
La media de $T_{a,b}(z)$ es $(\bar z-a)/b$, así que sus desviaciones
respecto de la media son $(z_k-\bar z)/b$ y $m_r(T_{a,b}(z))=b^{-r}m_r(z)$.
Entonces
$\gamma_1(T_{a,b}(z))=b^{-3}m_3/(b^{-2}m_2)^{3/2}=m_3/m_2^{3/2}$, donde se
usó $(b^{-2})^{3/2}=b^{-3}$, válido porque $b>0$; para $\gamma_2$,
$b^{-4}m_4/(b^{-2}m_2)^2=m_4/m_2^2$. La covarianza de $T_{a,b}(z)$ con $w$
es la de $z$ con $w$ multiplicada por $1/b$, y la desviación estándar de
$T_{a,b}(z)$ es la de $z$ multiplicada por $1/b$; el factor se cancela en
el cociente.
\end{proof}

Si $b<0$ el mismo cálculo da $\gamma_1(T_{a,b}(z))=-\gamma_1(z)$: la
hipótesis $b>0$ es necesaria para la invariancia de la asimetría. La
proposición es el fundamento de la afirmación del capítulo de que ningún
escalado lineal corrige la asimetría.

\begin{proposicion}[Compresión por un valor extremo en MinMax]\label{prop:minmax}
Fíjense $z_1,\dots,z_{n-1}$ y sea $z_n=M$. Cuando $M\to\infty$, la
normalización MinMax de cada $z_k$, $k<n$, tiende a $0$.
\end{proposicion}

\begin{proof}
Si $M>\max_{k<n}z_k$, entonces $z_{(n)}=M$ y $z_{(1)}=\min_{k<n}z_k$ no
depende de $M$; el valor normalizado $(z_k-z_{(1)})/(M-z_{(1)})$ tiene
numerador fijo y denominador que tiende a infinito.
\end{proof}

En la columna de energía del capítulo, sin el valor $23$ los otros nueve
ocupan $[0,1]$ tras MinMax; con él, ocupan $[0,0.426]$. El fenómeno se
explica con la Proposición~\ref{prop:ruptura}: los parámetros de MinMax,
MaxAbs y la estandarización tienen punto de ruptura $1/n$, los del
escalado robusto $\approx1/4$. El escalado robusto no elimina el valor
extremo: por ser creciente, lo conserva como el mayor (en el ejemplo vale
$2.74$ y el siguiente $0.79$); lo que protege es la escala a la que quedan
los demás.

\begin{definicion}
Una matriz $A\in\R^{n\times p}$ es \emph{dispersa} cuando la mayoría de sus
entradas son nulas; un \emph{formato disperso} almacena solo las entradas
no nulas con sus posiciones, con memoria proporcional a su número
$\mathrm{nnz}(A)$ en lugar de a $np$.
\end{definicion}

Un ejemplo es la matriz documento-término de \S\ref{sec:tfidf}, donde la
entrada $(d,t)$ cuenta las apariciones del término $t$ en el documento
$d$: cada documento usa una fracción ínfima del vocabulario.

\begin{proposicion}[Conservación de la dispersión]\label{prop:sparse}
$T_{a,b}(0)=0$ si y solo si $a=0$. Si $a\ne0$ y la columna tiene $n_0$
ceros, la columna transformada tiene $n_0$ entradas iguales a $-a/b\neq0$.
En la Definición~\ref{def:escala}, solo MaxAbs tiene $a=0$ para toda
columna.
\end{proposicion}

\begin{proof}
$T_{a,b}(0)=-a/b$, nulo si y solo si $a=0$ porque $b>0$. En los demás
casos de la tabla, $a$ es la media, el mínimo o la mediana, que en general
son no nulos (por ejemplo, la media de una columna no negativa y no nula es
positiva).
\end{proof}

Con $10^6$ documentos, $10^5$ términos y 100 términos distintos por
documento, la matriz tiene $10^8$ entradas no nulas; centrarla produce
$10^{11}$. Dividir por $s_n$ sin centrar ($a=0$,
\texttt{StandardScaler(with\_mean=False)}) también conserva la
dispersión.

\subsection{Escala y descenso de gradiente}

El capítulo afirma que escalar acelera la convergencia del descenso de
gradiente. La afirmación se precisa para mínimos cuadrados.

\begin{proposicion}\label{prop:cond}
Sea $f(\beta)=\frac1{2n}\|y-X\beta\|^2$, $X\in\R^{n\times p}$, y
supóngase que $H=X^\top X/n$ es definida positiva, con autovalores
$\lambda_{\max}\ge\dots\ge\lambda_{\min}>0$; sea
$\kappa=\lambda_{\max}/\lambda_{\min}$ (\emph{número de condición}). Entonces
$f$ tiene un único minimizador $\hat\beta$, y el descenso de gradiente
$\beta_{t+1}=\beta_t-\eta\nabla f(\beta_t)$ con paso $\eta=1/\lambda_{\max}$
cumple, para todo $\beta_0$,
$\|\beta_t-\hat\beta\|\le(1-1/\kappa)^t\|\beta_0-\hat\beta\|$.
\end{proposicion}

\begin{proof}
$\nabla f(\beta)=H\beta-X^\top y/n$; como $H$ es invertible, $\nabla f$ se
anula solo en $\hat\beta=H^{-1}X^\top y/n$, y como la hessiana $H$ es
definida positiva, $f$ es estrictamente convexa y $\hat\beta$ es su único
minimizador. Usando $X^\top y/n=H\hat\beta$, $\nabla
f(\beta)=H(\beta-\hat\beta)$, y
$\beta_{t+1}-\hat\beta=(I-\eta H)(\beta_t-\hat\beta)$. La matriz
$I-H/\lambda_{\max}$ es simétrica con autovalores
$1-\lambda_i/\lambda_{\max}\in[0,1-1/\kappa]$, de modo que su norma
espectral es $1-1/\kappa$ y
$\|\beta_{t+1}-\hat\beta\|\le(1-1/\kappa)\|\beta_t-\hat\beta\|$; se itera.
\end{proof}

Supóngase dos características con media cero y covarianza muestral nula,
de varianzas $s_1^2$ y $s_2^2$ (y la respuesta centrada, de modo que no
hace falta intercepto). Entonces $H=\operatorname{diag}(s_1^2,s_2^2)$ y
$\kappa=s_1^2/s_2^2$ si $s_1\ge s_2$. Si ambas características son
longitudes con la misma dispersión, pero una se registra en metros
($s_2=1$) y la otra en centímetros ($s_1=100$), $\kappa=10^4$ y la
cota garantiza una reducción del error de solo $1-10^{-4}$ por iteración.
Tras estandarizar, $H$ es la matriz de correlaciones muestrales; con dos
características de correlación $\rho$, sus autovalores son $1\pm\rho$ (los
de $\begin{psmallmatrix}1&\rho\\\rho&1\end{psmallmatrix}$, con
autovectores $(1,\pm1)$) y $\kappa=(1+|\rho|)/(1-|\rho|)$, que depende de
la redundancia entre las características y no de sus unidades.

\paragraph{Contraste.} Ningún escalamiento lineal cambia la forma de la
distribución (Proposición~\ref{prop:forma}); la elección depende de la
propiedad requerida: rango acotado conocido en entrenamiento (MinMax),
varianzas iguales para el condicionamiento y para que una penalización
trate igual a todos los coeficientes (estandarización,
Proposición~\ref{prop:cond} y Contraejemplo~\ref{ce:ridge}), parámetros
poco afectados por valores extremos (robusto, con la salvedad de columnas
con $\IQR=0$), conservación de la dispersión (MaxAbs,
Proposición~\ref{prop:sparse}). Para árboles y regresión lineal sin
penalizar el escalamiento no cambia las predicciones
(Proposiciones~\ref{prop:arbol} y~\ref{prop:ols}).

\begin{aws}
\textbf{Data Wrangler}, \emph{Process numeric} $\to$ \emph{Scale values},
con \emph{Standard scaler}, \emph{Robust scaler}, \emph{Min max scaler}
(rango configurable) y \emph{Max absolute scaler}; opera sobre una columna
por paso. \textbf{Glue DataBrew}, acción \texttt{SCALE} (en el JSON la
operación se llama \texttt{NORMALIZATION}) con \texttt{strategy} igual a
\texttt{MIN\_MAX}, \texttt{SCALE\_BETWEEN} (a un rango dado),
\texttt{MEAN\_NORMALIZATION} o \texttt{Z\_SCORE}:
\begin{lstlisting}
{"Action": {"Operation": "NORMALIZATION",
  "Parameters": {"sourceColumn": "all_votes", "strategy": "MIN_MAX",
                 "targetColumn": "all_votes_normalized"}}}
\end{lstlisting}
La documentación de DataBrew describe \texttt{MEAN\_NORMALIZATION} como
«media 0 y desviación 1 en $[-1,1]$»; con la fórmula de la
Definición~\ref{def:escala} la media es $0$ y los valores están en
$[-1,1]$, pero la desviación estándar no es $1$ en general. Como toda
acción de DataBrew, calcula sus estadísticos sobre el dataset al que se
aplica la receta. \textbf{SageMaker Processing}: \texttt{StandardScaler},
\texttt{RobustScaler}, \texttt{MinMaxScaler} y \texttt{MaxAbsScaler} de
scikit-learn, con \texttt{fit} sobre entrenamiento.
\end{aws}

% =====================================================================
\section{Características numéricas: transformaciones no lineales}
% =====================================================================

Las transformaciones de esta sección son funciones estrictamente
crecientes aplicadas a toda la columna. Por la
Proposición~\ref{prop:arbol} no alteran los árboles; su efecto recae sobre
modelos lineales, métodos de distancia y redes neuronales. Se estudian tres
propósitos distintos: linealizar relaciones, estabilizar la varianza y
reducir la asimetría.

\subsection{Transformación logarítmica}\label{sec:log}

\paragraph{Linealización.} Si $y=a\,e^{bx}\varepsilon$ con $a>0$ y un
error multiplicativo $\varepsilon>0$ con $\E[\log\varepsilon]=0$, entonces
$\log y=\log a+bx+\log\varepsilon$: un modelo lineal para $\log y$ con
error aditivo de media cero. Si $y=a\,x^b\varepsilon$ con $x>0$, la
relación es lineal entre $\log y$ y $\log x$. En el ejemplo del capítulo,
$y=10^{x-3}$ da $\log_{10}y=x-3$.

\paragraph{Estabilización de la varianza.} Muchas medidas positivas
(ingresos, tiempos, concentraciones) tienen desviación estándar
aproximadamente proporcional a su media. El modelo exacto de esa situación
es una familia de escala.

\begin{definicion}
Sea $G>0$ una variable aleatoria con distribución fija. La \emph{familia de
escala} generada por $G$ es la familia de distribuciones de $X=\theta G$,
$\theta>0$.
\end{definicion}

Son familias de escala la exponencial, la gamma con parámetro de forma fijo
y la lognormal con $\sigma$ fijo (si $\log G\sim N(0,\sigma^2)$, entonces
$\log(\theta G)\sim N(\log\theta,\sigma^2)$).

\begin{proposicion}\label{prop:escala}
Si $X=\theta G$ con $\E[G^2]<\infty$ y $\E[(\log G)^2]<\infty$, entonces
$\E[X]=\theta\E[G]$ y $\sqrt{\Var X}=\theta\sqrt{\Var G}$, de modo que el
cociente desviación/media no depende de $\theta$; y
$\Var(\log X)=\Var(\log G)$ no depende de $\theta$.
\end{proposicion}

\begin{proof}
Las dos primeras igualdades son la linealidad de la esperanza y
$\Var(\theta G)=\theta^2\Var G$. Para la tercera,
$\log X=\log\theta+\log G$, y sumar la constante $\log\theta$ no cambia la
varianza.
\end{proof}

Es decir: si varios grupos difieren solo en la escala, en la escala
original sus varianzas son proporcionales a sus medias al cuadrado, y tras
el logaritmo son iguales. Así, con $G$ gamma de forma $3$, $\Var(\log
X)=\psi'(3)\approx0.395$ para cualquier $\theta$ ($\psi'$ es la función
trigamma).

\paragraph{Asimetría.} Si $X$ es lognormal, $\log X$ es exactamente
normal, simétrica. Que el logaritmo reduzca la asimetría no es, sin
embargo, una propiedad general:

\begin{contraejemplo}\label{ce:logunif}
Si $U$ es uniforme en $(0,1)$ (asimetría $0$), entonces $-\log U$ es
exponencial de media $1$, porque $\Prob(-\log U>t)=\Prob(U<e^{-t})=e^{-t}$
para $t\ge0$. La exponencial de media $1$ tiene varianza $1$ y tercer
momento central $2$, luego asimetría $2$; por tanto $\log U$ tiene
asimetría $-2$. El logaritmo convirtió una variable simétrica en una muy
asimétrica hacia la izquierda.
\end{contraejemplo}

El logaritmo comprime más los valores grandes que los pequeños (su
derivada $1/x$ decrece), así que acorta colas derechas largas; aplicado a
una variable sin esa cola, estira la izquierda. Su uso se justifica en
variables con asimetría positiva o de la forma de la
Proposición~\ref{prop:escala}, y conviene comprobar la asimetría resultante.

\begin{proposicion}[Sesgo de retransformación]\label{prop:retrans}
Si $\log Y\sim N(\mu,\sigma^2)$ con $\sigma>0$, entonces
$\E[Y]=e^{\mu+\sigma^2/2}>e^{\mu}$, y $e^\mu$ es la mediana de $Y$.
\end{proposicion}

\begin{proof}
$\E[Y]=\E[e^W]$ con $W\sim N(\mu,\sigma^2)$ es la función generadora de
momentos de $W$, $M_W(t)=e^{\mu t+\sigma^2t^2/2}$, en $t=1$. Como $\exp$ es
creciente, $\Prob(Y\le e^\mu)=\Prob(W\le\mu)=1/2$.
\end{proof}

Si un modelo se ajusta para $\log y$ y se reporta $\exp$ de su predicción,
se estima la mediana de $Y$, no su media. Bajo el modelo lognormal, la
media se recupera multiplicando por $e^{\hat\sigma^2/2}$, con $\hat\sigma^2$
la varianza residual estimada.

\paragraph{Ceros.} El logaritmo exige valores positivos. El
desplazamiento $\log(1+x)$ (\texttt{log1p}) admite ceros, pero su efecto
depende de la unidad de medida de $x$: para $u\ge0$,
$0\le u-\log(1+u)\le u^2/2$ (la función $h(u)=u-\log(1+u)$ cumple
$h(0)=0$ y $0\le h'(u)=u/(1+u)\le u$). Si $x$ se mide en una unidad
$c$ veces mayor, $\log(1+x/c)$ difiere de la transformación lineal $x/c$
en a lo sumo $(x/c)^2/2$, que es despreciable cuando $x/c$ es pequeño:
con la unidad «grande», \texttt{log1p} es casi lineal y no reduce la
asimetría. La constante de desplazamiento debe elegirse en relación con la
escala de los datos.

\subsection{Raíces cuadrada y cúbica}\label{sec:raiz}

Para conteos, el modelo de referencia es Poisson, cuya varianza es igual a
su media; la raíz cuadrada estabiliza esa varianza. El resultado se
obtiene con el método delta.

\begin{proposicion}[Método delta]\label{prop:delta}
Sean $T_n$ variables aleatorias con
$\sqrt n(T_n-\mu)\xrightarrow{d}N(0,\sigma^2)$ y $g$ derivable en $\mu$.
Entonces
$\sqrt n\bigl(g(T_n)-g(\mu)\bigr)\xrightarrow{d}N\bigl(0,g'(\mu)^2\sigma^2\bigr)$.
\end{proposicion}

\begin{proof}
Sea $\rho(t)=\frac{g(t)-g(\mu)}{t-\mu}-g'(\mu)$ para $t\ne\mu$ y
$\rho(\mu)=0$; por la definición de derivada, $\rho$ es continua en $\mu$.
Se tiene
$\sqrt n(g(T_n)-g(\mu))=\bigl(g'(\mu)+\rho(T_n)\bigr)\sqrt n(T_n-\mu)$.
Como $\sqrt n(T_n-\mu)$ converge en distribución, está acotada en
probabilidad, y $T_n-\mu=n^{-1/2}\sqrt n(T_n-\mu)\to0$ en probabilidad;
por la continuidad de $\rho$ en $\mu$, $\rho(T_n)\to0$ en probabilidad. El
teorema de Slutsky da que el producto converge en distribución a
$g'(\mu)$ por una $N(0,\sigma^2)$.
\end{proof}

\begin{corolario}\label{cor:poisson}
Si $X_\mu\sim\text{Poisson}(\mu)$ con $\mu\in\N$, entonces
$\sqrt{X_\mu}-\sqrt\mu\xrightarrow{d}N(0,1/4)$ cuando $\mu\to\infty$.
\end{corolario}

\begin{proof}
La suma de variables de Poisson independientes es de Poisson con la suma
de los parámetros, así que $X_\mu$ tiene la distribución de
$\sum_{i=1}^\mu V_i$ con $V_i$ Poisson$(1)$ independientes, de media y
varianza $1$. Por el teorema central del límite,
$\sqrt\mu(X_\mu/\mu-1)\xrightarrow{d}N(0,1)$. La
Proposición~\ref{prop:delta} con $n=\mu$, $T_n=X_\mu/\mu$ y $g=\sqrt{\cdot}$
(con $g'(1)=1/2$) da $\sqrt\mu(\sqrt{X_\mu/\mu}-1)\xrightarrow{d}N(0,1/4)$,
y $\sqrt\mu\sqrt{X_\mu/\mu}=\sqrt{X_\mu}$.
\end{proof}

La varianza de $X_\mu$ es $\mu$; la de $\sqrt{X_\mu}$ se aproxima a $1/4$
sin importar $\mu$ (por simulación, $0.252$ con $\mu=50$ y $0.286$ con
$\mu=5$; Anscombe mostró que $\sqrt{x+3/8}$ mejora la aproximación con
$\mu$ pequeño). La raíz cúbica $x^{1/3}$ está definida en todo $\R$, es
impar y estrictamente creciente, así que se aplica a datos con signo sin
desplazamientos.

La intensidad con que estas funciones comprimen los valores grandes se
ordena por su crecimiento: para $0<\lambda_1<\lambda_2$,
$x^{\lambda_1}/x^{\lambda_2}=x^{\lambda_1-\lambda_2}\to0$, y para todo
$\lambda>0$, $\log x/x^\lambda\to0$ cuando $x\to\infty$ (por la regla de
L'Hôpital, el cociente de derivadas es $1/(\lambda x^\lambda)$). Así, para
valores grandes, $\log$ crece más despacio que $\sqrt[3]{\cdot}$, y esta
más despacio que $\sqrt{\cdot}$: es el orden de intensidad que da el
capítulo («escalera de potencias» de Tukey).

\subsection{Transformación de Box--Cox}

Las transformaciones anteriores se eligen de antemano. Box y Cox (1964)
plantearon elegir la potencia con los datos, dentro de una familia que
contiene a $\log$, $\sqrt{\cdot}$ y $1/x$ salvo constantes afines.

\begin{definicion}[Box--Cox]
Para $x>0$ y $\lambda\in\R$,
$x^{(\lambda)}=(x^\lambda-1)/\lambda$ si $\lambda\ne0$ y
$x^{(0)}=\log x$.
\end{definicion}

\begin{proposicion}\label{prop:bc-basicas}
Para cada $x>0$, $\lambda\mapsto x^{(\lambda)}$ es continua en
$\lambda=0$. Para cada $\lambda$, $x\mapsto x^{(\lambda)}$ es
estrictamente creciente en $(0,\infty)$.
\end{proposicion}

\begin{proof}
$(x^\lambda-1)/\lambda=(e^{\lambda\log x}-e^{0})/\lambda$ es el cociente
incremental en $0$ de la función $\lambda\mapsto e^{\lambda\log x}$, cuya
derivada en $0$ es $\log x$; luego el límite cuando $\lambda\to0$ es
$\log x=x^{(0)}$. La derivada respecto de $x$ es $x^{\lambda-1}>0$ (para
$\lambda=0$, $1/x$).
\end{proof}

Ambas propiedades justifican la forma de la definición: sin restar $1$ y
dividir por $\lambda$, $x^\lambda$ sería decreciente para $\lambda<0$
(invertiría el orden) y constante para $\lambda=0$, y la familia no
variaría con continuidad. La restricción $x>0$ es necesaria: $x^\lambda$
no es real para $x<0$ y $\lambda$ no entero, y $\log0$ no existe. Los
valores $\lambda=1$, $1/2$, $0$ y $-1$ corresponden, salvo transformaciones
afines, a la identidad, la raíz cuadrada, el logaritmo y el recíproco.

\begin{proposicion}[Estimación de $\lambda$]\label{prop:bc}
Supóngase como modelo de trabajo que existen $\lambda,\mu$ y $\sigma^2>0$
tales que $x_1^{(\lambda)},\dots,x_n^{(\lambda)}$ son i.i.d.\
$N(\mu,\sigma^2)$. La log-verosimilitud de $\lambda$, maximizada en
$(\mu,\sigma^2)$ para cada $\lambda$ (log-verosimilitud perfilada), es,
salvo una constante,
\[
\ell_p(\lambda)=-\frac n2\log\hat\sigma^2(\lambda)+(\lambda-1)\sum_{k=1}^n\log x_k,
\]
donde $\hat\sigma^2(\lambda)$ es la varianza, con divisor $n$, de los
$x_k^{(\lambda)}$. Si todos los $x_k$ se multiplican por una constante
$c>0$, $\ell_p$ cambia en la constante $-n\log c$ y su maximizador
$\hat\lambda$ no cambia.
\end{proposicion}

\begin{proof}
Por la Proposición~\ref{prop:bc-basicas}, $u\mapsto u^{(\lambda)}$ es
diferenciable y estrictamente creciente, con derivada $u^{\lambda-1}$; por
la fórmula de cambio de variable, la densidad de cada $x_k$ es la normal
evaluada en $x_k^{(\lambda)}$ multiplicada por $x_k^{\lambda-1}$. La
log-verosimilitud es
\[
-\frac n2\log(2\pi\sigma^2)-\frac1{2\sigma^2}\sum_k\bigl(x_k^{(\lambda)}-\mu\bigr)^2+(\lambda-1)\sum_k\log x_k .
\]
Para $\lambda$ fijo, es la log-verosimilitud normal de los
$x_k^{(\lambda)}$ más un término que no depende de $(\mu,\sigma^2)$, así
que se maximiza en la media y la varianza muestrales (divisor $n$) de los
$x_k^{(\lambda)}$; con esos valores, el término cuadrático vale $-n/2$, y
queda $\ell_p(\lambda)$ más constantes.

Para la invariancia: si $\lambda\ne0$,
$(cx)^{(\lambda)}=c^\lambda x^{(\lambda)}+(c^\lambda-1)/\lambda$, y si
$\lambda=0$, $\log(cx)=\log x+\log c$; en ambos casos una transformación
afín de $x^{(\lambda)}$ con pendiente $c^\lambda$, luego
$\hat\sigma_c^2(\lambda)=c^{2\lambda}\hat\sigma^2(\lambda)$. Además
$\sum_k\log(cx_k)=n\log c+\sum_k\log x_k$. Sustituyendo,
$\ell_{p,c}(\lambda)=\ell_p(\lambda)-n\lambda\log c+(\lambda-1)n\log
c=\ell_p(\lambda)-n\log c$.
\end{proof}

Dos observaciones delimitan el resultado. Primero, el modelo de trabajo
no puede cumplirse exactamente si $\lambda\ne0$: para $\lambda>0$,
$x^{(\lambda)}>-1/\lambda$, y una normal pone probabilidad positiva debajo
de cualquier número; Box y Cox lo presentan como una aproximación, y
$\hat\lambda$ se interpreta como la potencia que más acerca los datos a la
normalidad dentro de la familia. Segundo, el término
$(\lambda-1)\sum\log x_k$, que proviene del cambio de variable, es
indispensable: sin él se compararían varianzas medidas en escalas
distintas. Por ejemplo, si todos los $x_k>1$ y $\lambda\to-\infty$,
$x_k^\lambda\to0$ y $x_k^{(\lambda)}=-1/\lambda+x_k^\lambda/\lambda$, cuya
varianza $\Var(x^\lambda)/\lambda^2$ tiende a $0$: minimizar
$\hat\sigma^2(\lambda)$ sin el término elegiría siempre un $\lambda$
extremo. Por último, el estimador es invariante a la escala pero no a la
traslación: $\hat\lambda$ calculado para $x+c$ difiere en general, así que
desplazar datos para volverlos positivos cambia el resultado.

\begin{contraejemplo}\label{ce:gap}
Ninguna transformación de la familia (ni ninguna función continua y
estrictamente creciente) vuelve normal a una variable cuyo soporte tenga un
hueco. Sea $X$ con densidad positiva solo en $[1,2]\cup[10,11]$ y $g$
continua y estrictamente creciente. Entonces $g(X)\in[g(1),g(2)]\cup[g(10),g(11)]$
y $\Prob\bigl(g(2)<g(X)<g(10)\bigr)=\Prob(2<X<10)=0$, mientras que una
normal asigna probabilidad positiva al intervalo no vacío $(g(2),g(10))$.
\end{contraejemplo}

Si la función de distribución $F$ de $X$ es continua, la
\emph{transformación cuantil} $\Phi^{-1}\circ F$ sí produce una $N(0,1)$
exacta: para $u\in(0,1)$ sea $x_u$ el mayor $x$ con $F(x)=u$ (existe por
la continuidad de $F$ y porque $F(x)\to1$); como $F$ es no decreciente,
$F(X)\le u$ si y solo si $X\le x_u$, luego $\Prob(F(X)\le u)=F(x_u)=u$, es
decir, $F(X)$ es uniforme en $(0,1)$, y $\Prob(\Phi^{-1}(F(X))\le
t)=\Prob(F(X)\le\Phi(t))=\Phi(t)$. El precio es estimar $F$ completa (en
la práctica, la función de distribución empírica de entrenamiento), frente
a un solo parámetro $\lambda$.

\subsection{Transformación de Yeo--Johnson}

Yeo y Johnson (2000) extendieron la idea a datos con ceros y negativos,
que Box--Cox no admite.

\begin{definicion}[Yeo--Johnson]
Para $x\in\R$ y $\lambda\in\R$,
\[
\psi(\lambda,x)=\begin{cases}
\bigl((x+1)^\lambda-1\bigr)/\lambda,& x\ge0,\ \lambda\neq0,\\
\log(x+1),& x\ge0,\ \lambda=0,\\
-\bigl((1-x)^{2-\lambda}-1\bigr)/(2-\lambda),& x<0,\ \lambda\neq2,\\
-\log(1-x),& x<0,\ \lambda=2.
\end{cases}
\]
\end{definicion}

\begin{proposicion}
Para cada $\lambda$, $x\mapsto\psi(\lambda,x)$ es derivable en $\R$ con
derivada positiva: $(x+1)^{\lambda-1}$ si $x\ge0$ y $(1-x)^{1-\lambda}$ si
$x<0$. Además $\psi(\lambda,0)=0$, para $x\ge0$ coincide con la Box--Cox
de $x+1$, y $\psi(1,x)=x$.
\end{proposicion}

\begin{proof}
Las derivadas de cada rama se calculan directamente (también en los casos
logarítmicos: $1/(x+1)$ y $1/(1-x)$, que coinciden con las fórmulas en
$\lambda=0$ y $\lambda=2$); son positivas porque $x+1>0$ para $x\ge0$ y
$1-x>0$ para $x<0$. En $x=0$ ambas ramas valen $0$ y sus derivadas
laterales valen $1$, luego $\psi$ es derivable en $0$. La coincidencia
con Box--Cox es la definición de la rama $x\ge0$. Con $\lambda=1$:
$(x+1)-1=x$ si $x\ge0$ y $-((1-x)-1)=x$ si $x<0$.
\end{proof}

La estimación de $\lambda$ sigue la Proposición~\ref{prop:bc} con el
término de cambio de variable
$(\lambda-1)\sum_k\operatorname{sgn}(x_k)\log(|x_k|+1)$, que resulta de
sumar los logaritmos de las derivadas anteriores. En el ejemplo del
capítulo (datos simétricos), scikit-learn estima $\hat\lambda\approx0.92$ y
$1.02$, cercanos al valor $1$ de la identidad.

\paragraph{Contraste.}
\begin{center}
\begin{tabularx}{\textwidth}{@{}lYYY@{}}
\toprule
& Dominio & Parámetro & Situación que la justifica\\
\midrule
$\log x$ & $x>0$ & ninguno & familia de escala (Prop.~\ref{prop:escala}); relación multiplicativa\\
$\log(1+x/c)$ & $x>-c$ & $c$, elegido según la unidad & ídem con ceros\\
$\sqrt x$ & $x\ge0$ & ninguno & conteos de Poisson (Cor.~\ref{cor:poisson})\\
$\sqrt[3]{x}$ & $\R$ & ninguno & asimetría con valores de ambos signos\\
Box--Cox & $x>0$ & $\lambda$ por máxima verosimilitud & asimetría de grado desconocido\\
Yeo--Johnson & $\R$ & $\lambda$ por máxima verosimilitud & ídem, con ceros y negativos\\
$\Phi^{-1}\circ\hat F$ & $\R$ & la distribución empírica & normalidad aun con formas no alcanzables por potencias\\
\bottomrule
\end{tabularx}
\end{center}
Las transformaciones sin parámetros no pueden producir fuga; las de
potencia y la cuantil estiman parámetros y se ajustan con entrenamiento.
Todas cambian la interpretación de los coeficientes: en un modelo para
$\log y$, un coeficiente $b$ de $x$ significa que aumentar $x$ en una
unidad multiplica la mediana de $y$ por $e^b$ (bajo el modelo de la
Proposición~\ref{prop:retrans}).

\begin{aws}
\textbf{Glue DataBrew}, acción \texttt{SKEWNESS} con \texttt{skewFunction}
igual a \texttt{LOG} (base en \texttt{value}), \texttt{ROOT} (índice de la
raíz en \texttt{value}) o \texttt{SQUARE}:
\begin{lstlisting}
{"RecipeAction": {"Operation": "SKEWNESS",
  "Parameters": {"sourceColumn": "ingreso", "targetColumn": "log_ingreso",
                 "skewFunction": "LOG", "value": "2.718281828"}}}
\end{lstlisting}
DataBrew no tiene Box--Cox, Yeo--Johnson ni transformación cuantil.
\textbf{Data Wrangler} no tiene un paso dedicado: se usa \emph{Custom
formula}, con expresiones de Spark SQL como \texttt{log(x)},
\texttt{log1p(x)}, \texttt{sqrt(x)} o \texttt{cbrt(x)}, o \emph{Custom
transform} en pandas o PySpark (la versión de scikit-learn disponible ahí
incluye \texttt{PowerTransformer}). \textbf{SageMaker Processing}:
\texttt{PowerTransformer(method="box-cox"|"yeo-johnson")}, que estima
$\lambda$ por máxima verosimilitud y estandariza la salida por defecto, y
\texttt{QuantileTransformer(output\_distribution="normal")}.
\end{aws}

% =====================================================================
\section{Discretización}\label{sec:discretizacion}
% =====================================================================

\begin{definicion}
Dados puntos de corte $-\infty=t_0<t_1<\dots<t_{B-1}<t_B=+\infty$, la
\emph{discretización} asigna a $u\in\R$ el índice $b\in\{1,\dots,B\}$ tal
que $t_{b-1}<u\le t_b$. Los cortes interiores se calculan con la columna
de entrenamiento $z$: con \emph{ancho igual},
$t_b=z_{(1)}+\frac bB(z_{(n)}-z_{(1)})$; con \emph{frecuencia igual},
$t_b=Q(b/B)$.
\end{definicion}

Con cortes por cuantiles, cada intervalo contiene aproximadamente $n/B$
observaciones de entrenamiento; con ancho igual, los intervalos de las
colas pueden quedar casi vacíos si la distribución es asimétrica.

\begin{proposicion}\label{prop:bin-cuantil}
Supóngase que los valores de $z$ son distintos y sea $g$ estrictamente
creciente. La discretización por frecuencia igual de $g(z)$ asigna a cada
observación de entrenamiento el mismo índice que la de $z$. La
discretización por ancho igual no tiene esta propiedad.
\end{proposicion}

\begin{proof}
El corte $t_b=Q(b/B)$ es un estadístico de orden $z_{(r)}$ o un punto
estrictamente entre $z_{(r)}$ y $z_{(r+1)}$, con $r$ determinado solo por
$b$, $B$ y $n$; como los valores son distintos, en ambos casos
$\{k:z_k\le t_b\}$ es el conjunto de índices de rango $\le r$. Como $g$ es
estrictamente creciente, los rangos de $g(z)$ son los de $z$, y el corte
de $g(z)$ separa los mismos índices. Para ancho igual basta un ejemplo:
$z=(1,10,20,100)$ con $B=2$ corta en $50.5$ y agrupa $\{1,10,20\}$,
mientras que $\log_{10}z=(0,1,1.30,2)$ corta en $1$ y separa $\{1,10\}$ de
$\{20,100\}$.
\end{proof}

La discretización pierde información. El caso extremo se cuantifica.

\begin{proposicion}[Costo de dicotomizar]\label{prop:dicot}
Sea $X\sim N(0,1)$, $D=\ind{X>0}$ e $Y=\beta X+\varepsilon$ con
$\beta\ne0$ y $\varepsilon$ independiente de $X$ con varianza finita.
Entonces $\Corr(Y,D)=\sqrt{2/\pi}\;\Corr(Y,X)\approx0.80\,\Corr(Y,X)$.
\end{proposicion}

\begin{proof}
$\Cov(X,D)=\E[X\ind{X>0}]=\int_0^\infty x\varphi(x)\,dx=\varphi(0)=1/\sqrt{2\pi}$,
usando $\varphi'(x)=-x\varphi(x)$; $\Var D=1/4$ y $\Var X=1$, así que
$\Corr(X,D)=\frac{1/\sqrt{2\pi}}{1/2}=\sqrt{2/\pi}$. Como $\varepsilon$ es
independiente de $X$, también lo es de $D$, y
$\Cov(Y,D)=\beta\Cov(X,D)$ y $\Cov(Y,X)=\beta$. Entonces, con $\sigma_Y$ y
$\sigma_D$ las desviaciones estándar de $Y$ y de $D$,
$\Corr(Y,D)=\frac{\beta\Cov(X,D)}{\sigma_Y\sigma_D}
=\frac{\beta}{\sigma_Y}\Corr(X,D)=\Corr(Y,X)\Corr(X,D)$, pues
$\Corr(Y,X)=\beta/\sigma_Y$ (y $\sigma_X=1$, así que $\Corr(X,D)=\Cov(X,D)/\sigma_D$).
\end{proof}

Una lectura aproximada en términos de tamaño muestral: para probar
asociación con la correlación muestral $r$, se usa que $\sqrt nr$ es
aproximadamente $N(0,1)$ bajo independencia (con momentos cuartos
finitos), y bajo la alternativa su centro es aproximadamente $\sqrt
n\rho$. Mantener $\sqrt n\rho$ cuando $\rho$ se reduce por el factor
$\sqrt{2/\pi}$ exige multiplicar $n$ por $\pi/2\approx1.57$: dicotomizar
equivale, en esta aproximación, a descartar un 36\,\% de la muestra.

\begin{proposicion}[Error de aproximación]\label{prop:lipschitz}
Sea $f$ lipschitziana de constante $L$ en $[t_{b-1},t_b]$, de ancho
$w=t_b-t_{b-1}$, y $c$ el punto medio del intervalo. Entonces
$|f(u)-f(c)|\le Lw/2$ para todo $u$ del intervalo.
\end{proposicion}

\begin{proof}
$|f(u)-f(c)|\le L|u-c|\le Lw/2$, porque ningún punto del intervalo dista
más de $w/2$ de su punto medio.
\end{proof}

Si la media de la respuesta es una función $f$ lipschitziana de $x$, la
mejor constante en cada intervalo aproxima $f$ con error a lo sumo $Lw/2$,
y la media de $f$ sobre los puntos del intervalo, que es lo que estima un
modelo con una constante por intervalo (discretización seguida de
\emph{one-hot}, \S\ref{sec:onehot}), con error a lo sumo $Lw$ (porque esa
media está entre el mínimo y el máximo de $f$ en el intervalo, que difieren
en a lo sumo $Lw$). El sesgo disminuye con $w$; pero con intervalos más
estrechos cada constante se estima con menos observaciones y con más
varianza. A cambio
de la pérdida de información, la discretización acota la influencia de los
valores extremos (todos caen en el último intervalo, con independencia de
su magnitud), permite a un modelo lineal representar relaciones no
lineales mediante funciones escalonadas y es el paso previo del WoE para
variables continuas (\S\ref{sec:woe}). En árboles no aporta: los cortes
posibles pasan a ser a lo sumo $B-1$, un subconjunto de los que el árbol
examina sobre la variable original (Lema~\ref{lem:cortes}).

\begin{aws}
\textbf{Glue DataBrew}, acción \texttt{BUCKETIZATION}, con cortes
explícitos en \texttt{splits} (admite \texttt{-Infinity} e
\texttt{Infinity}) o con \texttt{percentage} para intervalos de igual
frecuencia; \texttt{BINARIZATION} para un único umbral:
\begin{lstlisting}
{"Action": {"Operation": "BUCKETIZATION",
  "Parameters": {"sourceColumn": "peso", "targetColumn": "peso_bin",
                 "bucketNames": "[\"Bin1\",\"Bin2\",\"Bin3\"]",
                 "splits": "[\"-Infinity\",\"2\",\"20\",\"Infinity\"]"}}}
\end{lstlisting}
\textbf{Data Wrangler} no tiene paso de discretización en su lista vigente;
se implementa con \emph{Custom formula} (\texttt{CASE WHEN}) o
\emph{Custom transform} en PySpark (\texttt{Bucketizer},
\texttt{QuantileDiscretizer}), seguido de \emph{Encode categorical}
$\to$ \emph{One-hot encode}. \textbf{SageMaker Processing}:
\texttt{KBinsDiscretizer(strategy="uniform"|"quantile"|"kmeans",
encode="onehot")}.
\end{aws}

% =====================================================================
\section{Características categóricas: codificaciones no supervisadas}
% =====================================================================

\subsection{Planteamiento}

Sea $X$ una característica categórica con niveles
$\mathcal C=\{c_1,\dots,c_K\}$, los observados en entrenamiento. Una
\emph{codificación} es una función $\phi:\mathcal C\to\R^d$; el modelo
recibe $\phi(X)$ en lugar de $X$. Cinco rasgos distinguen las
codificaciones: si $\phi$ es inyectiva (si puede confundir niveles), la
dimensión $d$, la geometría que induce (las distancias
$\|\phi(c_i)-\phi(c_j)\|$, que usan los métodos basados en distancias y
las redes), si usa la respuesta, y qué valor asigna en inferencia a un
nivel que no está en $\mathcal C$. Para modelos lineales, el rasgo decisivo
es el siguiente.

\begin{definicion}
Un \emph{predictor lineal} sobre $\phi$ es una función
$\eta(c)=\beta_0+\beta^\top\phi(c)$, $(\beta_0,\beta)\in\R^{1+d}$. Su
\emph{vector de efectos} es $(\eta(c_1),\dots,\eta(c_K))\in\R^K$. Sea
$\Phi\in\R^{K\times d}$ la matriz cuya fila $i$ es $\phi(c_i)^\top$.
\end{definicion}

\begin{proposicion}[Efectos representables]\label{prop:repr}
El conjunto de vectores de efectos de los predictores lineales sobre
$\phi$ es el espacio columna de la matriz $[\1\;\Phi]\in\R^{K\times(1+d)}$;
su dimensión es $\rango[\1\;\Phi]\le\min(K,d+1)$. En particular, todo
vector de $\R^K$ es un vector de efectos si y solo si
$\rango[\1\;\Phi]=K$.
\end{proposicion}

\begin{proof}
La coordenada $i$ del vector de efectos es
$\beta_0+\phi(c_i)^\top\beta$, la coordenada $i$ de
$[\1\;\Phi](\beta_0,\beta^\top)^\top$. El conjunto de vectores de efectos
es, por tanto, la imagen de la aplicación lineal
$(\beta_0,\beta)\mapsto[\1\;\Phi](\beta_0,\beta^\top)^\top$, que es el
espacio columna; su dimensión es el rango, acotado por el número de filas
y de columnas. Un subespacio de $\R^K$ es todo $\R^K$ si y solo si su
dimensión es $K$.
\end{proof}

Si la respuesta media depende de los niveles de forma arbitraria, una
codificación con rango menor que $K$ impide a un modelo lineal
representarla. Los árboles y las redes no están sujetos a esta
proposición, pero sí al orden y a la geometría que $\phi$ induce.

\subsection{Codificación ordinal}

\begin{definicion}
Dada una biyección $\sigma:\mathcal C\to\{1,\dots,K\}$ (un orden de los
niveles), la \emph{codificación ordinal} es $\phi(c)=\sigma(c)\in\R$. El
capítulo la llama \emph{label encoding}.
\end{definicion}

Si la variable es ordinal (tallas S $<$ M $<$ L $<$ XL), $\sigma$ es el orden
natural y la codificación lo conserva; si es nominal, $\sigma$ es
arbitraria (\texttt{LabelEncoder} de scikit-learn usa el orden
alfabético). Con $K\ge2$, $\rango[\1\;\Phi]=2$ (las columnas $\1$ y
$(1,2,\dots,K)^\top$ no son proporcionales), así que, por la
Proposición~\ref{prop:repr}, un predictor lineal solo representa efectos de
la forma $\beta_0+\beta\sigma(c)$: equiespaciados en el orden elegido.

\begin{contraejemplo}
Si los efectos de las tallas S, M, L, XL son $(0,0,0,10)$, ningún
$(\beta_0,\beta)$ los representa: $\beta_0+\beta=\beta_0+2\beta$ obliga a
$\beta=0$, y entonces el cuarto efecto sería $\beta_0=0\ne10$. Aun en una
variable ordinal, la codificación ordinal restringe la forma del efecto.
\end{contraejemplo}

Para árboles, el Lema~\ref{lem:cortes} dice que un corte sobre
$\sigma(X)$ separa los niveles en $\{c:\sigma(c)\le t\}$ y su
complemento: hay $K-1$ particiones posibles, mientras que el número de
particiones de $\mathcal C$ en dos conjuntos no vacíos es
$(2^K-2)/2=2^{K-1}-1$ (los $2^K$ subconjuntos menos $\emptyset$ y
$\mathcal C$, contando cada partición dos veces). Aislar un nivel de
posición $t$ requiere a lo sumo dos cortes: $\sigma\le t-1$ y, en el nodo
$\{\sigma\ge t\}$, $\sigma\le t$. Así, un árbol puede representar
cualquier función de los niveles con suficientes cortes, pero el orden
$\sigma$ determina cuántos necesita; el Teorema~\ref{teo:fisher} identifica
un orden con el que basta un corte para la mejor partición en dos.

\begin{aws}
\textbf{Data Wrangler}, \emph{Encode categorical} $\to$ \emph{Ordinal
encode} (enteros desde $0$), con manejo de niveles no vistos \emph{Skip},
\emph{Keep}, \emph{Error} o \emph{Replace with NaN}; paso con parámetros
ajustados (\emph{refit}). \textbf{Glue DataBrew},
\texttt{CATEGORICAL\_MAPPING}, que permite fijar $\sigma$ explícitamente:
\begin{lstlisting}
{"Action": {"Operation": "CATEGORICAL_MAPPING",
  "Parameters": {"sourceColumn": "talla", "targetColumn": "talla_ord",
                 "categoryMap": "{\"S\":\"1\",\"M\":\"2\",\"L\":\"3\",\"XL\":\"4\"}",
                 "mapType": "NUMERIC", "keepOthers": "false",
                 "deleteOtherRows": "false", "other": "0"}}}
\end{lstlisting}
\textbf{SageMaker Processing}: \texttt{OrdinalEncoder(categories=[[...]])}
con el orden explícito.
\end{aws}

\subsection{Codificación \emph{one-hot}}\label{sec:onehot}

\begin{definicion}
La codificación \emph{one-hot} es $\phi(c_i)=e_i\in\R^K$, el $i$-ésimo
vector de la base canónica. La \emph{codificación de referencia} elimina
la coordenada de un nivel fijado (el de referencia), con $d=K-1$. Para una
muestra, la matriz \emph{one-hot} $D\in\{0,1\}^{n\times K}$ tiene
$D_{ki}=1$ si y solo si la observación $k$ tiene el nivel $c_i$.
\end{definicion}

\begin{proposicion}[Trampa de las variables ficticias]\label{prop:trampa}
La matriz $Z=[\1\;D]\in\R^{n\times(K+1)}$ no tiene rango completo por
columnas, y $Z^\top Z$ no es invertible.
\end{proposicion}

\begin{proof}
Cada fila de $D$ tiene exactamente un $1$, luego $D\1_K=\1_n$ y
$Z\,(1,-\1_K^\top)^\top=\1_n-D\1_K=0$ con $(1,-\1_K^\top)\ne0$. Como
$Z^\top Zv=0$ implica $\|Zv\|^2=v^\top Z^\top Zv=0$, el núcleo de
$Z^\top Z$ coincide con el de $Z$, que contiene un vector no nulo.
\end{proof}

Las salidas son omitir un nivel (codificación de referencia), omitir el
intercepto o penalizar: con penalización \emph{ridge}, la matriz $Z^\top
Z+\lambda I$ con $\lambda>0$ es definida positiva ($v^\top(Z^\top
Z+\lambda I)v=\|Zv\|^2+\lambda\|v\|^2>0$ si $v\ne0$) y la solución es única.

\begin{proposicion}[Interpretación]\label{prop:onehot}
Supóngase que todos los niveles aparecen en la muestra. En la regresión por
mínimos cuadrados de $y$ sobre el intercepto y la codificación de
referencia (referencia $c_K$), el valor ajustado de toda observación del
nivel $c_i$ es la media $\bar y_i$ de las respuestas de ese nivel; los
coeficientes son $\hat\beta_0=\bar y_K$ y $\hat\beta_i=\bar y_i-\bar y_K$.
\end{proposicion}

\begin{proof}
El espacio columna de $[\1\;D_{\cdot1}\cdots D_{\cdot,K-1}]$ contiene a
$D_{\cdot K}=\1-\sum_{i<K}D_{\cdot i}$, así que coincide con $\col(D)$, el
conjunto de vectores que son constantes sobre las observaciones de cada
nivel. Por el Lema~\ref{lem:proy}, los valores ajustados minimizan
$\sum_i\sum_{k:x_k=c_i}(y_k-a_i)^2$ sobre las constantes $(a_i)$, y cada
sumando se minimiza por separado con $a_i=\bar y_i$. Las columnas del
diseño de referencia son linealmente independientes (si
$\beta_0\1+\sum_{i<K}\beta_iD_{\cdot i}=0$, evaluando en una observación de
$c_K$ se obtiene $\beta_0=0$ y luego, en una de cada $c_i$, $\beta_i=0$),
así que los coeficientes son únicos, y de $\hat\beta_0=\bar y_K$ y
$\hat\beta_0+\hat\beta_i=\bar y_i$ se obtienen los valores del enunciado.
\end{proof}

El rango de $[\1\;\Phi]$ es $K$ (las filas $(1,e_i^\top)$ son linealmente
independientes), de modo que la codificación es \emph{saturada}: no impone
restricción alguna (Proposición~\ref{prop:repr}). Su geometría es la
opuesta a la ordinal: $\|e_i-e_j\|=\sqrt2$ y $\langle e_i,e_j\rangle=0$
para $i\ne j$, así que no introduce orden ni similitudes artificiales,
pero tampoco puede expresar similitudes reales. Su costo es la dimensión y
la varianza: bajo el modelo $y_k=\theta_{x_k}+\varepsilon_k$ con errores
independientes de varianza $\sigma^2$, el valor ajustado del nivel $c_i$
es $\bar y_i$, con varianza $\sigma^2/n_i$, que es grande para niveles
raros. Un nivel no visto en entrenamiento recibe el vector nulo
(\texttt{handle\_unknown="ignore"}), es decir, la predicción del nivel de
referencia o del intercepto. La Proposición~\ref{prop:ridge-te} muestra que
la penalización \emph{ridge} sobre \emph{one-hot} reduce esa varianza con
el mismo mecanismo que el \emph{target encoding}.

\begin{aws}
\textbf{Data Wrangler}, \emph{Encode categorical} $\to$ \emph{One-hot
encode}, con la opción \emph{Drop last category} (codificación de
referencia, Proposición~\ref{prop:trampa}), manejo de niveles no vistos y
salida como vector disperso o columnas; el paso \emph{Replace rare}
(grupo \emph{Handle outliers}) agrupa antes los niveles poco frecuentes en
uno solo. \textbf{Glue DataBrew}, \texttt{ONE\_HOT\_ENCODING} con
\texttt{sourceColumn}, que admite como máximo 10 valores distintos.
\textbf{SageMaker Processing}: \texttt{OneHotEncoder(handle\_unknown="ignore",
drop="first", min\_frequency=...)}. Los algoritmos integrados
\emph{Linear Learner} y \emph{Factorization Machines} aceptan entrada
dispersa en formato \texttt{recordIO-protobuf}.
\end{aws}

\subsection{Codificación binaria}

\begin{definicion}
Dada una numeración $\sigma:\mathcal C\to\{1,\dots,K\}$, sea $d$ el menor
entero con $2^d-1\ge K$. La \emph{codificación binaria} asigna a $c$ el
vector $\phi(c)\in\{0,1\}^d$ de los dígitos de $\sigma(c)$ en base $2$.
\end{definicion}

Con $d$ dígitos binarios se escriben los enteros de $0$ a $2^d-1$; como
la numeración empieza en $1$ (convención de \texttt{category\_encoders}),
se necesita $2^d-1\ge K$, es decir, $d=\lceil\log_2(K+1)\rceil$. Con los
8 colores del capítulo, $d=4$; con $K=1000$, $d=10$. La codificación es
inyectiva, porque la expansión binaria de un entero es única: no pierde
información.

\begin{proposicion}\label{prop:binaria}
Si $K\ge5$, $\rango[\1\;\Phi]\le d+1<K$: un predictor lineal sobre la
codificación binaria no representa todos los vectores de efectos.
\end{proposicion}

\begin{proof}
$[\1\;\Phi]$ tiene $d+1$ columnas. Como
$d=\lceil\log_2(K+1)\rceil<\log_2(K+1)+1$, basta ver que
$\log_2(K+1)+2\le K$ para $K\ge5$: vale en $K=5$ ($\log_26+2\approx4.58$) y
la función $K-\log_2(K+1)$ es creciente para $K\ge1$ (su derivada es
$1-1/((K+1)\ln 2)>0$). La conclusión se sigue de la
Proposición~\ref{prop:repr}.
\end{proof}

La geometría inducida depende de la numeración. Con la numeración por
orden de aparición del capítulo (White $=1$, \dots, Brown $=8$), Red
($0011$) y Pink ($0111$) difieren en un bit y White ($0001$) y Brown
($1000$) en dos; con otra numeración, las distancias serían otras. En
árboles, un corte sobre un bit separa los niveles según ese dígito, otra
partición fijada por la numeración. La codificación binaria reduce la
dimensión de \emph{one-hot} de $K$ a $\approx\log_2K$ a cambio de una
estructura que no proviene de los datos.

\begin{aws}
No hay paso integrado en Data Wrangler ni en DataBrew. En \textbf{Data
Wrangler}, \emph{Ordinal encode} seguido de una \emph{Custom formula} de
Spark SQL por bit, \texttt{shiftright(color\_idx, j) \& 1},
$j=0,\dots,d-1$. En \textbf{SageMaker Processing},
\texttt{category\_encoders.BinaryEncoder}, instalado en el contenedor.
\end{aws}

\subsection{\emph{Hashing} de características}\label{sec:hashing}

\paragraph{Qué se representa.} El \emph{hashing} se aplica a vectores
indexados por \emph{símbolos}. Sea $\mathcal T$ el conjunto de símbolos
posibles, que puede ser infinito o desconocido de antemano: todas las
cadenas de caracteres que podrían aparecer como nivel de una categórica, o
todos los tokens que podrían aparecer en un texto (\S\ref{sec:texto}). Sea
$\R^{(\mathcal T)}$ el espacio de vectores $x=(x_t)_{t\in\mathcal T}$ con
solo un número finito de coordenadas no nulas, con producto interno
$\langle x,x'\rangle=\sum_tx_tx'_t$ (suma finita). Los dos casos del
capítulo son:
\begin{enumerate}[label=(\alph*)]
\item una característica categórica: la observación con nivel $c$ se
representa por $x=e_c$, con $x_c=1$ y $x_t=0$ para $t\ne c$ (la
codificación \emph{one-hot} indexada por todos los niveles posibles);
\item un documento: tras tokenizarlo, $x_t$ es el número de veces que el
token $t$ aparece en él (la bolsa de palabras de \S\ref{sec:tfidf}).
\end{enumerate}

\paragraph{El problema.} Usar $x$ directamente requiere un diccionario que
asigne a cada símbolo observado en entrenamiento una columna. Su memoria
crece con el número de símbolos distintos (millones en filtros de
\emph{spam} o en identificadores de usuario), y un símbolo que aparece por
primera vez en producción no tiene columna. Weinberger et al.\ (2009)
propusieron fijar de antemano el número de columnas.

\begin{definicion}[\emph{Hashing} de características]
Sean $m\in\N$, $h:\mathcal T\to\{1,\dots,m\}$ y
$\xi:\mathcal T\to\{-1,+1\}$ dos funciones. La representación con
\emph{hashing} de $x\in\R^{(\mathcal T)}$ es $\phi(x)\in\R^m$ con
\[
\phi_j(x)=\sum_{t:\,h(t)=j}\xi(t)\,x_t,\qquad j=1,\dots,m,
\]
suma finita porque $x$ tiene finitas coordenadas no nulas. En el caso (a),
$\phi(e_c)=\xi(c)\,e_{h(c)}$: un vector con un solo $\pm1$ en la
posición $h(c)$. En el caso (b), la coordenada $j$ suma, con signo, los
conteos de los tokens que $h$ envía a $j$.
\end{definicion}

Dos símbolos $t\ne u$ \emph{colisionan} si $h(t)=h(u)$. En la práctica,
$h$ y $\xi$ se calculan con una función hash de cadenas: scikit-learn
aplica MurmurHash3 con semilla fija a la cadena del símbolo, toma
$h(t)$ como el valor absoluto del resultado módulo $m$ y $\xi(t)$ como su
signo. Son funciones deterministas: el mismo símbolo produce siempre la
misma columna y el mismo signo, sin necesidad de guardar nada. Para
analizarlas se usa el modelo siguiente, que describe una elección
aleatoria de la función (por ejemplo, de su semilla), con $x$ fijo.

\paragraph{Modelo de análisis.} (H1) Para símbolos distintos, los valores
$h(t)$ son independientes y uniformes en $\{1,\dots,m\}$. (H2) Los
valores $\xi(t)$ son independientes y uniformes en $\{-1,+1\}$. (H3) $h$ y
$\xi$ son independientes. (En scikit-learn $h$ y $\xi$ provienen del mismo
valor hash, así que (H3) es una idealización.)

\begin{proposicion}[Insesgadez del producto interno]\label{prop:hash}
Bajo (H2) y (H3), para todo par $x,x'\in\R^{(\mathcal T)}$,
$\E\bigl[\langle\phi(x),\phi(x')\rangle\bigr]=\langle x,x'\rangle$.
\end{proposicion}

\begin{proof}
Sea $S$ el conjunto finito de símbolos con $x_t\ne0$ o $x'_t\ne0$.
Entonces
$\langle\phi(x),\phi(x')\rangle=\sum_{j=1}^m\sum_{t,u\in S}\ind{h(t)=j}\ind{h(u)=j}\xi(t)\xi(u)x_tx'_u
=\sum_{t,u\in S}\ind{h(t)=h(u)}\xi(t)\xi(u)x_tx'_u$. Por (H3), la
esperanza de cada sumando es $\Prob(h(t)=h(u))\,\E[\xi(t)\xi(u)]\,x_tx'_u$.
Si $t\ne u$, por (H2) $\E[\xi(t)\xi(u)]=\E[\xi(t)]\E[\xi(u)]=0$; si $t=u$,
$\xi(t)^2=1$ y $\Prob(h(t)=h(t))=1$. Queda $\sum_{t\in S}x_tx'_t$.
\end{proof}

En el caso (a) con $c\ne c'$, $\langle e_c,e_{c'}\rangle=0$: en promedio
sobre la elección de la función, los niveles distintos siguen siendo
ortogonales, como en \emph{one-hot}. El signo es la causa: si
$\xi\equiv1$, el sumando de un par $t\ne u$ tiene esperanza
$\Prob(h(t)=h(u))x_tx'_u=x_tx'_u/m$ bajo (H1), y el producto interno queda
sesgado en $\sum_{t\ne u}x_tx'_u/m$, positivo para conteos: las colisiones
siempre sumarían. Con signo, se cancelan en promedio; por eso
\texttt{FeatureHasher} usa \texttt{alternate\_sign=True} y en la salida
del capítulo aparecen valores $\pm1$.

\begin{proposicion}[Número de colisiones]\label{prop:colis}
Bajo (H1), entre $K$ símbolos distintos el número esperado de pares que
colisionan es $\binom K2/m$, y la probabilidad de que no haya ninguna
colisión es $\prod_{i=1}^{K-1}(1-i/m)\le\exp\bigl(-K(K-1)/(2m)\bigr)$.
Un símbolo dado colisiona con algún otro con probabilidad
$1-(1-1/m)^{K-1}$.
\end{proposicion}

\begin{proof}
Para un par, $\Prob(h(t)=h(u))=\sum_{j=1}^m\Prob(h(t)=j)\Prob(h(u)=j)=m\cdot
m^{-2}=1/m$; la esperanza del número de pares en colisión es la suma sobre
los $\binom K2$ pares de sus indicadoras. Para la segunda afirmación, se
asignan los símbolos en orden: dado que los $i$ primeros ocupan $i$
valores distintos, el siguiente evita esos valores con probabilidad
$1-i/m$ por la independencia de (H1); el producto de estas probabilidades
condicionales es la probabilidad de no colisión, y la cota usa $1-v\le
e^{-v}$ y $\sum_{i=1}^{K-1}i=K(K-1)/2$. La tercera: el símbolo evita el
valor de cada uno de los otros $K-1$ independientemente con probabilidad
$1-1/m$.
\end{proof}

\begin{ejemplo}
En el ejemplo del capítulo hay $K=8$ colores y $m=3$. El número esperado
de pares en colisión es $28/3\approx9.3$. En la salida real, la columna
\texttt{feature\_1} recibe White, Black y Blue (3 pares), la columna
\texttt{feature\_2} recibe Green, Yellow, Pink y Brown (6 pares) y Red
queda sola en \texttt{feature\_0}: 9 pares. El signo distingue a
Black ($+1$) de White y Blue ($-1$), y a Yellow ($+1$) de Green, Pink y
Brown ($-1$); pero White y Blue quedan con el mismo vector, y también
Green, Pink y Brown. Por el principio del palomar, con $m=3$ columnas y
dos signos hay a lo sumo $6$ vectores distintos para $8$ colores, así que
alguna coincidencia era inevitable.
\end{ejemplo}

Si $K$ es del orden de $m$, la tercera fórmula da
$1-(1-1/m)^{K-1}\approx1-e^{-1}\approx0.63$: la mayoría de los símbolos
comparte columna. Por eso se usan valores de $m$ del orden de $2^{18}$ a
$2^{20}$ y se acepta que colisionen sobre todo los símbolos raros.

\paragraph{Contraste.} Frente a \emph{one-hot}, el \emph{hashing} fija la
dimensión, no guarda diccionario y procesa símbolos nuevos sin
reentrenar; a cambio no es inyectivo ni invertible (una columna no indica
qué símbolos la forman) y $m$ es un hiperparámetro. Frente a la
codificación binaria, ambas reducen dimensión, pero la binaria es inyectiva
y el \emph{hashing} no; el \emph{hashing} conserva en promedio los
productos internos de \emph{one-hot} (Proposición~\ref{prop:hash}) y la
binaria introduce similitudes fijadas por la numeración. Como no tiene
parámetros que ajustar, el \emph{hashing} no puede producir fuga, siempre
que $h$, $\xi$ y $m$ sean los mismos en entrenamiento y en producción.

\subsubsection*{Variante: codificación por similitud con \emph{min-hash}}

Cuando los niveles son cadenas con variantes de escritura («juan perez»,
«juan peres»), \emph{one-hot} y \emph{hashing} las tratan como símbolos
ajenos. La codificación por similitud (Cerda y Varoquaux, 2022) representa
cada nivel por sus fragmentos de tres caracteres.

\begin{definicion}
Sea $\mathcal A$ un alfabeto finito y $U=\mathcal A^3$ el conjunto de
cadenas de longitud 3. Para una cadena $c$ de longitud $\ge3$, sea
$S(c)\subset U$ el conjunto de sus subcadenas de tres caracteres
consecutivos (sus \emph{3-gramas}). La \emph{similitud de Jaccard} de dos
conjuntos no vacíos es $J(A,B)=|A\cap B|/|A\cup B|$. Para una biyección
$\pi:U\to\{1,\dots,|U|\}$ (una permutación de $U$), el \emph{min-hash} de
$S\subset U$ no vacío es $h_\pi(S)=\min_{g\in S}\pi(g)$.
\end{definicion}

Por ejemplo, «juan perez» y «juan peres» tienen 7 3-gramas en común de 9
distintos en total, $J=7/9\approx0.78$; «juan perez» y «maria lopez» no
comparten ninguno, $J=0$.

\begin{proposicion}[Broder, 1997]\label{prop:minhash}
Si $\pi$ es una permutación aleatoria uniforme de $U$, entonces
$\Prob\bigl(h_\pi(A)=h_\pi(B)\bigr)=J(A,B)$ para $A,B\subset U$ no vacíos.
\end{proposicion}

\begin{proof}
Sea $g^*$ el elemento de $A\cup B$ con menor valor de $\pi$. Como $\pi$ es
uniforme, todos los elementos de $A\cup B$ tienen la misma probabilidad de
ser $g^*$, así que $\Prob(g^*\in A\cap B)=|A\cap B|/|A\cup B|$. Si $g^*\in
A\cap B$, es el mínimo de ambos conjuntos y $h_\pi(A)=h_\pi(B)=\pi(g^*)$.
Si $g^*\in A\setminus B$, entonces $h_\pi(A)=\pi(g^*)$, mientras que
$h_\pi(B)>\pi(g^*)$ porque todo elemento de $B$ está en $A\cup B$ y es
distinto de $g^*$; análogamente si $g^*\in B\setminus A$. Luego
$h_\pi(A)=h_\pi(B)$ si y solo si $g^*\in A\cap B$.
\end{proof}

La codificación con $d$ permutaciones independientes
$\pi_1,\dots,\pi_d$ es
\[
\phi(c)=\bigl(h_{\pi_1}(S(c)),\dots,h_{\pi_d}(S(c))\bigr).
\]
La fracción
de coordenadas en que coinciden $\phi(c)$ y $\phi(c')$ es un promedio de
$d$ indicadoras independientes de media $J(S(c),S(c'))$, luego estima esa
similitud sin sesgo y con varianza $J(1-J)/d$. Niveles con muchos 3-gramas
comunes reciben vectores que coinciden en muchas coordenadas: la geometría
refleja similitud de escritura, lo contrario de las colisiones arbitrarias
del \emph{hashing}. Las implementaciones sustituyen las permutaciones
aleatorias por funciones hash, que las aproximan.

\begin{aws}
\textbf{Data Wrangler}: \emph{Encode categorical} $\to$ \emph{Similarity
encode} implementa la codificación por similitud (3-gramas y
\emph{min-hash}, dimensión $d=30$ por defecto, salida como vector o
columnas). El \emph{hashing} está en \emph{Featurize text} $\to$
\emph{Vectorize} con el vectorizador \emph{Hashing} (número de columnas
$m$ configurable), pensado para el caso (b); para una categórica (caso
(a)) se usa \emph{Custom transform} con \texttt{FeatureHasher} de PySpark o
de scikit-learn. \textbf{SageMaker Processing}:
\texttt{FeatureHasher(n\_features=2**20, input\_type="string")}.
\end{aws}

% =====================================================================
\section{Características categóricas: codificaciones supervisadas}
% =====================================================================

Con cardinalidad alta, las codificaciones anteriores fallan de algún modo:
\emph{one-hot} estima un efecto por nivel con varianza $\sigma^2/n_i$; la
ordinal impone un orden arbitrario; la binaria y el \emph{hashing}
introducen estructura ajena a los datos. Para pérdida cuadrática, la
función del nivel que mejor predice $Y$ es $\E[Y\mid X=c_i]$ (por el
argumento de la demostración de la Proposición~\ref{prop:ar}, con la
historia reemplazada por $X$), de modo que codificar cada nivel por una
estimación de esa cantidad resume en un número lo que el nivel dice sobre
$Y$. Es la idea del \emph{target encoding}; el \emph{weight of evidence}
es su versión en la escala de log-momios.

\subsection{\emph{Target encoding}}\label{sec:target}

\paragraph{Notación.} En la muestra de entrenamiento, para el nivel $c_i$:
$n_i$ es su número de observaciones, $L_i=\{k:x_k=c_i\}$ su conjunto de
índices y $\bar y_i=n_i^{-1}\sum_{k\in L_i}y_k$ su media de la respuesta;
$\bar y$ es la media global e $i(k)$ el índice del nivel de la
observación $k$. Si $y\in\{0,1\}$, $\bar y_i$ es la proporción de eventos
del nivel.

\begin{definicion}[\emph{Target encoding}]
Para $m\ge0$, la codificación del nivel $c_i$ es
\[
S_i=\lambda_i\bar y_i+(1-\lambda_i)\bar y,\qquad\lambda_i=\frac{n_i}{n_i+m}.
\]
Con $m=0$ es la media del nivel sin suavizar. A un nivel no visto en
entrenamiento ($n_i=0$, $\lambda_i=0$) le corresponde $\bar y$. Si la
respuesta tiene $C$ clases, se calcula una columna por clase $j$ con $y$
reemplazada por $\ind{y=j}$.
\end{definicion}

El suavizado responde al problema señalado por Micci-Barreca (2001): un
nivel observado una vez con $y=1$ tendría $\bar y_i=1$, una estimación
basada en un solo dato. El peso $\lambda_i$ interpola entre la media del
nivel, fiable si $n_i$ es grande, y la global, preferible si $n_i$ es
pequeño; $m$ es el número de observaciones con el que ambas pesan lo
mismo. Los dos resultados siguientes muestran que esta forma es la media
posterior de modelos jerárquicos estándar.

\begin{proposicion}[Respuesta binaria]\label{prop:betabin}
Supóngase que, dado $\theta_i$, las respuestas del nivel $c_i$ son
Bernoulli$(\theta_i)$ independientes, y que a priori
$\theta_i\sim\text{Beta}(a,b)$ con $a,b>0$. Entonces
\[
\E[\theta_i\mid(y_k)_{k\in L_i}]=\frac{a+n_i\bar y_i}{a+b+n_i}
=\lambda_i\bar y_i+(1-\lambda_i)\frac a{a+b},\qquad
\lambda_i=\frac{n_i}{n_i+a+b}.
\]
\end{proposicion}

\begin{proof}
Con $s=n_i\bar y_i$ eventos, la verosimilitud es
$\theta_i^s(1-\theta_i)^{n_i-s}$ y la densidad a priori es proporcional a
$\theta_i^{a-1}(1-\theta_i)^{b-1}$; la posterior es proporcional a su
producto, $\theta_i^{a+s-1}(1-\theta_i)^{b+n_i-s-1}$, que es la densidad
Beta$(a+s,b+n_i-s)$, de media $(a+s)/(a+b+n_i)$. Escribiendo
$a+s=(a+b)\frac a{a+b}+n_i\bar y_i$ y dividiendo por $a+b+n_i$ se obtiene
la segunda forma.
\end{proof}

La codificación es la media posterior con media a priori $a/(a+b)$
(estimada por $\bar y$) y $m=a+b$, que actúa como un número de
observaciones previas.

\begin{proposicion}[Respuesta continua]\label{prop:normnorm}
Supóngase que, dado $\theta_i$, las respuestas del nivel son
$N(\theta_i,\sigma^2)$ independientes, y que $\theta_i\sim N(\mu,\tau^2)$,
con $\mu$, $\sigma^2>0$ y $\tau^2>0$ conocidos. Entonces la distribución
de $\theta_i$ dada la muestra del nivel es $N(S_i,v_i)$ con
\[
S_i=\lambda_i\bar y_i+(1-\lambda_i)\mu,\qquad
\lambda_i=\frac{n_i}{n_i+\sigma^2/\tau^2},\qquad
v_i=\Bigl(\frac1{\tau^2}+\frac{n_i}{\sigma^2}\Bigr)^{-1}.
\]
Además $\E[(S_i-\theta_i)^2]=v_i<\sigma^2/n_i=\E[(\bar y_i-\theta_i)^2]$.
\end{proposicion}

\begin{proof}
La verosimilitud de $\theta_i$ depende de los datos solo a través de
$\bar y_i$, que dado $\theta_i$ es $N(\theta_i,\sigma^2/n_i)$. El
logaritmo de la densidad posterior es, salvo constantes,
$-\frac{n_i}{2\sigma^2}(\bar y_i-\theta_i)^2-\frac1{2\tau^2}(\theta_i-\mu)^2$,
un polinomio de grado 2 en $\theta_i$ con coeficiente cuadrático
$-\frac12\bigl(\frac{n_i}{\sigma^2}+\frac1{\tau^2}\bigr)=-\frac1{2v_i}$ y
coeficiente lineal $\frac{n_i\bar y_i}{\sigma^2}+\frac\mu{\tau^2}$.
Completando el cuadrado, la posterior es normal con varianza $v_i$ y media
$v_i\bigl(\frac{n_i\bar y_i}{\sigma^2}+\frac\mu{\tau^2}\bigr)
=\frac{n_i\tau^2\bar y_i+\sigma^2\mu}{n_i\tau^2+\sigma^2}=\lambda_i\bar
y_i+(1-\lambda_i)\mu$. Para los errores: condicionando en los datos,
$\E[(S_i-\theta_i)^2\mid\text{datos}]=v_i$ (la varianza posterior), que
no depende de los datos; condicionando en $\theta_i$,
$\E[(\bar y_i-\theta_i)^2\mid\theta_i]=\sigma^2/n_i$. La desigualdad
$v_i<\sigma^2/n_i$ equivale a $\frac1{\tau^2}+\frac{n_i}{\sigma^2}>\frac{n_i}{\sigma^2}$.
\end{proof}

La proposición identifica $m=\sigma^2/\tau^2$: la dispersión dentro de los
niveles frente a la dispersión entre sus medias. En la práctica $\mu$,
$\sigma^2$ y $\tau^2$ se estiman con los datos (Bayes empírico).
\texttt{TargetEncoder(smooth="auto")} de scikit-learn usa, según su
documentación, $m_i=\sigma_i^2/\tau^2$, con $\sigma_i^2$ la varianza de
$y$ dentro del nivel y $\tau^2$ la varianza de $y$ en toda la muestra; en
el modelo, $\tau^2$ es la varianza entre medias de nivel, y por la ley de
la varianza total la de $y$ es $\sigma^2+\tau^2>\tau^2$, de modo que esa
elección suaviza menos que la regla del modelo.

\begin{proposicion}[\emph{Ridge} sobre \emph{one-hot} es un suavizado]\label{prop:ridge-te}
Supóngase que todos los niveles aparecen y sea $\lambda>0$. La función
$F(\beta_0,\beta)=\sum_k(y_k-\beta_0-\beta_{i(k)})^2+\lambda\sum_{i=1}^K\beta_i^2$
tiene un único minimizador, y en él el valor ajustado para el nivel $c_i$
es $\hat\beta_0+\lambda_i(\bar y_i-\hat\beta_0)$ con
$\lambda_i=n_i/(n_i+\lambda)$.
\end{proposicion}

\begin{proof}
$F$ es cuadrática con hessiana $2(Z^\top Z+\operatorname{diag}(0,\lambda
I_K))$, con $Z=[\1\;D]$. Si $v=(v_0,v_1,\dots,v_K)$ anula la forma
cuadrática, $\|Zv\|^2+\lambda\sum_{i\ge1}v_i^2=0$, así que $v_i=0$ para
$i\ge1$ y $Zv=v_0\1=0$, es decir, $v=0$: la hessiana es definida positiva,
$F$ es estrictamente convexa y su único minimizador es el punto donde se
anula el gradiente. La derivada respecto de $\beta_i$ es
$-2\sum_{k\in L_i}(y_k-\beta_0-\beta_i)+2\lambda\beta_i$; igualada a cero,
$(n_i+\lambda)\beta_i=n_i(\bar y_i-\beta_0)$, de donde
$\hat\beta_i=\lambda_i(\bar y_i-\hat\beta_0)$, y el valor ajustado del
nivel es $\hat\beta_0+\hat\beta_i$.
\end{proof}

La penalización \emph{ridge} sobre $K$ columnas \emph{one-hot} y el
\emph{target encoding} en una columna producen la misma contracción, con
$m=\lambda$ y $\hat\beta_0$ en el papel de la media a priori. La
diferencia práctica es que el \emph{target encoding} reduce la dimensión a
$1$ antes de modelar, y la columna resultante puede usarse con cualquier
algoritmo.

\subsubsection{Fuga de información dentro del entrenamiento}

La Proposición~\ref{prop:test} exige no usar la muestra de prueba para
ajustar; el resultado siguiente muestra que, además, usar la propia
observación en su codificación distorsiona el entrenamiento.

\begin{proposicion}[Codificación dentro de la muestra]\label{prop:te-fuga}
Supóngase que $y_k=\theta_{i(k)}+\varepsilon_k$, con constantes
$\theta_1,\dots,\theta_K$ y errores $\varepsilon_1,\dots,\varepsilon_n$
independientes, de media $0$ y varianza $\sigma^2$, y considérense fijos
los niveles de las observaciones. Si $S_{i(k)}$ se calcula con toda la
muestra de entrenamiento, entonces, con $i=i(k)$,
\[
\Cov\bigl(S_{i},\varepsilon_k\bigr)=\lambda_i\frac{\sigma^2}{n_i}+(1-\lambda_i)\frac{\sigma^2}{n}>0,
\]
mientras que para una observación nueva, cuyo error es independiente de
los de entrenamiento, la covarianza correspondiente es $0$.
\end{proposicion}

\begin{proof}
$\bar y_i=\theta_i+n_i^{-1}\sum_{l\in L_i}\varepsilon_l$ y
$\bar y=n^{-1}\sum_l\theta_{i(l)}+n^{-1}\sum_l\varepsilon_l$. Por la
independencia de los errores, $\Cov(\varepsilon_l,\varepsilon_k)$ vale
$\sigma^2$ si $l=k$ y $0$ si no, de modo que
$\Cov(\bar y_i,\varepsilon_k)=\sigma^2/n_i$ (pues $k\in L_i$) y
$\Cov(\bar y,\varepsilon_k)=\sigma^2/n$; $S_i$ es la combinación lineal
de $\bar y_i$ y $\bar y$ con coeficientes $\lambda_i$ y $1-\lambda_i$. Para
una observación nueva, su error no aparece en $\bar y_i$ ni en $\bar y$ y
es independiente de todos los que aparecen.
\end{proof}

En entrenamiento, la característica contiene parte del error de la propia
observación y el modelo aprende a apoyarse en ella; en producción esa parte
no existe. Con $n_i=1$ y $m=0$, $S_{i(k)}=y_k$: la característica es la
respuesta.

\begin{contraejemplo}[\emph{Leave-one-out}]
Codificar cada observación con la media de las demás de su nivel,
$S_{-k}=(T_i-y_k)/(n_i-1)$ con $T_i=\sum_{l\in L_i}y_l$ y $n_i\ge2$,
elimina la covarianza anterior, porque $\varepsilon_k$ ya no interviene.
Pero dentro del nivel $S_{-k}$ es una función afín estrictamente
decreciente de $y_k$, con los mismos $T_i$ y $n_i$ para todas las
observaciones del nivel. En un nivel con respuestas $(1,1,0,0)$, las
observaciones con $y=1$ reciben $S_{-k}=1/3$ y las de $y=0$, $S_{-k}=2/3$:
el corte $S_{-k}\le1/2$ las clasifica todas correctamente en
entrenamiento, mientras que en producción toda observación del nivel
recibe la media del nivel, $1/2$, y el corte no aporta nada.
\end{contraejemplo}

\begin{proposicion}[Codificación cruzada]\label{prop:crossfit}
Particiónese el conjunto de índices de entrenamiento en $V$ pliegues
$F_1,\dots,F_V$ y, para $k\in F_v$, sea $S^{\mathrm{cf}}_k$ la codificación
del nivel de $k$ calculada, con la fórmula de la definición, usando solo las
observaciones fuera de $F_v$. Bajo el modelo de la
Proposición~\ref{prop:te-fuga} y con los pliegues fijos:
(a) $S^{\mathrm{cf}}_k$ es independiente de $\varepsilon_k$; (b) dos
observaciones del mismo pliegue y el mismo nivel reciben el mismo valor.
\end{proposicion}

\begin{proof}
(a) $S^{\mathrm{cf}}_k$ es función de las constantes $\theta_i$ y de los
errores $\{\varepsilon_l:l\notin F_v\}$, independientes de
$\varepsilon_k$ porque $k\in F_v$. (b) Si $k,k'\in F_v$ tienen el mismo
nivel, ambas codificaciones se calculan con las mismas observaciones y la
misma fórmula.
\end{proof}

La parte (b) es la que falla en \emph{leave-one-out}: dentro de un pliegue
la codificación de un nivel es constante y no puede ordenar las respuestas
de ese nivel. Este esquema es el de \texttt{TargetEncoder.fit\_transform}
de scikit-learn ($V=5$ por defecto); por eso \texttt{fit(X,y).transform(X)}
no coincide con \texttt{fit\_transform(X,y)}. En validación, prueba y
producción se usa \texttt{transform}, con la codificación calculada sobre
todo el entrenamiento, independiente del error de la observación nueva
(Proposición~\ref{prop:te-fuga}).

\begin{observacion}[Estadísticos ordenados de CatBoost]
Prokhorenkova et al.\ (2018) toman una permutación aleatoria
$\varsigma$ de los índices de entrenamiento, una constante $\mu_0$ fijada
de antemano, y codifican la observación $k$ con
\[
S^{\mathrm{ord}}_k=\frac{\sum_{l:\,\varsigma(l)<\varsigma(k),\,i(l)=i(k)}y_l+m\,\mu_0}{\#\{l:\varsigma(l)<\varsigma(k),\,i(l)=i(k)\}+m}.
\]
$S^{\mathrm{ord}}_k$ solo depende de respuestas de observaciones
anteriores a $k$ en la permutación, y el argumento de la
Proposición~\ref{prop:crossfit}(a) da su independencia de $\varepsilon_k$.
A cambio, las primeras observaciones de la permutación se codifican con
pocos datos; CatBoost usa varias permutaciones para reducir esa varianza.
\end{observacion}

\subsubsection{El orden de las medias y los árboles}

El resultado siguiente, de Fisher (1958) para regresión y de Breiman et
al.\ (1984) para clasificación binaria, explica por qué el \emph{target
encoding} es adecuado para árboles.

\begin{lema}[Descomposición de la suma de cuadrados]\label{lem:anova}
Sea $N$ un conjunto de índices particionado en $N_A$ y $N_B$ no vacíos,
de tamaños $n_A,n_B$ y medias $\bar y_A,\bar y_B$, y sea $\bar y_N$ la
media en $N$. Entonces
$\SC(N)=\SC(N_A)+\SC(N_B)+n_A(\bar y_A-\bar y_N)^2+n_B(\bar y_B-\bar y_N)^2$.
\end{lema}

\begin{proof}
Para $k\in N_A$, $y_k-\bar y_N=(y_k-\bar y_A)+(\bar y_A-\bar y_N)$; al
elevar al cuadrado y sumar sobre $N_A$, el término cruzado es
$2(\bar y_A-\bar y_N)\sum_{k\in N_A}(y_k-\bar y_A)=0$. Se obtiene
$\SC(N_A)+n_A(\bar y_A-\bar y_N)^2$; lo mismo para $N_B$, y se suma.
\end{proof}

Así, la reducción $\Delta$ de la Definición~\ref{def:cart} para una
partición de $N$ en $N_A$ y $N_B$ es
$\mathcal B=n_A(\bar y_A-\bar y_N)^2+n_B(\bar y_B-\bar y_N)^2$.

\begin{teorema}\label{teo:fisher}
Sea $N$ un nodo en el que aparecen los niveles $c_1,\dots,c_K$, con
$n_i>0$ observaciones y medias $\bar y_i$ en $N$. Para $A\subset\{1,\dots,K\}$
no vacío y distinto del total, sea $N_A$ el conjunto de observaciones de
$N$ con nivel en $A$, $N_B=N\setminus N_A$, y $\mathcal B(A)$ la cantidad
anterior. Existe un maximizador de $\mathcal B$ formado por los primeros
niveles de una ordenación por $\bar y_i$ no decreciente. En consecuencia,
examinando los $K-1$ cortes de la codificación $\bar y_i$ (calculada en
el nodo) se encuentra la mejor de las $2^{K-1}-1$ particiones.
\end{teorema}

\begin{proof}
Si todas las $\bar y_i$ son iguales, $\bar y_A=\bar y_B=\bar y_N$ para
todo $A$ y todo $A$ es maximizador. En otro caso, sea $A$ un maximizador.
El conjunto $A_0$ de niveles con media mínima cumple
$\bar y_{A_0}<\bar y_N$, luego $\mathcal B(A_0)>0$; por tanto
$\mathcal B(A)>0$ y $\bar y_A\ne\bar y_B$ (si fueran iguales, ambas
coincidirían con $\bar y_N$ y $\mathcal B(A)=0$). Intercambiando los
nombres si hace falta, $\bar y_A<\bar y_B$. Sean $\delta_i=\bar y_i-\bar y_N$,
$\alpha=\bar y_A-\bar y_N$, $\beta=\bar y_B-\bar y_N$ (luego
$\alpha<\beta$), $s_A=n_A\alpha$, $s_B=n_B\beta$ y $G(s,w)=s^2/w$ para
$w>0$, de modo que $\mathcal B(A)=G(s_A,n_A)+G(s_B,n_B)$.

\emph{Convexidad de $G$.} Su hessiana es
$\begin{psmallmatrix}2/w&-2s/w^2\\-2s/w^2&2s^2/w^3\end{psmallmatrix}$, de
traza positiva y determinante $4s^2/w^4-4s^2/w^4=0$; sus autovalores son
$0$ y un número positivo, y $G$ es convexa en $\{w>0\}$.

\emph{Traslado parcial de un nivel.} Para $i\in A$ y $t\in[0,1]$, sea
$\Psi(t)=G(s_A-tn_i\delta_i,\,n_A-tn_i)+G(s_B+tn_i\delta_i,\,n_B+tn_i)$. Como el nivel
$i$ aporta $n_i$ observaciones y $n_i\delta_i$ a la suma
$s_A=\sum_{k\in N_A}(y_k-\bar y_N)$, $\Psi(1)$ es el valor de $\mathcal
B$ con el nivel $i$ trasladado a $B$, es decir, $\mathcal B(A\setminus\{i\})$
(o $0$ si $A=\{i\}$, por la continuidad que se ve enseguida), y
$\Psi(0)=\mathcal B(A)$. Cada sumando es $G$ compuesta con una función
afín de $t$, así que $\Psi$ es convexa en $[0,1)$; si $A=\{i\}$, el primer
sumando vale $(1-t)n_i\delta_i^2$ para $t<1$ y tiende a $0$, y definiendo
$\Psi(1)$ por continuidad la convexidad se conserva en $[0,1]$. La
maximalidad de $A$ da $\Psi(1)\le\Psi(0)$, y por convexidad la derivada
por la derecha en $0$ cumple $\Psi'(0^+)\le\Psi(1)-\Psi(0)\le0$.

\emph{Cálculo.} $\partial_sG=2s/w$ y $\partial_wG=-s^2/w^2$ valen $2\alpha$ y
$-\alpha^2$ en $(s_A,n_A)$, y $2\beta$ y $-\beta^2$ en $(s_B,n_B)$. Por la
regla de la cadena,
\[
\Psi'(0^+)=n_i\bigl(-2\alpha \delta_i+\alpha^2+2\beta \delta_i-\beta^2\bigr)=n_i(\beta-\alpha)(2\delta_i-\alpha-\beta).
\]
Como $\beta>\alpha$, $\Psi'(0^+)\le0$ equivale a $\bar y_i\le\tau$, con
$\tau=\bar y_N+\frac{\alpha+\beta}2=\frac{\bar y_A+\bar y_B}2$. El mismo
cálculo trasladando un nivel $j\in B$ hacia $A$ da
$n_j(\beta-\alpha)(\alpha+\beta-2\delta_j)\le0$, es decir, $\bar y_j\ge\tau$.

\emph{Conclusión.} Los niveles de $A$ tienen media $\le\tau$ y los de $B$,
media $\ge\tau$. Ordenando por $\bar y_i$ y colocando, entre los empatados
en $\tau$, primero los de $A$, el conjunto $A$ ocupa las primeras
posiciones. Si $y\in\{0,1\}$, $\SC(N)=|N|\hat p(1-\hat p)$ con $\hat p$ la
proporción de unos en $N$, que es $|N|/2$ veces el índice de Gini
$2\hat p(1-\hat p)$; el criterio de Gini es entonces proporcional al de
suma de cuadrados, y el resultado vale igual.
\end{proof}

Con una codificación ordinal en un orden arbitrario, un corte examina
$K-1$ de las $2^{K-1}-1$ particiones y en general no la mejor; con el
orden de las medias, esas $K-1$ contienen a la mejor. LightGBM usa esta
idea para sus cortes categóricos. El teorema se refiere a las medias
\emph{del nodo}; una columna codificada una vez con toda la muestra las
da exactamente solo en la raíz.

\paragraph{Ventajas y desventajas.} A favor: dimensión $1$ para cualquier
$K$; contracción óptima bajo los modelos jerárquicos
(Proposiciones~\ref{prop:betabin} y~\ref{prop:normnorm}); valor definido
para niveles nuevos; orden óptimo para cortes de árboles
(Teorema~\ref{teo:fisher}). En contra: usa la respuesta, así que requiere
codificación cruzada en entrenamiento y ajuste sin la muestra de prueba
(Proposiciones~\ref{prop:test} y~\ref{prop:te-fuga}); un predictor lineal
sobre $S$ solo representa efectos $\beta_0+\beta S_i$
(Proposición~\ref{prop:repr} con $d=1$); no capta interacciones del nivel
con otras variables; y la tabla de codificación es un artefacto que debe
versionarse y aplicarse igual en producción.

\begin{aws}
Ni Data Wrangler (Canvas) ni Glue DataBrew incluyen \emph{target encoding}
entre sus transformaciones (documentación revisada el 24 de septiembre de
2026). Un \emph{custom transform} de Data Wrangler que calcule medias por
nivel lo haría con todos los datos del flujo, sin pliegues, en la situación
de la Proposición~\ref{prop:te-fuga}. Tres vías respetan los resultados
anteriores.
\begin{enumerate}[leftmargin=*]
\item \textbf{SageMaker Processing} con un script que use
\texttt{TargetEncoder} (scikit-learn $\ge1.3$; se verifica la versión del
contenedor o se instala):
\begin{lstlisting}[language=Python]
from sklearn.preprocessing import TargetEncoder
import joblib, pandas as pd
tr = pd.read_parquet("/opt/ml/processing/train/")
te = pd.read_parquet("/opt/ml/processing/test/")
enc = TargetEncoder(smooth="auto", cv=5, target_type="binary")
# entrenamiento: codificación cruzada
tr["cp_te"] = enc.fit_transform(tr[["codigo_postal"]], tr["y"]).ravel()
# prueba: tabla calculada con todo el entrenamiento
te["cp_te"] = enc.transform(te[["codigo_postal"]]).ravel()
joblib.dump(enc, "/opt/ml/processing/model/target_encoder.joblib")
\end{lstlisting}
El codificador guardado en S3 se carga en el contenedor de inferencia (por
ejemplo, el primero de un \emph{inference pipeline}).
\item \textbf{Algoritmo integrado CatBoost} de SageMaker, que aplica
internamente los estadísticos ordenados. Las columnas categóricas se
declaran en \texttt{categorical\_index.json}, junto a los datos de
entrenamiento (CSV sin encabezado con la respuesta en la primera columna;
el canal de entrenamiento debe apuntar a un prefijo de S3):
\begin{lstlisting}
{"cat_index_list": [1, 3, 5]}
\end{lstlisting}
\item \textbf{AWS Glue} (PySpark) para volúmenes grandes: se asigna un
pliegue a cada fila, se calculan con \texttt{groupBy} los conteos y sumas
por nivel y pliegue, y cada fila se codifica con las sumas de los demás
pliegues. La tabla final (nivel, $n_i$, $S_i$) puede publicarse en un
\emph{feature group} de \textbf{Feature Store} con el nivel como
identificador, para consultarla desde el \emph{online store} en inferencia.
\end{enumerate}
\end{aws}

\subsection{\emph{Weight of evidence} e \emph{information value}}\label{sec:woe}

Turing y Good (1950) llamaron peso de la evidencia al logaritmo del
cociente de verosimilitudes que una observación aporta a favor de una
hipótesis frente a otra. La banca lo adoptó para construir
\emph{scorecards} de crédito, que debían ser modelos aditivos, con una
contribución por variable legible por un analista; la forma que sigue es la
de Siddiqi (2006).

\paragraph{Notación.} La respuesta es binaria: $Y=1$ denota el evento
(«malo») e $Y=0$ su ausencia («bueno»). $X$ es categórica o una variable
continua discretizada (\S\ref{sec:discretizacion}). En entrenamiento,
$n_i^{(0)}$ y $n_i^{(1)}$ son los números de buenos y de malos con nivel
$c_i$, y $N^{(0)}=\sum_in_i^{(0)}$, $N^{(1)}=\sum_in_i^{(1)}$.

\begin{definicion}
Supóngase $n_i^{(0)}>0$ y $n_i^{(1)}>0$ para todo $i$. Con
$g_i=n_i^{(0)}/N^{(0)}$ y $b_i=n_i^{(1)}/N^{(1)}$ (las distribuciones de
$X$ entre buenos y entre malos),
\[
\WoE_i=\log\frac{g_i}{b_i},\qquad\IV=\sum_{i=1}^K(g_i-b_i)\WoE_i .
\]
\end{definicion}

Algunos textos usan el cociente inverso, lo que cambia el signo del WoE y
no el IV. La hipótesis de conteos positivos es necesaria: un nivel sin
malos tendría $\WoE=+\infty$. Las soluciones habituales son sumar $0.5$ a
cada conteo o fusionar el nivel con otro; ambas son formas de suavizado,
comparables con $m$ en el \emph{target encoding}.

\begin{proposicion}[WoE y log-momios]\label{prop:woe-bayes}
En la población, sean $\pi=\Prob(Y=1)\in(0,1)$, $G_i=\Prob(X=c_i\mid Y=0)>0$,
$B_i=\Prob(X=c_i\mid Y=1)>0$ y $W_i=\log(G_i/B_i)$. Entonces
$\logit\Prob(Y=1\mid X=c_i)=\logit\pi-W_i$, donde $\logit
u=\log(u/(1-u))$. En la muestra, con $\hat\pi=N^{(1)}/(N^{(0)}+N^{(1)})$ y
$\bar y_i=n_i^{(1)}/(n_i^{(0)}+n_i^{(1)})$,
$\WoE_i=\logit\hat\pi-\logit\bar y_i$.
\end{proposicion}

\begin{proof}
Por la regla de Bayes,
$\frac{\Prob(Y=1\mid X=c_i)}{\Prob(Y=0\mid X=c_i)}=\frac{\pi B_i}{(1-\pi)G_i}$;
tomando logaritmos, $\logit\Prob(Y=1\mid X=c_i)=\logit\pi-\log(G_i/B_i)$.
En la muestra,
$\frac{\bar y_i}{1-\bar y_i}=\frac{n_i^{(1)}}{n_i^{(0)}}=\frac{b_iN^{(1)}}{g_iN^{(0)}}
=\frac{b_i}{g_i}\cdot\frac{\hat\pi}{1-\hat\pi}$, y se toman logaritmos.
\end{proof}

\begin{corolario}\label{cor:woe-te}
$\WoE_i$ es una función estrictamente decreciente de $\bar y_i$, el
\emph{target encoding} sin suavizar. Ambas codificaciones ordenan los
niveles de la misma manera (en sentido inverso), el
Teorema~\ref{teo:fisher} se aplica igual a las dos, y un árbol produce las
mismas hojas con cualquiera de ellas.
\end{corolario}

\begin{proof}
$\logit$ es estrictamente creciente en $(0,1)$ y
$\WoE_i=\logit\hat\pi-\logit\bar y_i$, luego $-\WoE_i$ es una función
estrictamente creciente de $\bar y_i$ y la Proposición~\ref{prop:arbol} da
las mismas hojas con $\bar y$ y con $-\WoE$. Falta ver que $\WoE$ y
$-\WoE$ dan las mismas hojas. Si $v_1<\dots<v_r$ son los valores distintos
de una columna en un nodo, sus cortes producen las particiones
$\{\{k:v\le v_i\},\{k:v>v_i\}\}$, $i<r$ (Lema~\ref{lem:cortes}); los
de la columna opuesta producen $\{\{k:v\ge v_{i+1}\},\{k:v<v_{i+1}\}\}$,
que como particiones no ordenadas son las mismas. Como el criterio y el
desempate de la Definición~\ref{def:cart} solo dependen de la partición no
ordenada, el argumento de la Proposición~\ref{prop:arbol} se aplica
sin cambios.
\end{proof}

La diferencia está en la escala: el \emph{target encoding} está en la
escala de probabilidad y el WoE en la de log-momios, que es la escala
lineal de la regresión logística.

\begin{proposicion}[Regresión logística sobre el WoE]\label{prop:woe-logit}
Supóngase $0<\bar y_i<1$ para todo $i$. En la regresión logística de $y$
sobre la columna $\WoE_{i(k)}$, los coeficientes
$(\hat\beta_0,\hat\beta_1)=(\logit\hat\pi,-1)$ maximizan la verosimilitud,
y las probabilidades ajustadas son las proporciones $\bar y_i$.
\end{proposicion}

\begin{proof}
Todo modelo que asigne a las observaciones del nivel $c_i$ una misma
probabilidad $p_i$ tiene log-verosimilitud
$\sum_in_i\bigl(\bar y_i\log p_i+(1-\bar y_i)\log(1-p_i)\bigr)$, con
$n_i=n_i^{(0)}+n_i^{(1)}$. Cada sumando es una función estrictamente cóncava
de $p_i\in(0,1)$ con derivada nula en $p_i=\bar y_i$, así que el máximo
sobre todos esos modelos se alcanza en $p_i=\bar y_i$. La regresión
logística sobre la columna WoE asigna a cada nivel una probabilidad común,
luego su verosimilitud no supera ese máximo; con
$(\logit\hat\pi,-1)$ da $\logit p_i=\logit\hat\pi-\WoE_i=\logit\bar y_i$
(Proposición~\ref{prop:woe-bayes}) y lo alcanza.
\end{proof}

\begin{proposicion}[WoE y clasificador bayesiano ingenuo]\label{prop:woe-nb}
Sean $X_1,\dots,X_p$ categóricas condicionalmente independientes dado
$Y$, con $\Prob(X_j=x_j\mid Y=y)>0$, y
$W_j(x_j)=\log\frac{\Prob(X_j=x_j\mid Y=0)}{\Prob(X_j=x_j\mid Y=1)}$.
Entonces
\[
\logit\Prob(Y=1\mid X_1=x_1,\dots,X_p=x_p)=\logit\pi-\sum_{j=1}^pW_j(x_j).
\]
\end{proposicion}

\begin{proof}
Por la independencia condicional,
$\Prob(X_1=x_1,\dots,X_p=x_p\mid Y=y)=\prod_j\Prob(X_j=x_j\mid Y=y)$. Por
la regla de Bayes, el cociente $\Prob(Y=1\mid x)/\Prob(Y=0\mid x)$ es
$\frac{\pi}{1-\pi}\prod_j\frac{\Prob(X_j=x_j\mid Y=1)}{\Prob(X_j=x_j\mid Y=0)}$,
cuyo logaritmo es $\logit\pi-\sum_jW_j(x_j)$.
\end{proof}

Una regresión logística sobre las columnas $\WoE_j$ con coeficientes
libres $\beta_j$ contiene como caso particular ($\beta_j=-1$ para todo $j$)
al clasificador bayesiano ingenuo; coeficientes alejados de $-1$ corrigen
la dependencia entre variables. Esa aditividad, un término por variable en
la escala de log-momios, es la que permite traducir el modelo a una tabla
de puntos.

\begin{proposicion}[El IV es una divergencia simétrica]\label{prop:iv}
$\IV=\mathrm{KL}(g\|b)+\mathrm{KL}(b\|g)$, donde
$\mathrm{KL}(g\|b)=\sum_ig_i\log(g_i/b_i)$ es la divergencia de
Kullback--Leibler. En consecuencia $\IV\ge0$, con igualdad si y solo si
$g_i=b_i$ para todo $i$. Además, cada sumando $(g_i-b_i)\WoE_i$ es $\ge0$.
\end{proposicion}

\begin{proof}
La primera igualdad es separar $\sum_i(g_i-b_i)\log(g_i/b_i)$ en
$\sum_ig_i\log(g_i/b_i)+\sum_ib_i\log(b_i/g_i)$. Por la desigualdad de
Jensen para la función cóncava $\log$ y la distribución $(g_i)$,
$-\mathrm{KL}(g\|b)=\sum_ig_i\log(b_i/g_i)\le\log\sum_ig_i(b_i/g_i)=\log\sum_ib_i=0$,
con igualdad si y solo si $b_i/g_i$ es constante, y la constante es $1$
porque $\sum_ig_i=\sum_ib_i=1$; lo mismo para $\mathrm{KL}(b\|g)$. En
cuanto al signo, $g_i-b_i$ y $\log(g_i/b_i)$ son ambos positivos, ambos
negativos o ambos nulos.
\end{proof}

La condición $g_i=b_i$ para todo $i$ significa que en la muestra la
distribución de $X$ es la misma entre buenos y malos: independencia
empírica. De aquí sale un sesgo del IV estimado. Supóngase que en la
población $X$ e $Y$ son independientes, de modo que el IV poblacional es
$0$, que $K\ge2$ y que todas las combinaciones (nivel, clase) tienen
probabilidad positiva. El IV muestral es $\ge0$, y hay configuraciones de
probabilidad positiva en que no vale $0$ (por ejemplo, todos los buenos en
$c_1$ y todos los malos en $c_2$), así que su esperanza es estrictamente
positiva. Como cada nivel aporta un sumando no negativo, más niveles dan
más ocasiones de sumar ruido positivo. Los umbrales usados en la industria
(Siddiqi: $<0.02$ inútil, $0.02$--$0.1$ débil, $0.1$--$0.3$ medio, $>0.3$
fuerte) son convenciones empíricas.

\begin{ejemplo}\label{ej:woe}
Mil solicitudes con 100 malos ($\hat\pi=0.1$) y una variable con tres
niveles:
\begin{center}
\begin{tabular}{@{}lrrrrrrr@{}}
\toprule
Nivel & $n_i$ & malos & $\bar y_i$ & $g_i$ & $b_i$ & $\WoE_i$ & $(g_i-b_i)\WoE_i$\\
\midrule
A & 500 & 30 & 0.060 & 0.5222 & 0.30 & $0.554$ & 0.123\\
B & 400 & 50 & 0.125 & 0.3889 & 0.50 & $-0.251$ & 0.028\\
C & 100 & 20 & 0.200 & 0.0889 & 0.20 & $-0.811$ & 0.090\\
\midrule
&&&&&&$\IV$&0.241\\
\bottomrule
\end{tabular}
\end{center}
Para A, $g=470/900$ y $b=30/100$. Se verifica la
Proposición~\ref{prop:woe-bayes}:
$\logit0.1-\logit0.06=-2.197-(-2.752)=0.554$. El \emph{target encoding}
con $m=10$ da $(0.0608,\,0.1244,\,0.1909)$: el mismo orden, con el nivel
pequeño C más contraído hacia $0.1$; con $m=100$, $(0.0667,\,0.12,\,0.15)$.
\end{ejemplo}

\paragraph{Contraste entre \emph{target encoding} y WoE.}
\begin{center}
\begin{tabularx}{\textwidth}{@{}lYY@{}}
\toprule
& \emph{Target encoding} & WoE\\
\midrule
Respuesta & binaria, multiclase o continua & binaria\\
Escala & probabilidad o media & log-momios\\
Suavizado & media posterior; $m$ estimable (Prop.~\ref{prop:normnorm}) & $+0.5$ en los conteos o fusión de niveles\\
Modelo natural & árboles (Teor.~\ref{teo:fisher}) y cualquier otro & regresión logística (Prop.~\ref{prop:woe-logit} y~\ref{prop:woe-nb})\\
Variables continuas & no se aplica & tras discretizar\\
Diagnóstico asociado & --- & IV por variable\\
Fuga & sí; requiere codificación cruzada & sí; el mismo tratamiento\\
\bottomrule
\end{tabularx}
\end{center}
Ambas inducen el mismo orden de niveles (Corolario~\ref{cor:woe-te}).

\begin{aws}
No hay transformación de WoE ni de IV en Data Wrangler, DataBrew ni
Canvas, y scikit-learn tampoco la incluye. Una implementación consta de:
(i) discretizar las continuas con \texttt{BUCKETIZATION} de DataBrew o
\texttt{KBinsDiscretizer}; (ii) un trabajo de \textbf{SageMaker
Processing} (o de \textbf{AWS Glue} en PySpark para volúmenes grandes) que
calcule con entrenamiento la tabla de conteos y el WoE, con $0.5$ sumado a
cada conteo, y la guarde en S3 o en Feature Store:
\begin{lstlisting}[language=Python]
import numpy as np, pandas as pd
def tabla_woe(x, y, eps=0.5):
    t = pd.crosstab(x, y)                  # columnas: 0 (buenos), 1 (malos)
    g = (t[0] + eps) / (t[0] + eps).sum()
    b = (t[1] + eps) / (t[1] + eps).sum()
    woe = np.log(g / b)
    return woe, ((g - b) * woe).sum()     # (WoE por nivel, IV)
woe, iv = tabla_woe(train["nivel"], train["malo"])
train["nivel_woe"] = train["nivel"].map(woe)
test["nivel_woe"]  = test["nivel"].map(woe).fillna(0.0)
\end{lstlisting}
A un nivel no visto se le asigna $0$: por la
Proposición~\ref{prop:woe-bayes}, $\WoE=0$ deja la probabilidad en su
valor a priori, es decir, no aporta evidencia. (iii) La regresión
logística, con el algoritmo integrado \textbf{Linear Learner}
(\texttt{predictor\_type="binary\_classifier"}) o con scikit-learn. Con
alta cardinalidad se usa la codificación cruzada de la
Proposición~\ref{prop:crossfit}.
\end{aws}

\subsection{Resumen de codificaciones categóricas}

\begin{center}
\small
\begin{tabularx}{\textwidth}{@{}lcccY@{}}
\toprule
Codificación & $d$ & $\rango[\1\;\Phi]$ & Usa $y$ & Observaciones\\
\midrule
Ordinal & 1 & 2 & no & efectos lineales en el orden elegido\\
\emph{One-hot} & $K$ ($K-1$) & $K$ & no & saturada; niveles equidistantes; varianza $\sigma^2/n_i$\\
Binaria & $\lceil\log_2(K+1)\rceil$ & $\le d+1$ & no & inyectiva; similitudes fijadas por la numeración\\
\emph{Hashing} & $m$ & $\le m+1$ & no & colisiones; sin diccionario; símbolos nuevos sin costo\\
Similitud & $d$ fijo & --- & no & coincidencias según 3-gramas comunes\\
\emph{Target encoding} & 1 (o $C$) & 2 & sí & orden óptimo para árboles; codificación cruzada\\
WoE & 1 & 2 & sí & escala logit; respuesta binaria\\
\bottomrule
\end{tabularx}
\end{center}

% =====================================================================
\section{Series de tiempo}
% =====================================================================

Sea $(y_t)_{t\in\Z}$ un proceso real observado en instantes
equiespaciados $t=1,\dots,T$, y sea $\mathcal H_t=(y_s)_{s<t}$ su
\emph{historia} antes de $t$. Para predecir $y_t$ solo pueden usarse
funciones de $\mathcal H_t$ y funciones conocidas de $t$ (el calendario).
Una característica que dependa de $y_t$ o de valores posteriores es fuga de
información: en producción no existe cuando hay que predecir.

\subsection{Descomposición de la marca de tiempo}

De la marca de tiempo de $t$ se extraen componentes de calendario: mes,
día del mes, día de la semana, hora, trimestre. Cada una es una variable
\emph{cíclica}: toma valores $u\in\{0,1,\dots,P-1\}$ y después de $P-1$
vuelve a $0$ (mes: $P=12$; hora: $P=24$). La codificación ordinal
$\phi(u)=u$ coloca a los valores $P-1$ y $0$ (diciembre y enero) a la
máxima distancia, aunque son consecutivos.

\begin{definicion}[Codificación cíclica]
$e(u)=\bigl(\sin(2\pi u/P),\cos(2\pi u/P)\bigr)\in\R^2$. La
\emph{distancia circular} entre $u$ y $v$ es
$\delta(u,v)=\min(|u-v|,\,P-|u-v|)$.
\end{definicion}

\begin{proposicion}\label{prop:ciclica}
$\|e(u)-e(v)\|=2\sin\bigl(\pi\,\delta(u,v)/P\bigr)$; en particular, la
distancia euclídea entre codificaciones es una función estrictamente
creciente de la distancia circular.
\end{proposicion}

\begin{proof}
Con $a=2\pi u/P$ y $b=2\pi v/P$,
$\|e(u)-e(v)\|^2=2-2(\sin a\sin b+\cos a\cos b)=2-2\cos(a-b)=4\sin^2\bigl(\tfrac{a-b}2\bigr)$,
de modo que $\|e(u)-e(v)\|=2|\sin(\pi(u-v)/P)|$. Sea $\Delta=|u-v|\in[0,P)$.
Como $\sin(\pi-x)=\sin x$, $\sin(\pi\Delta/P)=\sin(\pi(P-\Delta)/P)$, así que
el valor coincide en $\Delta$ y en $P-\Delta$, y por tanto en
$\delta(u,v)\in[0,P/2]$; en $[0,P/2]$, $\delta\mapsto\sin(\pi\delta/P)$ es
no negativa y estrictamente creciente.
\end{proof}

Con meses, diciembre y enero quedan a distancia $2\sin(\pi/12)\approx0.52$;
enero y junio, a $2\sin(5\pi/12)\approx1.93$. Un predictor lineal sobre
$e(u)$ es $\beta_0+\beta_1\sin a+\beta_2\cos a$ con $a=2\pi u/P$; con
$R=\sqrt{\beta_1^2+\beta_2^2}$ y $\varphi$ tal que $\cos\varphi=\beta_2/R$ y
$\sin\varphi=\beta_1/R$ (si $R>0$), es igual a $\beta_0+R\cos(a-\varphi)$
por la fórmula del coseno de una diferencia: una sola sinusoide de período
$P$. Perfiles estacionales más generales requieren más armónicos,
$(\sin(2\pi hu/P),\cos(2\pi hu/P))$ con $h=2,3,\dots$, o \emph{one-hot} de
$u$, saturado con $P-1$ columnas (Proposición~\ref{prop:repr}).

Para árboles la codificación cíclica es incómoda. Como
$\sin(2\pi(P/2-u)/P)=\sin(\pi-2\pi u/P)=\sin(2\pi u/P)$, los valores $u$ y
$P/2-u$ (módulo $P$) tienen el mismo seno, y ningún corte sobre esa
coordenada los separa; análogamente, $u$ y $P-u$ tienen el mismo coseno.
Con $P=12$, los cortes sobre el seno no separan enero ($u=1$) de mayo
($u=5$), y los cortes sobre el coseno no separan enero de noviembre
($u=11$): aislar enero exige cortar en las dos coordenadas. Con la
codificación ordinal bastan dos cortes sobre una. La documentación de
Data Wrangler recomienda el modo cíclico para modelos lineales y redes, y
el ordinal para árboles, en coherencia con este análisis.

\begin{aws}
\textbf{Data Wrangler}, \emph{Featurize datetime}: acepta cadenas de fecha
o marcas Unix, infiere o recibe el formato (\texttt{strftime}) y produce
los componentes en modo \emph{Cyclic} u \emph{Ordinal}, como vector o en
columnas. Para tratar un componente como categórico se encadena
\emph{Encode categorical} $\to$ \emph{One-hot encode} (por ejemplo, el
trimestre en cuatro columnas). En \textbf{SageMaker Processing}, con
\texttt{pandas.Series.dt} y las fórmulas anteriores.
\end{aws}

\subsection{Características de rezago}

\begin{definicion}
El \emph{rezago} de orden $h\ge1$ en el instante $t$ es $y_{t-h}$; las
características de rezago de orden $1$ a $p$ son
$(y_{t-1},\dots,y_{t-p})$, y los rezagos estacionales son
$y_{t-P},y_{t-2P},\dots$
\end{definicion}

Su justificación proviene de los modelos autorregresivos.

\begin{proposicion}\label{prop:ar}
Supóngase $\E[y_t^2]<\infty$ y
$y_t=\phi_0+\sum_{h=1}^p\phi_hy_{t-h}+\varepsilon_t$, con $\varepsilon_t$
de media $0$ e independiente de $\mathcal H_t$. Entre todas las funciones
$f$ de la historia con $\E[f^2]<\infty$, la que minimiza
$\E[(y_t-f)^2]$ es $f^*=\phi_0+\sum_{h=1}^p\phi_hy_{t-h}$.
\end{proposicion}

\begin{proof}
Sea $m_t=\E[y_t\mid\mathcal H_t]$. Para cualquier $f$ función de la
historia,
$\E[(y_t-f)^2]=\E[(y_t-m_t)^2]+\E[(m_t-f)^2]+2\E[(y_t-m_t)(m_t-f)]$, y el
último término es nulo: condicionando en $\mathcal H_t$, el factor
$m_t-f$ es conocido y $\E[y_t-m_t\mid\mathcal H_t]=0$. Luego
$\E[(y_t-f)^2]\ge\E[(y_t-m_t)^2]$, con igualdad si $f=m_t$. Como los
$y_{t-h}$ son funciones de la historia y $\varepsilon_t$ es independiente
de ella,
$m_t=\phi_0+\sum_h\phi_hy_{t-h}+\E[\varepsilon_t]=f^*$.
\end{proof}

Bajo el modelo, los rezagos $1,\dots,p$ contienen toda la información de
la historia útil para predecir en error cuadrático, y los de orden mayor
no la mejoran. Esto justifica elegir $p$ con la \emph{autocorrelación
parcial}: si el proceso es estacionario (medias y covarianzas invariantes
en el tiempo) con matrices de covarianza no singulares, la autocorrelación
parcial de orden $k$ es el coeficiente de $y_{t-k}$ en el mejor predictor
lineal de $y_t$ a partir de $(1,y_{t-1},\dots,y_{t-k})$. Si $k\ge p+1$,
$f^*$ es lineal en esas variables y óptimo entre todas las funciones de la
historia, así que también es el mejor predictor lineal a partir de ellas;
por la unicidad de este (consecuencia de la no singularidad), el
coeficiente de $y_{t-k}$ es $0$. Las primeras $p$ observaciones de la
serie no tienen historia completa: se descartan o se imputan sus rezagos.

\begin{aws}
\textbf{Data Wrangler}, \emph{Lag features}: genera los rezagos hasta una
duración dada respecto de una columna de tiempo (por serie, si antes se
aplica \emph{Group by}), con opciones para incluir la ventana completa,
aplanarla en columnas y descartar las filas sin historia. Antes conviene
\emph{Validate timestamps} y \emph{Resample}, para que un rezago
corresponda a un intervalo fijo. En pandas,
\texttt{groupby(...).shift(h)}.
\end{aws}

\subsection{Características de ventana móvil}

\begin{definicion}
Para una ventana de tamaño $w$ que termina en $t-1$, la media móvil es
$\bar y^{(w)}_t=w^{-1}\sum_{h=1}^wy_{t-h}$; análogamente se definen el
mínimo, el máximo y la desviación estándar de $(y_{t-w},\dots,y_{t-1})$.
\end{definicion}

Que la ventana termine en $t-1$ es la condición de admisibilidad del
principio de la sección. El tamaño $w$ resuelve un compromiso entre sesgo
y varianza.

\begin{proposicion}\label{prop:ventana}
Sea $y_t=a+bt+\varepsilon_t$ con $\varepsilon_t$ independientes de media
$0$ y varianza $\sigma^2$. Entonces
\[
\E\bigl[(y_t-\bar y^{(w)}_t)^2\bigr]=\frac{b^2(w+1)^2}{4}+\frac{\sigma^2}w+\sigma^2 .
\]
\end{proposicion}

\begin{proof}
Como $\sum_{h=1}^wh=w(w+1)/2$,
$\bar y^{(w)}_t=a+b\bigl(t-\frac{w+1}2\bigr)+\bar\varepsilon$ con
$\bar\varepsilon=w^{-1}\sum_{h=1}^w\varepsilon_{t-h}$. Entonces
$y_t-\bar y^{(w)}_t=\frac{b(w+1)}2+\varepsilon_t-\bar\varepsilon$, donde
$\varepsilon_t$ y $\bar\varepsilon$ son independientes (involucran
errores de instantes distintos), de media $0$, con varianzas $\sigma^2$ y
$\sigma^2/w$. El error cuadrático medio es el cuadrado de la constante más
la suma de las varianzas.
\end{proof}

Sin tendencia ($b=0$) el error decrece con $w$; con tendencia, el término
de sesgo crece como $w^2$. Derivando respecto de $w$ la parte que depende
de $w$, el mínimo cumple $b^2(w+1)/2=\sigma^2/w^2$, y para $w$ grande
$w^*\approx(2\sigma^2/b^2)^{1/3}$: la ventana debe acortarse cuando la
tendencia crece respecto del ruido. El mínimo, el máximo y la desviación
estándar móviles describen el rango y la volatilidad recientes, que la
media no refleja. La extracción automática de muchos estadísticos de
ventana multiplica las columnas y conviene seguirla de selección
(\S\ref{sec:seleccion}).

\paragraph{Contraste.} Los rezagos representan exactamente la dinámica de
un modelo autorregresivo y permiten ponderar cada retardo; su número
crece con la memoria que se quiera capturar. Las ventanas resumen muchos
retardos en pocos números, con menos varianza y pérdida de detalle. Las
componentes de fecha no dependen de $y$, así que no producen fuga y existen
para cualquier horizonte; un rezago de orden $h$ solo está disponible
para predecir hasta $h$ pasos adelante sin usar predicciones intermedias.

\begin{aws}
\textbf{Data Wrangler}, \emph{Rolling window features}: tamaño de ventana,
columna de tiempo y conjunto de estadísticos (mínimo, eficiente, completo
o elegido a mano), con opción de aplanar; \emph{Extract features} calcula
estadísticos de la serie completa. \emph{Split data} en modo
\emph{Ordered} deja la prueba después del entrenamiento. En pandas,
\texttt{y.shift(1).rolling(w).mean()}: el \texttt{shift(1)} hace que la
ventana termine en $t-1$.
\end{aws}

% =====================================================================
\section{Imágenes}
% =====================================================================

Una imagen digital de $H\times W$ píxeles y $C$ canales (tres en color:
rojo, verde, azul) es un arreglo $I\in\{0,\dots,255\}^{H\times W\times C}$;
como vector, está en $\R^{HWC}$, y con $224\times224$ píxeles en color
tiene $150\,528$ coordenadas. Usar esas coordenadas como columnas tiene dos
problemas: la dimensión (\S\ref{sec:pca}) y la sensibilidad a
desplazamientos, porque mover un objeto un píxel cambia casi todas las
coordenadas. Las técnicas de esta sección atacan uno u otro.

\subsection{Normalización de píxeles}

Dividir por $255$ es un escalado MinMax (Definición~\ref{def:escala}) con
los límites posibles $0$ y $255$ en lugar de los observados: no se estima
nada, no hay fuga y un píxel extremo no comprime a los demás (a
diferencia de la Proposición~\ref{prop:minmax}). La estandarización por
canal, $(I_c/255-\mu_c)/\sigma_c$, usa medias y desviaciones fijas de un
corpus de referencia. Con una red preentrenada deben usarse las del
preentrenamiento (en ImageNet, $\mu=(0.485,0.456,0.406)$ y
$\sigma=(0.229,0.224,0.225)$): la red se ajustó sobre entradas con esa
escala, y otra escala en inferencia es un caso de \emph{training-serving
skew}.

\subsection{Extracción clásica: convolución y gradientes}

Para enunciar las propiedades sin efectos de borde, se considera una
imagen de un canal como una función $I:\Z^2\to\R$ nula fuera de un
conjunto finito (la imagen real, extendida con ceros).

\begin{definicion}
Sea $K:\Z^2\to\R$ nula fuera de un conjunto finito (el \emph{núcleo}). La
convolución, en la convención usada en redes neuronales, es
$(I*K)(p)=\sum_{q\in\Z^2}I(p+q)K(q)$. El núcleo de Sobel horizontal es
$K_x(q_1,q_2)=q_1w(q_2)$ para $q_1,q_2\in\{-1,0,1\}$, con $w(\pm1)=1$ y
$w(0)=2$, y $0$ en otro caso; en forma de matriz,
$\begin{psmallmatrix}-1&0&1\\-2&0&2\\-1&0&1\end{psmallmatrix}$.
\end{definicion}

$K_x$ combina una diferencia centrada en la primera coordenada con un
promedio ponderado $(1,2,1)$ en la segunda. En un borde vertical,
$I(p)=\ind{p_1\ge0}$, se obtiene
$(I*K_x)(p)=\sum_{q_2}w(q_2)\sum_{q_1}q_1\ind{p_1+q_1\ge0}
=4\bigl(\ind{p_1\ge-1}-\ind{p_1\ge1}\bigr)=4\cdot\ind{-1\le p_1\le0}$: la
respuesta es $4$ en las dos columnas contiguas al borde y $0$ en el resto.

\begin{proposicion}[Equivariancia a traslaciones]\label{prop:equiv}
Para $a\in\Z^2$, sea $(\tau_aI)(p)=I(p-a)$. Entonces
$(\tau_aI)*K=\tau_a(I*K)$. Si $\rho:\R\to\R$ cumple $\rho(0)=0$, la
cantidad $F(I)=\sum_{p\in\Z^2}\rho\bigl((I*K)(p)\bigr)$ es una suma finita
y $F(\tau_aI)=F(I)$.
\end{proposicion}

\begin{proof}
$((\tau_aI)*K)(p)=\sum_qI(p+q-a)K(q)=(I*K)(p-a)=(\tau_a(I*K))(p)$. La
función $I*K$ es nula fuera de un conjunto finito (el soporte de $I$ más
el de $K$), y $\rho(0)=0$, así que la suma que define $F$ tiene finitos
términos no nulos. Por la primera parte,
$F(\tau_aI)=\sum_p\rho((I*K)(p-a))$, que es la misma suma reindexada con
$p'=p-a$.
\end{proof}

Un detector de bordes responde igual a un borde en cualquier posición y
se desplaza con él; agregando sus respuestas se obtienen características
que no cambian al desplazar el contenido. En una imagen recortada a
$H\times W$ la igualdad vale mientras el desplazamiento no saque contenido
del recuadro. Los histogramas son el caso extremo de agregación:

\begin{ejemplo}[Histograma y su límite]
El histograma $h_I(v)=\#\{p:I(p)=v\}$, $v\in\{0,\dots,255\}$, de una imagen
de $H\times W$ píxeles no cambia si los píxeles se permutan, porque solo
cuenta valores; en particular, no cambia con traslaciones cíclicas ni con
rotaciones de $90^\circ$. Por la misma razón pierde toda la disposición
espacial: un tablero de $8\times8$ casillas blancas y negras y una imagen
con la mitad izquierda negra y la derecha blanca, del mismo tamaño, tienen
el mismo histograma.
\end{ejemplo}

\subsection{Vectores de características de una red preentrenada}

Una red convolucional (CNN) compone capas de convolución con núcleos
aprendidos (equivariantes por la Proposición~\ref{prop:equiv}), funciones
no lineales y agregaciones espaciales, y termina en una capa lineal
seguida de \emph{softmax}. Si $\psi:\R^{H\times W\times C}\to\R^D$ denota
la salida de la penúltima capa, la red calcula probabilidades de clase
$\operatorname{softmax}(A\psi(I)+c)$ con $A\in\R^{C'\times D}$ y
$c\in\R^{C'}$, que es un modelo logístico multinomial con características
$\psi(I)$; al entrenar la red se ajustan conjuntamente $\psi$ y $(A,c)$.
La \emph{transferencia} consiste en tomar $\psi$ de una red entrenada en un
corpus grande, calcular $\psi(I)$ ($D$ del orden de $10^3$) para las
imágenes de un problema nuevo y entrenar sobre esas columnas un modelo
sencillo. Que la $\psi$ aprendida para unas clases sirva para otras no es
un teorema, sino una observación empírica: Yosinski et al.\ (2014) midieron
que las capas iniciales aprenden detectores genéricos (bordes, texturas) y
que la transferibilidad disminuye en las capas finales y cuanto más
difieren las tareas.

\paragraph{Contraste.} Las características clásicas (gradientes,
histogramas) son interpretables, baratas y sin parámetros aprendidos,
pero no representan objetos. Los vectores de una CNN sí, a costa de
interpretabilidad, cómputo y dependencia del preprocesamiento del
preentrenamiento. Las etiquetas de un servicio de visión son pocas
columnas legibles, limitadas al vocabulario del servicio.

\begin{aws}
\textbf{SageMaker JumpStart} ofrece modelos de \emph{image embedding} (por
ejemplo, la familia \emph{Image Embedding -- TensorFlow}) que devuelven
$\psi(I)$ para cada imagen, desplegables como \emph{endpoint} o en
\emph{batch transform}. \textbf{Amazon Bedrock} ofrece modelos de
\emph{embeddings} multimodales. \textbf{Amazon Rekognition}
(\texttt{DetectLabels}) devuelve etiquetas con su confianza y, para
objetos, \emph{bounding boxes} con \texttt{Left}, \texttt{Top},
\texttt{Width} y \texttt{Height} expresados como fracciones del tamaño de
la imagen; con ellas se construyen columnas como la indicadora de una
etiqueta con confianza mayor que un umbral, el número de instancias o la
fracción de área $\texttt{Width}\cdot\texttt{Height}$. (Las funciones de
\emph{Streaming Events} y de moderación de contenido por lotes de
Rekognition dejaron de aceptar clientes nuevos el 30 de abril de 2026; la
detección de etiquetas no está afectada.) \textbf{Data Wrangler} incluye
transformaciones de imagen (cambio de tamaño, mejora, escala de grises,
rotación, eliminación de imágenes corruptas o duplicadas) y, sobre las
características ya tabulares, normalización y PCA.
\end{aws}

\subsection{Análisis de componentes principales}\label{sec:pca}

El capítulo propone PCA para reducir la dimensión de las características
de imagen y la de \emph{one-hot}. Sea $X\in\R^{n\times p}$ con columnas de
media cero y $S=X^\top X/n$ su matriz de covarianzas muestrales. Por el
teorema espectral, $S=\sum_{i=1}^p\lambda_iu_iu_i^\top$ con
$\lambda_1\ge\dots\ge\lambda_p\ge0$ y $\{u_1,\dots,u_p\}$ una base
ortonormal de autovectores; los autovalores son no negativos porque
$u^\top Su=\|Xu\|^2/n\ge0$.

\begin{teorema}[PCA como máxima varianza]\label{teo:pca}
Para $1\le d\le p$, el máximo de $\tr(V^\top SV)$ sobre las matrices
$V\in\R^{p\times d}$ con $V^\top V=I_d$ es $\lambda_1+\dots+\lambda_d$ y
se alcanza en $V=[u_1\cdots u_d]$. La cantidad $\tr(V^\top SV)$ es la suma
de las varianzas muestrales de las columnas de $XV$, las proyecciones de
los datos sobre las columnas de $V$.
\end{teorema}

\begin{proof}
$\tr(V^\top SV)=\sum_i\lambda_i\tr(V^\top u_iu_i^\top V)=\sum_i\lambda_iw_i$
con $w_i=\|V^\top u_i\|^2$. Por un lado, $VV^\top$ es la proyección
ortogonal sobre el espacio columna de $V$ (porque $V^\top V=I$), así que
$w_i=u_i^\top VV^\top u_i=\|VV^\top u_i\|^2\le\|u_i\|^2=1$. Por otro,
$\sum_iw_i=\tr\bigl(V^\top(\sum_iu_iu_i^\top)V\bigr)=\tr(V^\top V)=d$.
Como $w_i-1\le0$ y $\lambda_i\ge\lambda_d$ para $i\le d$, y $w_i\ge0$ y
$\lambda_i\le\lambda_d$ para $i>d$,
\[
\sum_i\lambda_iw_i-\sum_{i\le d}\lambda_i=\sum_{i\le
d}\lambda_i(w_i-1)+\sum_{i>d}\lambda_iw_i\le\lambda_d\Bigl(\sum_{i\le
d}(w_i-1)+\sum_{i>d}w_i\Bigr)=\lambda_d\Bigl(\sum_iw_i-d\Bigr)=0 .
\]
Con $V=[u_1\cdots u_d]$ se tiene $w_i=1$ para $i\le d$ y $w_i=0$ para
$i>d$, y se alcanza el valor. Finalmente, las columnas de $XV$ tienen
media cero (porque las de $X$ la tienen), su matriz de covarianzas es
$V^\top X^\top XV/n=V^\top SV$, y la traza es la suma de sus varianzas.
\end{proof}

\begin{proposicion}[Error de reconstrucción]\label{prop:reconst}
Con $V=[u_1\cdots u_d]$, el error cuadrático medio de reconstruir cada fila
$x_l$ por su proyección $VV^\top x_l$ es
$n^{-1}\sum_l\|x_l-VV^\top x_l\|^2=\sum_{i>d}\lambda_i$.
\end{proposicion}

\begin{proof}
Como $VV^\top$ es una proyección ortogonal, $x-VV^\top x$ es ortogonal a
$VV^\top x$ y, por el teorema de Pitágoras,
$\|x-VV^\top x\|^2=\|x\|^2-\|VV^\top x\|^2=\|x\|^2-\|V^\top x\|^2$ (la última
igualdad porque $\|Vv\|=\|v\|$ cuando $V^\top V=I$). Promediando sobre las filas,
$n^{-1}\sum_l\|x_l\|^2=\tr(S)=\sum_i\lambda_i$ y
$n^{-1}\sum_l\|V^\top x_l\|^2=\tr(V^\top SV)=\sum_{i\le d}\lambda_i$.
\end{proof}

La \emph{fracción de varianza explicada} por $d$ componentes es
$\sum_{i\le d}\lambda_i/\sum_{i\le p}\lambda_i$, y lo que falta es, por la
proposición, el error relativo de reconstrucción. En la práctica se usa la
descomposición en valores singulares $X=U\Sigma W^\top$ ($U$ y $W$ con
columnas ortonormales, $\Sigma$ diagonal con valores singulares
$\sigma_i\ge0$): entonces $S=W(\Sigma^2/n)W^\top$, de modo que las
columnas de $W$ son autovectores de $S$ con autovalores $\sigma_i^2/n$, y
no hace falta formar $S$.

\begin{ejemplo}
Si $S=\begin{psmallmatrix}2&1\\1&2\end{psmallmatrix}$, los autovalores son
$3$ y $1$, con autovectores $(1,1)/\sqrt2$ y $(1,-1)/\sqrt2$; la primera
componente, la suma normalizada de las dos variables, explica $3/4$ de la
varianza total.
\end{ejemplo}

\begin{observacion}[PCA depende de las unidades]
Si una columna se multiplica por $c$, sus entradas en $S$ se multiplican
por $c$ (covarianzas) o $c^2$ (varianza); medir una longitud en centímetros
en lugar de metros multiplica su varianza por $10^4$ y la primera
componente tiende a alinearse con ella. Estandarizar antes equivale a
hacer PCA sobre la matriz de correlaciones. Tras \emph{one-hot}, la
indicadora de un nivel de frecuencia relativa $f_i$ tiene varianza
$f_i(1-f_i)$ (varianza de una Bernoulli), pequeña para niveles raros; PCA,
que maximiza varianza y no usa $y$, tiende a descartarlos aunque sean
predictivos.
\end{observacion}

\begin{aws}
\textbf{Data Wrangler}, \emph{Reduce dimensionality} (PCA): número de
componentes fijo o umbral de fracción de varianza explicada (95\,\% por
defecto), opciones de centrar y de escalar a desviación $1$ (la
observación anterior indica cuándo usarlas), salida como vector o columnas
y opción de conservar las columnas originales. \textbf{Algoritmo
integrado PCA} de SageMaker: modo \emph{regular} (datos dispersos, tamaño
moderado) o \emph{randomized} (muchas filas y columnas, con un algoritmo
aproximado); calcula la matriz de covarianzas, o una aproximación
distribuida, y su descomposición en valores singulares; acepta
\texttt{recordIO-protobuf} y CSV.
\end{aws}

% =====================================================================
\section{Texto}\label{sec:texto}
% =====================================================================

Un \emph{corpus} es una colección de documentos $d_1,\dots,d_N$, cada uno
una cadena de caracteres. El objetivo es representar cada documento (o
cada palabra) por un vector. La secuencia habitual es: tokenizar,
normalizar (palabras vacías, \emph{stemming} o lematización) y representar
(bolsa de palabras, $n$-gramas, TF-IDF o \emph{embeddings}).

\subsection{Tokenización}

\begin{definicion}
Un \emph{tokenizador} es una función que asigna a cada cadena de
caracteres una sucesión finita $(t_1,\dots,t_L)$ de elementos de un
conjunto $\mathcal V$, el \emph{vocabulario}; los $t_\ell$ son los
\emph{tokens} y $L$ la longitud del documento en tokens. En la
tokenización por palabras, los tokens son las subcadenas delimitadas por
espacios y signos de puntuación, y $\mathcal V$ es el conjunto de palabras
vistas en entrenamiento.
\end{definicion}

Por ejemplo, «Machine learning is fascinating» da
(Machine, learning, is, fascinating). Con tokenización por palabras, una
palabra ausente de $\mathcal V$ no tiene representación y se sustituye por
un símbolo genérico, con lo que su información se pierde. Los modelos
fundacionales usan tokenización por subpalabras, que elimina ese problema.

\begin{definicion}[\emph{Byte-pair encoding}, Sennrich et al., 2016]
Sea $\mathcal V_0$ el conjunto de caracteres (o de bytes). Se segmenta cada
palabra del corpus de entrenamiento en caracteres y se repite $M$ veces:
contar, en la segmentación actual, la frecuencia de cada par de símbolos
adyacentes; agregar al vocabulario el símbolo formado por la concatenación
del par más frecuente (con una regla fija de desempate); reemplazar en la
segmentación todas sus apariciones. El vocabulario final tiene
$|\mathcal V_0|+M$ elementos, y la lista ordenada de fusiones define el
tokenizador: una cadena nueva se segmenta en caracteres y se le aplican
las fusiones en el mismo orden.
\end{definicion}

\begin{proposicion}
Si $\mathcal V_0$ contiene todos los bytes, toda cadena se tokeniza con
símbolos de $\mathcal V$.
\end{proposicion}

\begin{proof}
La segmentación inicial usa solo elementos de $\mathcal V_0\subset\mathcal
V$, y cada fusión reemplaza dos símbolos consecutivos por su
concatenación, que es un elemento de $\mathcal V$ por construcción; por
inducción sobre el número de fusiones aplicadas, todos los símbolos de la
segmentación están en $\mathcal V$.
\end{proof}

Las palabras frecuentes terminan como un solo token y las raras se
dividen en fragmentos. El número de tokens por palabra depende del idioma
y del tokenizador, y determina el costo en servicios de modelos
fundacionales de texto facturados por uso: $p_{\mathrm{in}}T_{\mathrm{in}}+p_{\mathrm{out}}T_{\mathrm{out}}$,
con $T_{\mathrm{in}},T_{\mathrm{out}}$ los números de tokens de entrada y
salida y $p_{\mathrm{in}},p_{\mathrm{out}}$ sus precios unitarios.

\begin{aws}
\textbf{Glue DataBrew}, acción \texttt{TOKENIZATION}, que en un paso
tokeniza con un delimitador, expande contracciones, aplica
\emph{stemming} y elimina palabras vacías:
\begin{lstlisting}
{"Action": {"Operation": "TOKENIZATION",
  "Parameters": {"sourceColumn": "resena", "targetColumn": "resena_tok",
                 "delimiter": " ", "expandContractions": "ENABLED",
                 "stemmingMode": "PORTER", "stopWordRemovalMode": "DEFAULT",
                 "customStopWords": "[]"}}}
\end{lstlisting}
\textbf{Data Wrangler}, \emph{Featurize text} $\to$ \emph{Vectorize}, con
tokenizador \emph{Standard} (separación por espacios) o \emph{Custom}
(expresión regular). En \textbf{Amazon Bedrock} la tokenización por
subpalabras es interna a cada modelo y es la unidad de facturación de los
modelos de texto bajo demanda.
\end{aws}

\subsection{Palabras vacías}

La \emph{ley de Zipf} es una regularidad empírica: en muchos corpus, la
frecuencia de la palabra que ocupa el lugar $r$ en orden de frecuencia es
aproximadamente proporcional a $r^{-s}$, con $s$ cercano a $1$. Bajo este
modelo con $s=1$ y $V$ palabras distintas, la fracción de los tokens que
corresponde a las $k$ más frecuentes es $H_k/H_V$, con
$H_k=\sum_{r=1}^k1/r$: la frecuencia de la palabra $r$ es $r^{-1}/H_V$ del
total. Con $V=50\,000$ y $k=100$, $H_{100}/H_{50\,000}\approx5.19/11.40\approx0.46$.
Las palabras más frecuentes son en su mayoría funcionales (artículos,
preposiciones, conjunciones), y se llaman \emph{palabras vacías}.
Eliminarlas reduce casi a la mitad el volumen de tokens y el peso de términos
que aparecen en todos los documentos, pero no es neutral: las listas
estándar incluyen negaciones, y eliminar «no» convierte «no bueno» en
«bueno». La ponderación TF-IDF (\S\ref{sec:tfidf}) hace una atenuación
basada en los datos.

\begin{aws}
\textbf{DataBrew}: \texttt{stopWordRemovalMode} (lista por defecto) y
\texttt{customStopWords} dentro de \texttt{TOKENIZATION}.
\textbf{Data Wrangler}: el \emph{Count vectorizer} admite una frecuencia
documental máxima, que elimina los términos presentes en demasiados
documentos: una lista de palabras vacías construida con los datos.
\end{aws}

\subsection{\emph{Stemming} y lematización}

\begin{definicion}
Una \emph{normalización} es una función $\psi:\mathcal V\to\mathcal V'$ que
asigna a cada palabra una forma representativa; las palabras con la misma
imagen se consideran la misma. El \emph{stemming} define $\psi$ con reglas
de supresión de sufijos (algoritmos de Porter, 1980, y de Lancaster); la
lematización, con un diccionario morfológico que usa la categoría
gramatical de la palabra.
\end{definicion}

\begin{proposicion}[Normalizar es restringir coeficientes]\label{prop:stem}
Sea $x\in\R^{\mathcal V}$ el vector de conteos de un documento ($x_v$ es el
número de apariciones de la palabra $v$) y $x'=Mx\in\R^{\mathcal V'}$ el
vector tras normalizar, con $M_{v'v}=\ind{\psi(v)=v'}$. Los predictores
lineales $\gamma^\top x'$ son exactamente los predictores lineales
$\beta^\top x$ cuyos coeficientes son iguales dentro de cada clase
$\psi^{-1}(v')$.
\end{proposicion}

\begin{proof}
$\gamma^\top x'=\gamma^\top Mx=(M^\top\gamma)^\top x$ y
$(M^\top\gamma)_v=\sum_{v'}M_{v'v}\gamma_{v'}=\gamma_{\psi(v)}$, constante en
cada clase. Recíprocamente, si $\beta_u=\beta_v$ siempre que
$\psi(u)=\psi(v)$, se toma $\gamma_{v'}$ igual al valor común de $\beta$
en $\psi^{-1}(v')$ (cualquiera si la clase es vacía) y se tiene
$\beta=M^\top\gamma$.
\end{proof}

Normalizar reduce el número de parámetros y agrega los conteos de formas
de la misma palabra (menos varianza), a cambio de sesgo cuando se fusionan
palabras con efectos distintos. Ejemplos reales (implementaciones de
NLTK): Porter lleva «studies» y «studying» a «studi», pero también
«university» y «universe» a «univers» y «organization» y «organ» a
«organ»; Lancaster, más agresivo, lleva «generous» y «generation» a
«gen» y «better» a «bet». La lematización lleva «studies» a «study» y,
sabiendo que es un adjetivo, «better» a «good», lo que ninguna regla de
supresión de sufijos logra. Porter y Lancaster son algoritmos para inglés;
para español se usa el \emph{stemmer} Snowball de español o un
lematizador.

\begin{aws}
\textbf{DataBrew}: \texttt{stemmingMode} igual a \texttt{PORTER} o
\texttt{LANCASTER} en \texttt{TOKENIZATION}. No hay lematización integrada
en DataBrew ni en Data Wrangler; se implementa en \textbf{SageMaker
Processing} con spaCy o NLTK instalados en el contenedor. \textbf{Amazon
Comprehend} (\texttt{DetectSyntax}) devuelve los tokens con su categoría
gramatical, que es el dato que necesita un lematizador.
\end{aws}

\subsection{Bolsa de palabras, $n$-gramas y TF-IDF}\label{sec:tfidf}

\begin{definicion}
Fijado un vocabulario $\mathcal V$ de \emph{términos}, la \emph{matriz
documento-término} $X\in\R^{N\times|\mathcal V|}$ tiene entradas
$X_{dt}=\tf_{t,d}$, el número de apariciones del término $t$ en el
documento $d$; la fila $d$ es la \emph{bolsa de palabras} del documento.
Un $n$-grama es una sucesión de $n$ tokens consecutivos; en la
representación por $n$-gramas, los términos son los $n$-gramas observados
en entrenamiento. Sea $\df_t$ el número de documentos de entrenamiento que
contienen $t$. Las ponderaciones TF-IDF de Spark ML y de scikit-learn
(con \texttt{smooth\_idf=True}) son $\tf_{t,d}\cdot\idf_t$ con
\[
\idf_t^{\text{Spark}}=\log\frac{N+1}{\df_t+1},\qquad
\idf_t^{\text{sklearn}}=\log\frac{N+1}{\df_t+1}+1,
\]
habitualmente seguidas de normalización euclídea de cada fila.
\end{definicion}

La bolsa de palabras solo cuenta, así que no cambia si se permutan los
tokens de un documento: «el perro muerde al hombre» y «el hombre muerde al
perro» tienen la misma fila. Con bigramas difieren («perro muerde» frente
a «hombre muerde»), y «no bueno» queda como un término. El costo es
combinatorio: hay $|\mathcal V|^n$ $n$-gramas posibles y un documento de
$L$ tokens contiene a lo sumo $L-n+1$, así que la matriz es más dispersa
al crecer $n$.

La ponderación IDF (Spärck Jones, 1972) responde a que un término presente
en muchos documentos discrimina poco entre ellos: $\log(N/\df_t)$ es la
cantidad de información, $-\log p$, del suceso «el término aparece en un
documento elegido al azar», cuya frecuencia es $p=\df_t/N$. Los $+1$
evitan divisiones por cero; el $+1$ final de scikit-learn impide que un
término presente en todos los documentos reciba peso nulo.

\begin{ejemplo}
Documentos «el gato come», «el perro come carne» y «el gato duerme»
($N=3$). Con la fórmula de Spark: $\idf(\text{el})=\log(4/4)=0$,
$\idf(\text{gato})=\idf(\text{come})=\log(4/3)\approx0.288$ y
$\idf(\text{perro})=\idf(\text{carne})=\idf(\text{duerme})=\log2\approx0.693$.
Con la de scikit-learn, los mismos valores más $1$: «el» conserva peso $1$.
\end{ejemplo}

\begin{proposicion}\label{prop:coseno}
Si las filas TF-IDF se normalizan a norma euclídea $1$, el producto
interno de dos filas normalizadas es el coseno del ángulo entre las filas
sin normalizar. Si todos los conteos de un documento se multiplican por
$c>0$, su fila normalizada no cambia.
\end{proposicion}

\begin{proof}
$\langle u/\|u\|,v/\|v\|\rangle=\langle u,v\rangle/(\|u\|\|v\|)$, que es
la definición del coseno. Si $\tf_{t,d}\mapsto c\,\tf_{t,d}$, la fila
TF-IDF se multiplica por $c$ (los $\idf$ son los de entrenamiento y no
cambian), y su normalización es la misma.
\end{proof}

Por ejemplo, un documento y su concatenación consigo mismo quedan
representados igual: la normalización elimina el efecto de la longitud.

\begin{aws}
\textbf{Data Wrangler}, \emph{Featurize text} $\to$ \emph{Vectorize}:
tokenizador; vectorizador \emph{Count vectorizer} (frecuencias mínima y
máxima de término y de documento, tamaño máximo de vocabulario, salida
binaria) o \emph{Hashing} (\S\ref{sec:hashing}, caso (b)); opción
\emph{Apply IDF} con frecuencia documental mínima ($5$ por defecto); salida
como vector disperso o aplanada. \emph{Character statistics} añade el
número de palabras y de caracteres y las proporciones de mayúsculas,
minúsculas, dígitos y otros símbolos. Es un paso con parámetros
ajustados (vocabulario, IDF), con \emph{refit}. \textbf{SageMaker
Processing}: \texttt{TfidfVectorizer(ngram\_range=(1,2), min\_df=...,
max\_df=...)}; si después se escala, \texttt{MaxAbsScaler} conserva la
dispersión (Proposición~\ref{prop:sparse}).
\end{aws}

\subsection{\emph{Word embeddings}}

En la bolsa de palabras cada término es una coordenada, lo que equivale a
representar la palabra $v$ por $e_v$; para $u\ne v$,
$\cos(e_u,e_v)=0$. «Gato» y «felino» son tan distintos como «gato» e
«hipoteca», y lo aprendido sobre una palabra no se transfiere a otra. La
hipótesis distribucional (Harris, 1954), según la cual palabras que
aparecen en contextos parecidos tienen significados parecidos, sugiere
representar cada palabra por un vector denso de dimensión $D$ (del orden de
$100$ a $300$) aprendido de las coocurrencias.

\paragraph{Pares palabra--contexto.} Se ve el corpus tokenizado como una
sucesión $w_1,\dots,w_{N_T}$ de tokens. Con un radio $L\ge1$, el
multiconjunto de pares es $\mathcal P=\{(w_a,w_b):1\le a,b\le N_T,\
0<|a-b|\le L\}$: cada token $w_a$ (la \emph{palabra}) con cada token a
distancia a lo sumo $L$ (su \emph{contexto}). Sean $\#(w,c)$ el número de
veces que el par $(w,c)$ aparece en $\mathcal P$, $\#(w)=\sum_c\#(w,c)$,
$\#(c)=\sum_w\#(w,c)$ y $|\mathcal P|$ el total de pares; las frecuencias
relativas $\hat P(w,c)=\#(w,c)/|\mathcal P|$,
$\hat P(w)=\#(w)/|\mathcal P|$ y $\hat P(c)=\#(c)/|\mathcal P|$ definen la
\emph{información mutua puntual}
$\PMI(w,c)=\log\frac{\hat P(w,c)}{\hat P(w)\hat P(c)}$.

\begin{definicion}[\emph{Skip-gram} con muestreo negativo, Mikolov et al., 2013]
Se asigna a cada palabra un vector $v_w\in\R^D$ y a cada contexto un
vector $u_c\in\R^D$. Con $\sigma(x)=1/(1+e^{-x})$, un entero $k\ge1$ y una
distribución $P_n$ sobre contextos (la distribución de ruido), se maximiza
\[
\ell=\sum_{w,c}\#(w,c)\log\sigma(u_c^\top v_w)+k\sum_w\#(w)\sum_cP_n(c)\log\sigma(-u_c^\top v_w).
\]
\end{definicion}

El primer término premia que los pares observados tengan producto interno
grande; el segundo es la esperanza, para cada aparición de $w$, de $k$
contextos falsos sorteados con $P_n$, y premia que tengan producto interno
pequeño. El problema que resolvió esta formulación es de cómputo: el
modelo original,
$p(c\mid w)=\exp(u_c^\top v_w)/\sum_{c'}\exp(u_{c'}^\top v_w)$, suma sobre
todo el vocabulario en cada evaluación, y $\ell$ solo evalúa $k+1$
términos por par. Lo que se aprende queda caracterizado así:

\begin{teorema}[Levy y Goldberg, 2014]\label{teo:sgns}
Supóngase que los números $x_{wc}=u_c^\top v_w$ pueden elegirse
libremente (lo que requiere $D$ suficientemente grande) y que
$\#(w,c)>0$ y $P_n(c)>0$ para todos los pares $(w,c)$ que intervienen en
$\ell$. Entonces el máximo de $\ell$ se alcanza en
\[
u_c^\top v_w=\log\frac{\#(w,c)}{\#(w)\,P_n(c)}-\log k .
\]
Si $P_n(c)=\hat P(c)$, el valor es $\PMI(w,c)-\log k$.
\end{teorema}

\begin{proof}
Con los $x_{wc}$ libres, $\ell=\sum_{w,c}\ell_{wc}(x_{wc})$ con
$\ell_{wc}(x)=\#(w,c)\log\sigma(x)+k\#(w)P_n(c)\log\sigma(-x)$, y se
maximiza cada sumando por separado. Como $\frac{d}{dx}\log\sigma(x)=\sigma(-x)$
y $\frac{d}{dx}\log\sigma(-x)=-\sigma(x)$, se tiene
$\ell_{wc}'(x)=\#(w,c)\sigma(-x)-k\#(w)P_n(c)\sigma(x)$ y
$\ell_{wc}''(x)=-\bigl(\#(w,c)+k\#(w)P_n(c)\bigr)\sigma(x)\sigma(-x)<0$;
$\ell_{wc}$ es estrictamente cóncava y su máximo es su único punto
crítico. Usando $\sigma(-x)/\sigma(x)=e^{-x}$, la condición $\ell_{wc}'(x)=0$
es $\#(w,c)e^{-x}=k\#(w)P_n(c)$, es decir,
$x=\log\frac{\#(w,c)}{\#(w)P_n(c)}-\log k$. Con $P_n(c)=\#(c)/|\mathcal
P|$, $\frac{\#(w,c)}{\#(w)\#(c)/|\mathcal P|}=\frac{\hat P(w,c)}{\hat
P(w)\hat P(c)}$.
\end{proof}

Con $D$ finito, los productos internos no pueden elegirse libremente y el
método ajusta una aproximación de rango $D$ a la matriz de PMI desplazada;
las palabras con filas de PMI parecidas reciben vectores parecidos, que es
el sentido en que el método generaliza. En la práctica $P_n$ es la
distribución unigrama elevada a $3/4$ y renormalizada, y el teorema vale
con esa $P_n$.

\begin{definicion}[GloVe, Pennington et al., 2014]
Con $X_{ij}$ el número de coocurrencias de las palabras $i$ y $j$ en
ventanas del corpus, se minimiza
$J=\sum_{i,j:X_{ij}>0}f(X_{ij})\bigl(v_i^\top\tilde v_j+b_i+\tilde
b_j-\log X_{ij}\bigr)^2$ sobre vectores $v_i,\tilde v_j\in\R^D$ y sesgos
$b_i,\tilde b_j\in\R$, con $f(x)=(x/x_{\max})^{3/4}$ si $x<x_{\max}$ y
$f(x)=1$ si no ($x_{\max}=100$).
\end{definicion}

GloVe hace explícita una factorización del mismo tipo. Si existen
vectores con $v_i^\top\tilde v_j=\PMI(i,j)$ para todos los pares con
$X_{ij}>0$ (con $\PMI$ calculada a partir de $X$), entonces con
$b_i=\log X_i$ y $\tilde b_j=\log X_j-\log X_{\cdot\cdot}$, donde
$X_i=\sum_jX_{ij}$ y $X_{\cdot\cdot}=\sum_{ij}X_{ij}$, cada término vale
$\PMI(i,j)+\log X_i+\log X_j-\log X_{\cdot\cdot}-\log X_{ij}=0$ por la
definición de PMI, y $J=0$ es el mínimo: los sesgos absorben las
frecuencias marginales y los productos internos representan la PMI. La
función $f$ reduce el peso de los pares rarísimos, cuyos conteos son
ruidosos, y acota el de los frecuentísimos.

\begin{proposicion}[Promedio de \emph{embeddings}]\label{prop:promedio}
Sea $E\in\R^{|\mathcal V|\times D}$ la matriz cuyas filas son los vectores
$v_w$ y $x$ la bolsa de palabras de un documento con $L$ tokens. El
promedio de los vectores de sus tokens es $L^{-1}E^\top x$.
\end{proposicion}

\begin{proof}
$\sum_{\ell=1}^Lv_{t_\ell}=\sum_wx_wv_w=E^\top x$, agrupando los tokens
iguales, pues $x_w$ cuenta las apariciones de $w$.
\end{proof}

Representar un documento por el promedio de sus \emph{embeddings} es, por
tanto, proyectar su bolsa de palabras (dividida por $L$) con una matriz
fija aprendida en otro corpus, y hereda la invariancia al orden de la
bolsa de palabras. El análisis semántico latente (LSA), en cambio,
proyecta sobre las primeras componentes principales de la matriz TF-IDF del
propio corpus (Teorema~\ref{teo:pca}). Los \emph{embeddings} contextuales
de los modelos basados en transformadores asignan a cada aparición de una
palabra un vector que depende de toda la oración, y pueden distinguir
sentidos de una misma palabra («banco» para sentarse y «banco» que presta
dinero), lo que un vector por palabra no puede hacer por definición.

\paragraph{Contraste.} Bolsa de palabras y TF-IDF: interpretables, sin
preentrenamiento, dispersas y de alta dimensión, sin similitud entre
palabras ni orden. $n$-gramas: orden local, a costa de dispersión.
Word2Vec y GloVe: densos, con similitud entre palabras y transferibles; un
vector por palabra, y sin representación para palabras nuevas salvo que se
modelen subpalabras. \emph{Embeddings} contextuales: resuelven la
polisemia; son costosos de calcular y opacos.

\begin{aws}
\textbf{BlazingText}, algoritmo integrado de SageMaker, implementa
Word2Vec en modos \texttt{cbow}, \texttt{skipgram} (una CPU o una GPU) y
\texttt{batch\_skipgram} (varias CPU), con vectores de subpalabras
(\emph{character $n$-grams}) que dan vectores a palabras fuera del
vocabulario. Espera un archivo de texto con una oración por línea y tokens
separados por espacios, y produce \texttt{vectors.txt} (compatible con
Gensim y spaCy) y \texttt{vectors.bin}. \textbf{Amazon Bedrock},
\emph{Titan Text Embeddings V2} (\texttt{amazon.titan-embed-text-v2:0}):
\emph{embeddings} contextuales de documento, con dimensión configurable
($1024$, $512$ o $256$) y normalización opcional, para entradas de hasta
unos $8000$ tokens. \textbf{Amazon Comprehend}: características de alto
nivel por API (sentimiento con puntuaciones por clase, entidades, frases
clave, idioma, sintaxis) que se agregan como columnas; su modelado de
tópicos dejó de aceptar clientes nuevos el 30 de abril de 2026, y para
tópicos quedan los algoritmos integrados \textbf{LDA} y \textbf{NTM} de
SageMaker. \textbf{Amazon Textract} convierte documentos escaneados en
texto antes de todo lo anterior.
\end{aws}

% =====================================================================
\section{Selección de características por filtros}\label{sec:seleccion}
% =====================================================================

Un \emph{método de filtro} ordena las características según una medida de
asociación con $y$, calculada en entrenamiento, y conserva las primeras.
Las medidas difieren en qué dependencia detectan.

\begin{contraejemplo}[Correlación nula no es independencia]
Si $X\sim N(0,1)$ e $Y=X^2$, entonces $\Cov(X,Y)=\E[X^3]-\E[X]\E[X^2]=0$
porque $X$ es simétrica, y $\Corr(X,Y)=0$, aunque $Y$ es función de $X$.
Un filtro por correlación de Pearson descartaría $X$.
\end{contraejemplo}

\begin{definicion}
Para $X,Y$ discretas con distribución conjunta $P_{XY}$ y marginales
$P_X,P_Y$, la \emph{información mutua} es
$I(X;Y)=\sum_{x,y:P_{XY}(x,y)>0}P_{XY}(x,y)\log\frac{P_{XY}(x,y)}{P_X(x)P_Y(y)}$,
es decir, $\mathrm{KL}(P_{XY}\|P_XP_Y)$.
\end{definicion}

\begin{proposicion}
$I(X;Y)\ge0$, con igualdad si y solo si $X$ e $Y$ son independientes.
\end{proposicion}

\begin{proof}
Por la desigualdad de Jensen para $\log$ (cóncava) y la distribución
$P_{XY}$,
$-I(X;Y)=\sum P_{XY}\log\frac{P_XP_Y}{P_{XY}}\le\log\sum_{P_{XY}>0}P_XP_Y\le\log1=0$,
con igualdad en la primera desigualdad si y solo si
$P_XP_Y/P_{XY}$ es constante donde $P_{XY}>0$, y en la segunda si y solo
si $P_XP_Y$ se anula donde $P_{XY}$ se anula; ambas juntas equivalen a
$P_{XY}=P_XP_Y$.
\end{proof}

La información mutua detecta cualquier dependencia, incluida la del
contraejemplo tras discretizar (para variables continuas, scikit-learn la
estima con vecinos más cercanos). Para dos categóricas está ligada a la
prueba de independencia:

\begin{proposicion}
En una tabla de contingencia con conteos $O_{ij}$, total $N$ y conteos
esperados bajo independencia $E_{ij}=O_{i\cdot}O_{\cdot j}/N$ (con
$O_{i\cdot},O_{\cdot j}$ los totales de fila y columna), el estadístico de
razón de verosimilitudes $G=2\sum_{ij:O_{ij}>0}O_{ij}\log(O_{ij}/E_{ij})$
vale $2N\hat I$, donde $\hat I$ es la información mutua de la distribución
empírica $\hat p_{ij}=O_{ij}/N$.
\end{proposicion}

\begin{proof}
$O_{ij}/E_{ij}=\hat p_{ij}/(\hat p_{i\cdot}\hat p_{\cdot j})$ con
$\hat p_{i\cdot}=O_{i\cdot}/N$ y $\hat p_{\cdot j}=O_{\cdot j}/N$; entonces
\[
\sum_{ij}O_{ij}\log\frac{O_{ij}}{E_{ij}}=N\sum_{ij}\hat p_{ij}\log\frac{\hat
p_{ij}}{\hat p_{i\cdot}\hat p_{\cdot j}}=N\hat I .
\]
\end{proof}

El estadístico de Pearson $\chi^2=\sum_{ij}(O_{ij}-E_{ij})^2/E_{ij}$
aproxima a $G$ cuando los $O_{ij}$ son cercanos a los $E_{ij}$: con
$u=(O_{ij}-E_{ij})/E_{ij}$, el desarrollo
$(1+u)\log(1+u)=u+u^2/2+O(u^3)$ da
$O_{ij}\log(O_{ij}/E_{ij})=(O_{ij}-E_{ij})+\frac{(O_{ij}-E_{ij})^2}{2E_{ij}}+O\bigl(|O_{ij}-E_{ij}|^3/E_{ij}^2\bigr)$,
y al sumar, como $\sum_{ij}(O_{ij}-E_{ij})=0$, queda
$G=\chi^2+\sum_{ij}O\bigl(|O_{ij}-E_{ij}|^3/E_{ij}^2\bigr)$.

\begin{definicion}[Factor de inflación de la varianza]
En una matriz de diseño con intercepto, sea $R_j^2$ el coeficiente de
determinación de la regresión de la columna $j$ sobre las demás columnas.
$\mathrm{VIF}_j=1/(1-R_j^2)$.
\end{definicion}

\begin{proposicion}
En el modelo $y=Z\beta+\varepsilon$ con $Z$ de rango completo por
columnas, que incluye la columna $\1$, y $\Var\varepsilon=\sigma^2I$, el
estimador de mínimos cuadrados cumple
$\Var(\hat\beta_j)=\sigma^2\,\mathrm{VIF}_j/\mathrm{SCT}_j$, con
$\mathrm{SCT}_j=\sum_k(z_{kj}-\bar z_j)^2$, para cada columna $j$ distinta
del intercepto.
\end{proposicion}

\begin{proof}
Por el teorema de Frisch--Waugh--Lovell (que requiere solo que $Z$ tenga
rango completo, lo que se supone), $\hat\beta_j=r_j^\top y/r_j^\top r_j$,
donde $r_j$ es el vector de residuos de la regresión de la columna $j$
sobre las demás. Como $r_j$ no depende de $y$,
$\Var(\hat\beta_j)=r_j^\top(\sigma^2I)r_j/(r_j^\top r_j)^2=\sigma^2/(r_j^\top
r_j)$. Como las demás columnas incluyen $\1$, la definición del
coeficiente de determinación da $r_j^\top r_j=(1-R_j^2)\mathrm{SCT}_j$.
\end{proof}

El VIF mide redundancia, no relevancia: una característica con VIF alto
puede ser muy predictiva, pero su coeficiente es inestable. Es la
formalización del ejemplo del capítulo de ingresos y ventas que «se mueven
en paralelo».

\begin{aws}
\textbf{Data Wrangler} ofrece análisis (no transformaciones) para
seleccionar: \emph{Quick Model} (bosque aleatorio con importancias de
Gini y partición 70/30); \emph{Target Leakage} (para cada columna, AUC o
$R^2$ de un modelo con esa sola columna por validación cruzada sobre una
muestra; valores cercanos a $1$ indican posible fuga, como la de los
ingresos del mismo día que señala el capítulo); \emph{Multicollinearity}
(VIF recortado a 50, valores singulares de PCA tras estandarizar y
coeficientes de lasso); y el \emph{Data Quality and Insights Report}, con
el poder predictivo de cada característica. \textbf{SageMaker Processing}:
\texttt{SelectKBest} con \texttt{mutual\_info\_classif}, \texttt{chi2} o
\texttt{f\_classif}, ajustado solo con entrenamiento.
\end{aws}

% =====================================================================
\section{Resumen: técnica y servicio}
% =====================================================================

\begin{center}
\footnotesize
\begin{tabularx}{\textwidth}{@{}>{\raggedright\arraybackslash}p{2.9cm}YYY@{}}
\toprule
Técnica & Data Wrangler (Canvas) & Glue DataBrew & Otra vía en AWS\\
\midrule
Imputación, indicadora & \emph{Handle missing} & \texttt{FILL\_WITH\_*} & Processing (\texttt{SimpleImputer})\\
Atípicos ($z$, IQR, MAD) & \emph{Handle outliers} & \texttt{FLAG\_OUTLIERS}, \texttt{REMOVE\_OUTLIERS}, etc. & ---\\
Escalado lineal & \emph{Process numeric} & \texttt{SCALE} & Processing\\
$\log$, raíces & \emph{Custom formula} & \texttt{SKEWNESS} & Processing\\
Box--Cox, Yeo--Johnson, cuantil & \emph{Custom transform} & --- & Processing (\texttt{PowerTransformer}, \texttt{QuantileTransformer})\\
Discretización & \emph{Custom formula}/\emph{transform} & \texttt{BUCKETIZATION}, \texttt{BINARIZATION} & Processing (\texttt{KBinsDiscretizer})\\
Ordinal & \emph{Ordinal encode} & \texttt{CATEGORICAL\_MAPPING} & Processing\\
\emph{One-hot} & \emph{One-hot encode} & \texttt{ONE\_HOT\_ENCODING} ($\le10$ niveles) & Processing\\
Binaria & \emph{Custom formula} (bits) & --- & Processing (\texttt{category\_encoders})\\
\emph{Hashing} & \emph{Vectorize}/\emph{Custom transform} & --- & Processing (\texttt{FeatureHasher})\\
Similitud (\emph{min-hash}) & \emph{Similarity encode} & --- & ---\\
\emph{Target encoding} & --- & --- & Processing, CatBoost integrado, Glue + Feature Store\\
WoE / IV & --- & (\texttt{BUCKETIZATION} previa) & Processing o Glue; Linear Learner\\
Fecha cíclica & \emph{Featurize datetime} & --- & Processing\\
Rezagos, ventanas & \emph{Lag}/\emph{Rolling window features} & --- & Processing\\
Imagen: vectores CNN & transformaciones de imagen & --- & JumpStart, Bedrock, Rekognition\\
PCA & \emph{Reduce dimensionality} & --- & PCA integrado\\
Tokenización, \emph{stemming} & \emph{Vectorize} (tokenizador) & \texttt{TOKENIZATION} & Processing (spaCy, NLTK)\\
TF-IDF, $n$-gramas & \emph{Vectorize} + IDF & --- & Processing (\texttt{TfidfVectorizer})\\
\emph{Word embeddings} & --- & --- & BlazingText, Bedrock (Titan), Comprehend\\
Selección & análisis \emph{Quick Model}, \emph{Target Leakage}, \emph{Multicollinearity} & --- & Processing (\texttt{SelectKBest})\\
\bottomrule
\end{tabularx}
\end{center}

En todos los casos, los parámetros ajustados se estiman con entrenamiento
(Proposición~\ref{prop:test}), se guardan como artefacto versionado y se
aplican sin cambios en validación, prueba e inferencia; Feature Store
ofrece el almacén común que evita que la definición de una característica
difiera entre entrenamiento y servicio.

\begin{thebibliography}{99}\small
\bibitem{bc} G.~E.~P. Box y D.~R. Cox (1964). An analysis of transformations. \emph{JRSS B} 26, 211--252.
\bibitem{breiman} L.~Breiman, J.~Friedman, R.~Olshen y C.~Stone (1984). \emph{Classification and Regression Trees}. Wadsworth.
\bibitem{broder} A.~Broder (1997). On the resemblance and containment of documents. \emph{Compression and Complexity of Sequences}.
\bibitem{cerda} P.~Cerda y G.~Varoquaux (2022). Encoding high-cardinality string categorical variables. \emph{IEEE TKDE} 34, 1164--1176.
\bibitem{dh} D.~L. Donoho y P.~J. Huber (1983). The notion of breakdown point. En \emph{A Festschrift for Erich L.\ Lehmann}. Wadsworth.
\bibitem{fisher} W.~D. Fisher (1958). On grouping for maximum homogeneity. \emph{JASA} 53, 789--798.
\bibitem{good} I.~J. Good (1950). \emph{Probability and the Weighing of Evidence}. Griffin.
\bibitem{esl} T.~Hastie, R.~Tibshirani y J.~Friedman (2009). \emph{The Elements of Statistical Learning}, 2.ª ed. Springer.
\bibitem{ih} B.~Iglewicz y D.~Hoaglin (1993). \emph{How to Detect and Handle Outliers}. ASQC.
\bibitem{jolliffe} I.~T. Jolliffe (2002). \emph{Principal Component Analysis}, 2.ª ed. Springer.
\bibitem{lg} O.~Levy e Y.~Goldberg (2014). Neural word embedding as implicit matrix factorization. \emph{NeurIPS}.
\bibitem{lr} R.~J.~A. Little y D.~B. Rubin (2019). \emph{Statistical Analysis with Missing Data}, 3.ª ed. Wiley.
\bibitem{mb} D.~Micci-Barreca (2001). A preprocessing scheme for high-cardinality categorical attributes. \emph{SIGKDD Explorations} 3, 27--32.
\bibitem{mikolov} T.~Mikolov et al.\ (2013). Distributed representations of words and phrases and their compositionality. \emph{NeurIPS}.
\bibitem{glove} J.~Pennington, R.~Socher y C.~Manning (2014). GloVe: Global vectors for word representation. \emph{EMNLP}.
\bibitem{porter} M.~F. Porter (1980). An algorithm for suffix stripping. \emph{Program} 14, 130--137.
\bibitem{catboost} L.~Prokhorenkova et al.\ (2018). CatBoost: unbiased boosting with categorical features. \emph{NeurIPS}.
\bibitem{rubin} D.~B. Rubin (1976). Inference and missing data. \emph{Biometrika} 63, 581--592.
\bibitem{bpe} R.~Sennrich, B.~Haddow y A.~Birch (2016). Neural machine translation of rare words with subword units. \emph{ACL}.
\bibitem{shiffler} R.~E. Shiffler (1988). Maximum Z scores and outliers. \emph{The American Statistician} 42, 79--80.
\bibitem{siddiqi} N.~Siddiqi (2006). \emph{Credit Risk Scorecards}. Wiley.
\bibitem{sj} K.~Spärck Jones (1972). A statistical interpretation of term specificity. \emph{Journal of Documentation} 28, 11--21.
\bibitem{tukey} J.~W. Tukey (1977). \emph{Exploratory Data Analysis}. Addison-Wesley.
\bibitem{hash} K.~Weinberger et al.\ (2009). Feature hashing for large scale multitask learning. \emph{ICML}.
\bibitem{yj} I.-K. Yeo y R.~A. Johnson (2000). A new family of power transformations. \emph{Biometrika} 87, 954--959.
\bibitem{yos} J.~Yosinski et al.\ (2014). How transferable are features in deep neural networks? \emph{NeurIPS}.
\end{thebibliography}

\end{document}
